% concatenated from main.tex
\documentclass[10pt]{article}%

\usepackage[utf8]{inputenc}

\usepackage{amsmath}
\usepackage{amsfonts}
 \usepackage{mathrsfs}
\usepackage{amssymb, color}
\usepackage[linkcolor=black,anchorcolor=black,citecolor=black]{hyperref}
\usepackage[numbers,sort&compress]{natbib}
\usepackage{graphicx}
\numberwithin{equation}{section}
\usepackage[body={15.5cm,21cm}, top=3cm]{geometry}%
\setcounter{MaxMatrixCols}{30}
%TCIDATA{OutputFilter=latex2.dll}
%TCIDATA{Version=5.50.0.2960}
%TCIDATA{Codepage=936}
%TCIDATA{LastRevised=Saturday, August 06, 2011 06:33:25}
%TCIDATA{<META NAME="GraphicsSave" CONTENT="32">}
%TCIDATA{<META NAME="SaveForMode" CONTENT="1">}
%TCIDATA{BibliographyScheme=Manual}
%TCIDATA{Language=American English}
%BeginMSIPreambleData
\providecommand{\U}[1]{\protect\rule{.1in}{.1in}}
%EndMSIPreambleData
\providecommand{\U}[1]{\protect \rule{.1in}{.1in}}
\newtheorem{theorem}{Theorem}[section]
\newtheorem{acknowledgement}[theorem]{Acknowledgement}
\newtheorem{algorithm}[theorem]{Algorithm}
\newtheorem{axiom}[theorem]{Axiom}
\newtheorem{case}[theorem]{Case}
\newtheorem{claim}[theorem]{Claim}
\newtheorem{conclusion}[theorem]{Conclusion}
\newtheorem{condition}[theorem]{Condition}
\newtheorem{conjecture}[theorem]{Conjecture}
\newtheorem{corollary}[theorem]{Corollary}
\newtheorem{criterion}[theorem]{Criterion}
\newtheorem{definition}[theorem]{Definition}
\newtheorem{example}[theorem]{Example}
\newtheorem{exercise}[theorem]{Exercise}
\newtheorem{lemma}[theorem]{Lemma}
\newtheorem{notation}[theorem]{Notation}
\newtheorem{problem}[theorem]{Problem}
\newtheorem{proposition}[theorem]{Proposition}
\newtheorem{remark}[theorem]{Remark}
\newtheorem{solution}[theorem]{Solution}
\newtheorem{summary}[theorem]{Summary}
\newtheorem{assumption}[theorem]{Assumption}
\newenvironment{proof}[1][Proof]{\noindent \textbf{#1.} }{\  \rule{0.5em}{0.5em}}
\DeclareMathOperator*{\esssup}{ess\,sup}
\DeclareMathOperator*{\essinf}{ess\,inf}

\newcommand{\cG}{{\mathcal G}}
\newcommand{\cH}{{\mathcal H}}
\newcommand{\cP}{{\mathcal P}}
\newcommand{\frank}[1]{{\textcolor{blue}{#1}}}
\def \cA{{\mathcal A}}
%\def \cF{{\mathcal F}}
\def \cL{{\mathcal L}}
\def \cC{{\mathcal C}}
\def \cS{{\mathcal S}}
\def \E{\mathbf{E}}
\def \d{\mathsf{d}}
\def \EE{\widehat{\mathsf{E}}}
\def \P{\mathbf{P}}
\def \R{\mathbb{R}}
\def \F{\mathbb{F}}
\def \Q{\mathsf{Q}}
\def\LL{{\mathbb{L} }}
\def\X{\widehat{X}}
\def\S{\widehat{S}}
\def\WW{\widetilde{W}}
\def\FF{\widehat{F}}
\usepackage [numbers,sort&compress] {natbib} 

\def\d{\mathrm{d}}



\begin{document}
	\title{Mean-Field Backward Stochastic Differential Equations with Nonlinear Resistance and Double Mean Reflections}
	\author{ 	Hanwu Li \thanks{Research Center for Mathematics and Interdisciplinary Sciences, Shandong University, Qingdao 266237, Shandong, China. lihanwu@sdu.edu.cn.}
	\thanks{Frontiers Science Center for Nonlinear Expectations (Ministry of Education), Shandong University, Qingdao 266237, Shandong, China.}
    \thanks{Shandong Province Key Laboratory of Financial Risk, Shandong University, Qingdao 266237, Shandong, China.}
    \and Jin Shi \thanks{Research Center for Mathematics and Interdisciplinary Sciences, Shandong University, Qingdao 266237, Shandong, China. }}
	%\and Shige Peng\thanks{School of Mathematics and Qilu Institute of Finance, Shandong University,
	%peng@sdu.edu.cn. Li and Peng's research was
	%partially supported by the Tian Yuan Projection of the National Nature Sciences Foundation of China (No. 11526205 and No. 11626247)  and by the 111
	%Project (No. B12023).}}
	\date{}
	\maketitle
	
	\begin{abstract}
      In this paper, we investigate mean-field backward stochastic differential equation ($\text{MFBSDE}$) with double mean reflections and nonlinear resistance. Specifically, the constraints are formulated in terms of the expectation of the solution, and a compensating term is incorporated into the generator.  We establish the existence and uniqueness for both the case of Lipschitz generator and the case where the generator is quadratic and the terminal value is bounded. Finally, when the compensating term is absolutely continuous, we study the well-posedness of a variant type of doubly mean reflected MFBSDE with nonlinear resistance, whose generator depends on the density function of the compensating term.
      % {\color{red} it is better to use density functions instead of processes since the density is deterministic}. {\color{red}Rewrite the abstract}
	\end{abstract}

    \textbf{Key words}: Backward stochastic differential equations, Double mean reflections, Nonlinear resistance, Skorokhod problem, Density function

    \textbf{MSC-classification}: 60H10
	
\section{Introduction}
Pardoux and Peng \cite{PP} first pioneered the study of backward stochastic differential equations (BSDEs), which take the following form:
\begin{equation}
 Y_t=\xi+\int_t^T f(s,Y_s,Z_s)ds-\int_t^T Z_s dB_s,\quad \forall t\in[0,T].   
\end{equation}
They established the existence and uniqueness of solutions for BSDEs under the assumptions that the generator $f$ is uniformly Lipschitz continuous and the terminal value $\xi$ is square integrable.  Over the past decades, BSDE theory has evolved into several major branches to accommodate diverse modeling requirements, including BSDEs with jumps, forward-backward SDEs (FBSDEs), and constrained BSDEs. Among these, to tackle optimal stopping problems in mathematical finance and provide probabilistic representations for partial differential equations (PDEs) with obstacles constraints,  El Karoui, Kapoudjian, Pardoux, Peng and Quenez \cite{EKPPQ} further proposed the notion of reflected backward stochastic differential equations (RBSDEs). %,  which satisfies the following equation: 
%\begin{equation}\label{intro1}
 %Y_t=\xi+\int_t^T f(s,Y_s,Z_s)ds-\int_t^T Z_s dB_s+(K_T-K_t),\quad \forall t\in[0,T].
%\end{equation}
 By incorporating a non-decreasing reflection process $K$ that forces the solution $Y$ to stay above a given stochastic barrier (obstacle) $L$, i.e. %for all $t\in[0,T]$, $$Y_t\geq L_t,$$ 
 \begin{align}\label{constraint pointwise}
   Y_t\geq L_t, \ \forall t\in[0,T],  
 \end{align}
 the primary aim is to identify the minimal solution $Y$, which is uniquely determined by the Skorokhod condition $\int_{0}^{T} (Y_t-L_t) dK_t=0$. In contrast to pointwise constraints in classical RBSDEs, Briand, Elie, and Hu  \cite{BEH} introduced $\text{BSDEs}$ with mean reflection, whose constraints are articulated in terms of the distribution of the solution. Specifically, given a loss function $l$, the mean reflection constraint is formulated as 
 \begin{align}\label{constraint expectation}
 \mathbf{E}[l(t,Y_t)]\geq 0, \ \forall t\in[0,T].
 \end{align}
 Unlike the classical reflected case, the force $K$ is required to be a deterministic function that adheres to the Skorokhod condition$\int_{0}^{T}\mathbf{E}[l(t,Y_t)]dK_t=0$. The authors proved the existence and uniqueness of the so-called deterministic flat solution for mean reflected $\text{BSDEs}$ and further applied the results to  super-hedging problems within the context of ongoing risk management. More recently, Li \cite{L} investigated a more general reflection constraint. For given two nonlinear loss functions $L,R$  with $L\leq R$, the first component $Y$  satisfies the following condition $$\mathbf{E}[L(t,Y_t)]\leq 0\leq \mathbf{E}[R(t,Y_t)], \ \forall t\in[0,T],$$ 
and the minimality conditions for the bounded variation function $K$ are defined accordingly.  For more relevant advances in mean reflected BSDEs, {we refer the readers to \cite{BH,FS,HHLLW,LW,NQW,QF}.} %{\color{red}The references are not suitable here.}

Motivated by the demand for modeling complex systems in which individual dynamics depend on the collective behavior of the population, %such as in mean-field games and economics, 
the study of mean-field backward stochastic differential equations (MFBSDEs) was pioneered by Buckdahn, Djehiche, Li and Peng \cite{BDL} in their seminal work, which incorporates the distribution of solution to the generators $f$ for BSDEs. This framework was developed in the context of analyzing the limiting behavior of a high-dimensional system of forward-backward BSDEs. Subsequently, Buckdahn, Li and Peng \cite{BLP} further investigated $\text{MFBSDEs}$ within a more general framework. They established key results including the existence and uniqueness of solution, the comparison theorem and the stochastic representation for a nonlocal $\text{PDE}$. {Li \cite{Li'} firstly introduced reflected MFBSDEs, where the constraint is written as in \eqref{constraint pointwise}. Recently, Hu, Moreau and Wang \cite{HMW}  investigated the solvability of MFBSDEs with mean reflection, whose constraint is given in terms of the expectation of the solution as in \eqref{constraint expectation}. Correspondingly,  Li and Shi \cite{LS} extended the MFBSDEs to the case of double mean reflections.} %Correspondingly, Li \cite{Li'} firstly introduced reflected MFBSDE, by means of a pure probabilistic method to prove the well-posedness for reflected MFBSDE. The author proved that this type of reflected mean-field BSDEs can also be obtained as the limit equation of the mean-field BSDEs by the penalization method and gave the probabilistic interpretation of the nonlinear and nonlocal PDEs with the obstacles by the solutions of reflected MFBSDEs. Hu, Moreau and Wang \cite{HMW} are devoted to investigate the solvability of a generalized MFBSDE, whose driver also depends on the distribution of solution term $Y$. Using a fixed-point argument, BMO martingale theory and the $\theta$-method, they establish existence and uniqueness results for such BSDEs in several typical situations, including the case where the generator is Lipschitz continuous or quadratic with bounded and unbounded terminal condition. Similarly, Li and Shi \cite{LS} analyze $\text{MFBSDEs}$ with double mean reflections, whose generator and constraints both depend on the distribution of the solution. When the generator is Lipschitz continuous, based on the backward Skorokhod problem with nonlinear constraints, they investigate the solvability of the doubly mean reflected $\text{MFBSDEs}$ by constructing a contraction mapping. Furthermore, if the constraints are linear, the solution can also be constructed by a penalization method. For the case of quadratic growth, they obtain the existence and uniqueness results by using a fixed-point argument, the BMO martingale theory and the $\theta$-method. 
For further insights into the reflected MFBSDEs, {we refer to works \cite{CHM,DD,DDZ,DEH,HMW',HMW} and the references therein.} %{\color{red}The references are not suitable here.}
%Motivated by specific problems in mathematical finance {\color{red}too brief here, what kind of problems? Besides, ``motivated by" has been used in the previous paragraph},

{Inspired by Skorokhod equation-based methods for obstacle problems, and the practical demands for financial pricing and risk hedging under wealth constraints,}  Qian and Xu \cite{QX}  first extended the classical framework of RBSDEs by introducing a novel nonlinear resistance mechanism, where the reflecting force $K$ not only acts directly but also feeds back into the system's generator $f$  through a nonlinear mapping $H(K)$. This introduces a resistance effect, where the force $K$ influences the solution's dynamics in a more complex, nonlinear manner. % The authors are further driven by the goal of providing a pathwise construction of solutions using Skorokhod's equation, as opposed to traditional penalization methods, and by the application of this theory to model super-hedging problems with wealth constraints and optimal stopping with recursive utilities in finance. 
{The authors  established a pathwise construction of solutions via Skorokhod’s equation, in contrast to conventional penalization methods. They also applied the developed theory to financial problems, including super-hedging under wealth constraints and optimal stopping with recursive utilities. The RBSDEs with resistance share some similarities with a variant of Skorokhod's obstacle problem introduced in \cite{BEK}, which was further studied by Ma and Wang \cite{MW}.} {Subsequently, Luo \cite{LP}  investigated the well-posedness of mean-reflected MFBSDEs with nonlinear resistance. The author addressed both Lipschitz and quadratic generator conditions and presented an application to superhedging problems subject to risk constraints.} %By the contraction mapping argument, the author first proves that the equation admits a unique deterministic flat local solution on a small time interval. Then by a two-step approach and further provide an application to the super-hedging problem with risk constraint.

The present paper is dedicated to the study of the following $\text{MFBSDE}$ with double mean reflections and nonlinear resistance over the time interval $[0,T]$ 
\begin{equation}\label{MFBSDEDMR}
\begin{cases}
Y_t=\xi+\int_t^T f(s,Y_s,\mathbf{P}_{Y_s},Z_s,\mathbf{P}_{Z_s},G_s(K))ds-\int_t^T Z_s dB_s+K_T-K_t, \\
\mathbf{E}[L(t,Y_t)]\leq 0\leq \mathbf{E}[R(t,Y_t)], \\
K_t=K^R_t-K^L_t,\ K^R,K^L\in I[0,T], \\
\int_0^T \mathbf{E}[R(t,Y_t)]dK_t^R=\int_0^T \mathbf{E}[L(t,Y_t)]dK^L_t=0,
\end{cases}
\end{equation}
where $G_s(\cdot)$ is a functional of $K$, which depends on the path of $K$ during time period $[0,s]$. To some extent, this feature makes the equation regarded as a forward-backward equation. This leads to the difficulty of constructing the solution by stitching the local ones on each small time obtained by the contraction mapping method. Motivated by \cite{LP}, we investigate the well-posedness of Eq. \eqref{MFBSDEDMR} using a so-called two-step method. More precisely, for both the Lipschitz case and the quadratic growth case, the key difference between the two cases lies in the time interval: in the Lipschitz case, the contraction mapping holds on a general fixed-length time interval, whereas in the quadratic growth case, the mapping is contractive only on a small time interval whose length depends on the bound of the component $Y$. Therefore, for the quadratic growth case, we first establish a uniform bound for the component $Y$ under an additional assumption on the generator $f$. Using this bound, we obtain the local existence and uniqueness of the solution for Eq. \eqref{MFBSDEDMR}. Furthermore, we have the global well-posedness of Eq. \eqref{MFBSDEDMR} under the given resistance term.
In contrast, for the Lipschitz case, the well-posedness of MFBSDE with double mean reflections (for the given resistance term) is already known, and we refer the reader to \cite{LS} for details. Finally, when the generator includes the nonlinear resistance term, we apply the contraction mapping principle to establish the well-posedness of the solution for Eq. \eqref{MFBSDEDMR} under both cases.
%{\color{blue}Motivated by \cite{LP}, according to the properties of solution to Skorokhod problem, we investigate the well-posedness of Eq. \eqref{MFBSDEDMR} using a so-called two-step method, which combines the stitching method and the fixed point method. More precisely, under the Lipschitz continuity condition on the generator, for a given resistance term, Eq. \eqref{MFBSDEDMR} is equivalent to MFBSDE with double mean reflections under Lipschitz case. By combining with the principle of standard contraction mapping with respect to the nonlinear resistance term, we obtain the existence and uniqueness for Eq. \eqref{MFBSDEDMR}. When the generator is quadratic and the terminal value is bounded, we first derive a uniform estimate for the component $Y$ on an arbitrary small time interval. Since the generator $f$ is quadratic with respect to $z$, we then establish the local well-posedness of Eq.\eqref{MFBSDEDMR} on a sufficiently small time interval, which  depends crucially on the uniform estimate for the component $Y$. For a given resistance term, we further establish the well-posedness of the global solution via the existence and uniqueness of local solutions as well as uniform estimates of solutions. Finally, for Eq.\eqref{MFBSDEDMR} defined on a general time horizon, we apply the contraction mapping principle to the nonlinear resistance term to derive its well-posedness.
%   {\color{red}Too brief. Why we need to use a two-step method? Usually, we construct the solution on small time interval by fixed point. Then, dividing the whole time interval into finite small interval. Finally, pasting the solution on each small interval, we obtain the solution on the whole interval. Explain here why this does not hold.} 

Lastly, we consider doubly mean reflected MFBSDEs with a density function, {where $G_{\cdot}(K)$ in Eq. \eqref{MFBSDEDMR} denotes the Radon-Nikodym derivative of 
$K$ with respect to the Lebesgue measure.  In fact, the targeted solution can be constructed based on the solution to doubly mean-reflected MFBSDEs without a resistance term. To this end, the first problem we aim to address is to identify the conditions ensuring the absolute continuity of $K$, where $(Y,Z,K)$ stands for the solution to the underlying doubly mean reflected MFBSDE. Roughly speaking, the desired result holds provided that the loss functions $L$ and $R$  are sufficiently smooth and $Z$ satisfies certain enhanced integrability conditions. Moreover, by virtue of Malliavin  calculus, we derive a sufficient condition that guarantees the integrability of $Z$. Further, under a specific structural assumption on the generator (see Eq. \eqref{k+k-}), we construct the solution to the density-dependent doubly mean reflected MFBSDE. In particular, if the loss functions are linear and the generator is separable with respect to the variable $k$ (see Eq. \eqref{repre of f}), the corresponding solution is  unique.
} %For reflected MFBSDE
%\begin{align}\label{introMFBSDE}
%    Y_t=\xi+\int_t^T f(s,Y_s,\mathbf{E}[Y_s],Z_s,\mathbf{E}[Z_s])ds-\int_t^T Z_s dB_s+K_T-K_t,
%\end{align}
%we obtain the solution $K$ for Eq.\eqref{introMFBSDE} is Lipschitz continuous and it has density process by applying a priori estimates for MFBSDEs and $p$-order moment estimation for the second component $Z$ of solution to MFBSDE. Combining with an additional known MFBSDE with double mean reflections, we obtain the existence for MFBSDE with terminal value $\xi$, the generator $\left\{f(s,Y_s,\mathbf{E}[Y_s],Z_s,\mathbf{E}[Z_s],k_s)\right\}_{s\in[0,T]}$, loss functions $L,R$ and density process $k$.  And when the generator $f$ has the following representation:
% \begin{align}
% f(t,\omega,y,y',z,z',k)=a(t,\omega,y,y',z,z')+b(t,y',z',k).
%  \end{align}
 % and the constraints are linear, the uniqueness for MFBSDE with double mean reflections and density process is derived.}  {\color{red}this paragraph is not good.}
 
This paper is organized as follows. In Section 2, we recall some basic results about the Skorokhod problem and a priori estimates for MFBSDE. In Section 3,  we study the existence and uniqueness of doubly mean reflected MFBSDEs with nonlinear resistance when the generators are Lipschitz continuous. Then, we consider the case of quadratic generators and bounded terminal values in Section 4. Finally, the well-posedness of the doubly mean reflected MFBSDEs with density functions is investigated in Section 5. 
 
\section{Preliminaries}

 Fix a finite time horizon $T>0$. Consider a filtered probability space $(\Omega,\mathcal{F},\mathbb{F},\P)$ satisfying the usual conditions of right continuity and completeness. Let $B$ be a $1$-dimensional standard Brownian motion. Actually, all the results in this paper still hold when we consider the multi-dimensional Brownian motion. For each $p>1$ and for any finite time interval $[u,v]\subset [0,T]$, we first introduce the following notations.
 
\begin{itemize}	
\item $\mathcal{L}^p_v$: the collection of real-valued $\mathcal{F}_{v}$-measurable random variable $\xi$ satisfying 
	$$\left \| \xi \right \|_{\mathcal{L}^p}=\mathbf{E} \left [ \left | \xi  \right | ^p \right ]^{\frac{1}{p} }< \infty. $$
	 
	\item $\mathcal{L}^{\infty}_v$: the collection of real-valued $\mathcal{F}_{v}$-measurable random variable $\xi$ satisfying
	 $$\left \| \xi  \right \|_{\mathcal{L}^\infty}=\operatorname{ess\,sup}\left |\mathbf{\xi} \left ( \omega  \right )   \right |<\infty. $$
	
	\item $\mathcal{H}^{p}\left[u,v\right]$: the collection of real-valued $\mathcal{F}$-progressively measurable processes $\left(z_{t}\right)_{u\leq t\leq v}$ satisfying 
	$$\left \| z \right \|_{\mathcal{H}^p }=\mathbf{E}\left [\left ( \int_{u}^{v}\left | z_{t} \right |^2 d t \right )^{\frac{p}{2}}   \right ]^{\frac{1}{p} }<\infty. $$
	
	\item $S^p\left[u,v\right]$: the collection of real-valued $\mathcal{F}$-adapted continuous processes $\left(y_{t}\right)_{u\leq t \leq v}$ satisfying 
	$$ \left \| y \right \|_{\mathcal{S}^p }=\mathbf{E}\left [ \sup _{t\in[u,v]} \left | y_{t} \right |^p \right ]^{\frac{1}{p}}  <\infty.$$
	
	\item $S^{\infty}\left[u,v\right]$: the collection of real-valued $\mathcal{F}$-adapted continuous processes $\left(y_{t}\right)_{u\leq t \leq v}$ satisfying
	$$\left \| y \right \|_{\mathcal{S}^{\infty } } =\operatorname{ess\,sup}_{(t,\omega)\in\left[u,v\right]\times\Omega}\left | y_t\left (\omega\right )  \right |  <\infty.$$
	
	\item $\mathcal{P}_{p}\left(\mathbb{R}\right)$: the collection of all probability measures over $\left(\mathbb{R},\mathcal{B}\left(\mathbb{R}\right)\right)$ with finite $p$-th moment, endowed with the $p$-Wasserstein distance $W_{p}$, where $p > 1$ is a positive integer and $$W_{p}(\mu,\nu)=\left(\inf_{X \sim \mu,Y \sim\nu}\mathbf{E}[\left|X-Y\right|]^p\right)^{\frac{1}{p}}.$$
	
	%\item $\mathbb{L}^p_v$: the collection of real-valued $\mathcal{F}_{v}$-measurable random variable $\xi$ satisfying $\mathbf{E}\left[e^{p\left|\xi\right|}\right]<\infty.$
	
	%\item $\mathbb{S}^p[u,v]$: the collection of all stochastic processes $Y$ such that $e^Y\in\mathcal{S}^p\left[u,v\right].$
	
	%\item $\mathbb{L}_v$: the collection of all random variable $\xi\in\mathbb{L}_v^p$ for any $p> 1$, and $\mathcal{H}\left[u,v\right]$ and $\mathbb{S}\left[u,v\right]$ are defined similarly.
	
	\item $\mathcal{T}_{t}[u,v]$: the collection of $\left[u,v\right]$-valued $\mathcal{F}$-stopping times $\tau$ such that $\tau \geq t$, $\mathbf{P}$-a.s. 
	
	\item $BMO\left[u,v\right]$: the collection of real-valued progressively measurable processes $\left(z_{t}\right)_{u\leq t \leq v}$ such that 
	$$\left \| z \right \|_{BMO}:=\sup_{\tau\in \mathcal{T}_{u}\left[u,v\right]}\operatorname{ess\,sup}_{\omega\in\Omega}\mathbf{E}_{\tau}\left[\int_{\tau}^{v}\left | z_{s} \right |^2 d s \right]^{\frac{1}{2}}<\infty. $$
   
	
	\item $C\left[u,v\right]$: the set of continuous functions from $\left[u,v\right]$ to $\mathbb{R}$.

	\item  $BV\left[u,v\right]$: the set of functions $K\in C\left[u,v\right]$ with $K_u=0$ and $K$ is of bounded variation on $\left[u,v\right]$.
	
	\item  $I\left[u,v\right]$: the set of functions in $C\left[u,v\right]$ starting from the origin which is nondecreasing.
 \end{itemize}
	
When the interval $[u,v]=[0,T]$, we always omit the time index. %We denote by $\ell$ the corresponding collections for the processes and functions with indices on $[0,T]$. 
 For example, we write $\mathcal{H}^{p},\mathcal{S}^p,\mathcal{S}^{\infty}$ and so on. 

 \iffalse
 We first introduce some notations which will be frequently used in this paper.

\begin{itemize}
\item $L^2(\mathcal{F}_t)$: the set of real-valued $\mathcal{F}_t$-measurable random variable $\xi$ such that $\E[|\xi|^2]<\infty$, $t\in[0,T]$;
\item $\mathcal{S}^2$: the set of real-valued adapted continuous processes $Y$ on $[0,T]$ with $\E[\sup_{t\in[0,T]}|Y_t|^2]<\infty$;
\item $\mathcal{H}^2$: the set of real-valued progressively measurable processes $Z$ such that $\E[\int_0^T|Z_t|^2dt]<\infty$;
\item $C[0,T]$: the set of continuous functions from $[0,T]$ to $\mathbb{R}$;
\item $BV[0,T]$: the set of functions in $C[0,T]$ starting from the origin with bounded variation on $[0,T]$;
\item $I[0,T]$: the set of functions in $C[0,T]$ starting from the origin which is nondecreasing;
\item $\mathcal{P}_p(\mathbb{R}^k)$: the collection of all probability measures on $(\mathbb{R}^k,\mathcal{B}(\mathbb{R}
^k))$ with finite $p$-th moment, endowed with the $p-\text{Wasserstein}$ distance $W_p$, where $p\geq 1$, $k$ is a positive integer.
\end{itemize}
\fi

\subsection{Skorokhod problem}
In this section, we present the Skorokhod problem with nonlinear constraints, systematically addressing its definition, well-posedness, and essential properties of the resulting solutions. 
\begin{definition}\label{def1}
Let $s\in C[0,T]$,  and $l,r:[0,T] \times \mathbb{R}\rightarrow \mathbb{R}$ be two functions with $l\leq r$.  A pair of functions $(x,K)\in C[0,T]\times BV[0,T]$ is called a solution to the  Skorokhod problem for $s$ with nonlinear constraints $l,r$ ($(x,K)=\mathbb{SP}_l^r(s)$ for short) if % and $L(0,S_0)\leq 0\leq R(0,S_0)$
\begin{itemize}
\item[(i)] $x_t=s_t+K_t$;
\item[(ii)] $l(t,x_t)\leq 0\leq r(t,x_t)$;
\item[(iii)]  $K_{0-}=0$ and $K$ has the decomposition $K=K^r-K^l$, where $K^r,K^l$ are nondecreasing functions satisfying  
\begin{align}
\int_0^{T} I_{\{l(s,x_s)<0\}}dK^l_s=0, \  \int_0^{T} I_{\{r(s,x_s)>0\}}dK^r_s=0.
\end{align}
\end{itemize}
\end{definition}

We propose the following assumption on the functions $l,r$.
\begin{assumption}\label{asslr}
The functions $l,r:[0,T] \times \mathbb{R}\rightarrow \mathbb{R}$ satisfy the following conditions
\begin{itemize}
\item[(i)] For each fixed $x\in\mathbb{R}$, $l(\cdot,x),r(\cdot,x)\in C[0,T]$;
\item[(ii)] For any fixed $t\geq 0$, $l(t,\cdot)$, $r(t,\cdot)$ are strictly increasing;
\item[(iii)] There exist two positive constants $0<c<C<\infty$, such that for any $t\in [0,T]$ and $x,y\in \mathbb{R}$,
\begin{align*}
&c|x-y|\leq |l(t,x)-l(t,y)|\leq C|x-y|,\\
&c|x-y|\leq |r(t,x)-r(t,y)|\leq C|x-y|.
\end{align*} 
\item[(iv)] $\inf_{(t,x)\in[0,T]\times\mathbb{R}}(r(t,x)-l(t,x))>0$.
\end{itemize}
\end{assumption}

\begin{theorem}[\cite{Li}]\label{SP}
Suppose that $l,r$ satisfy Assumption \ref{asslr}. For any given $s\in C[0,T]$, there exists a unique pair of solution to the Skorokhod problem $(x,K)=\mathbb{SP}_l^r(s)$. Moreover, for any $t\geq 0$, let $\phi_t,\psi_t$ be the unique solutions to the following equations, respectively, 
    \begin{align*}
        l(t,s_t+x)=0,\ r(t,s_t+x)=0.
    \end{align*}
    Then, we have 
    \begin{align*}
        &K_t=-\max\left((-\phi_0)^{+}\wedge  \inf_{r\in[0,T]}(-\psi_r), \sup_{s\in[0,T]}\left[(-\phi_s) \wedge  \inf_{r\in[s,t]}(-\psi_r)\right]\right),\\
    \end{align*}
    and 
   \begin{align*}
        &\sup_{t\in[0,T]}\left|K_t\right|\leq \sup_{t\in[0,T]}\left|\phi_t\right|+\sup_{t\in[0,T]}\left|\psi_t\right|.\\
    \end{align*} 
\end{theorem}

The following proposition provides the continuity property of the solution  to the Skorokhod problem with respect to input function $s$ and the reflecting boundary functions $l,r$.

\begin{proposition}[\cite{Li}]\label{continuity}
Suppose that $(l^i,r^i)$ satisfy Assumption $\ref{asslr}$, $i=1,2$.
Given  $s^i\in C[0,T]$, let $(x^i,K^i)$ be the solution to the Skorokhod problem $\mathbb{SP}_{l^i}^{r^i}(s^i)$. Then, we have
\begin{equation}\label{diffk}
\sup_{t\in[0,T]}\left|K^1_t-K^2_t\right|
\leq \frac{C}{c}\sup_{t\in[0,T]}\left|s^1_t-s^2_t\right|+\frac{1}{c}(\bar{L}_T\vee\bar{R}_T),
\end{equation}
where
\begin{align*}
\bar{L}_T=\sup_{(t,x)\in[0,T]\times \mathbb{R}}\left|{l}^1(t,x)-{l}^2(t,x)\right|,\
\bar{R}_T=\sup_{(t,x)\in[0,T]\times \mathbb{R}}\left|{r}^1(t,x)-{r}^2(t,x)\right|.
\end{align*}
\end{proposition}
%{\color{blue}\begin{proposition}[\cite{Li}]\label{respreK}
 %   Given $s\in C[0,T]$. For any $t\geq 0$, let $\phi_t,\psi_t$ be the unique solutions to the following equations, respectively, 
  %  \begin{align*}
   %     l(t,s_t+x)=0,\ r(t,s_t+x)=0.
    %\end{align*}
    %Therefore, we have 
    %\begin{align*}
     %   K_t=-\max\left((-\phi_0)^{+}\wedge  \inf_{r\in[0,T]}(-\psi_r), \sup_{s\in[0,T]}\left[(-\phi_s) \wedge  \inf_{r\in[s,t]}(-\psi)\right]\right)
    %\end{align*}
%\end{proposition}}
%{\color{red}I have put the above results into theorem \ref{SP}.}

%\begin{remark}
%    Actually, Theorem \ref{SP} and Proposition \ref{continuity} still hold if the input function and the reflecting boundary functions are c\`{a}dl\`{a}g. The main difference is that each component of the solution $(x,K)$ to the Skorokhod problem is c\`{a}dl\`{a}g.
%\end{remark}

\subsection{Backward Skorokhod problem}
In this section, we introduce the backward Skorokhod problem with nonlinear constraints, including its definition, well-posedness, and key properties of its solution, while also clarifying its relationship with the classical Skorokhod problem.

\begin{definition}\label{def2}
Let $s\in C[0,T]$, $a\in \mathbb{R}$ and $l,r:[0,T]\times \mathbb{R}\rightarrow \mathbb{R}$ be two functions such that $l\leq r$ and $l(T,a)\leq 0\leq r(T,a)$. A pair of functions $(x,K)\in C[0,T]\times BV[0,T]$ is called a solution of the backward Skorokhod problem for $s$ with nonlinear constraints $l,r$ ($(x,K)=\mathbb{BSP}_l^r(s,a)$ for short) if 
\begin{itemize}
\item[(i)] $x_t=a+s_T-s_t+K_T-K_t$;
\item[(ii)] $l(t,x_t)\leq 0\leq r(t,x_t)$, $t\in[0,T]$;
\item[(iii)]  $K$ has the decomposition $K=K^r-K^l$, where $K^r,K^l\in I[0,T]$ satisfy %are nondecreasing functions satisfying %$k_{0}=0$
\begin{align}\label{iii}
\int_0^T I_{\{l(s,x_s)<0\}}dK^l_s=0, \  \int_0^T I_{\{r(s,x_s)>0\}}dK^r_s=0.
\end{align}
\end{itemize}
\end{definition}

%Let us first introduce the assumption made for the reflecting boundary functions. 
%\begin{assumption}\label{asslr}
%The functions $l,r:[0,T]\times \mathbb{R}\rightarrow \mathbb{R}$ satisfy the following conditions
%\begin{itemize}
%\item[(i)] For each fixed $x\in\mathbb{R}$, $l(\cdot,x),r(\cdot,x)\in C$;
%\end{itemize}
%\quad(ii), (iii) and (iv) are same as Assumption $\ref{assgh}$.
%\end{assumption}

\begin{theorem}[\cite{L}]\label{BSP}
Let Assumption \ref{asslr} holds. For any given $s\in C[0,T]$ and $a\in \mathbb{R}$ with $l(T,a)\leq 0\leq r(T,a)$, there exists a unique solution to the backward Skorokhod problem $(x,k)=\mathbb{BSP}_l^r(s,a)$. %{\color{red}Add the relation with forward Skorokhod problem since it is used in Step 1 of the proof of Theorem \ref{thmbddterminal}.}
\end{theorem}
\begin{remark}\label{rem2.7}
Set
\begin{equation}
    \begin{split}
     &\bar{s}_t=a+s_T-s_{T-t},\ x_t=\bar{x}_{T-t},\ K_t=\bar{K}_{T}-\bar{K}_{T-t}, \ t\in[0,T],\\ 
     &\bar{l}\left(t,x\right)=l\left(T-t,x\right),\ \bar{r}\left(t,x\right)=r\left(T-t,x\right), \ (t,x)\in[0,T]\times\mathbb{R}.
    \end{split}
\end{equation}
According to Theorem 3.9 in \cite{L},  $\left(\bar{x},\bar{K}\right)$ is the unique solution to the Skorokhod problem $\mathbb{SP}^{\bar{r}}_{\bar{l}}\left(\bar{s}\right)$.
\end{remark}   
 
The following proposition provides the continuous dependence of the solution with respect to the input function $s$ and the reflecting boundary functions $l,r$.

\begin{proposition}[\cite{L}]\label{continuous}
Given $a^i\in\mathbb{R}$, $s^i\in C[0,T]$, $l^i,r^i$ satisfy Assumption \ref{asslr} and $l^i(T,a^i)\leq 0\leq r^i(T,a^i)$, $i=1,2$, 
 let $(x^i,k^i)$ be the solution to the backward Skorokhod problem $\mathbb{BSP}_{l^i}^{r^i}(s^i,a^i)$, $i=1,2$. Then, we have
\begin{equation}\label{diffK}
\sup_{t\in[0,T]}\left|K^1_t-K^2_t\right|
\leq  2\frac{C}{c}\left|a^1-a^2\right|+4\frac{C}{c}\sup_{t\in[0,T]}\left|s^1_t-s^2_t\right|+\frac{2}{c}(\bar{L}_T\vee\bar{R}_T),
\end{equation}
where $\bar{L}_T$ and $\bar{R}_T$ are the same as in Proposition \ref{continuity}.
%\begin{align*}
%&\bar{L}_T=\sup_{(t,x)\in[0,T]\times \mathbb{R}}\left|\bar{l}^1(t,x)-\bar{l}^2(t,x)\right|=\sup_{(t,x)\in[0,T]\times \mathbb{R}}\left|{l}^1(t,x)-{l}^2(t,x)\right|,\\
%&\bar{R}_T=\sup_{(t,x)\in[0,T]\times \mathbb{R}}\left|\bar{r}^1(t,x)-\bar{r}^2(t,x)\right|=\sup_{(t,x)\in[0,T]\times \mathbb{R}}\left|{r}^1(t,x)-{r}^2(t,x)\right|.
%\end{align*}
\end{proposition}
\subsection{A priori estimates for MFBSDEs}
%Similar to Proposition 3.3 in \cite{LX}, we have a priori estimates for MFBSDE, which satisfy the following equation:

In this subsection, we consider the MFBSDEs taking the following form:
\begin{equation}\label{MFBSDE}
    Y_t=\xi+\int_{t}^{T}f(s,Y_s, \mathbf{P}_{Y_s},Z_s,\mathbf{P}_{Z_s})ds-\int_{t}^{T}Z_sdB_s.
\end{equation}
We propose the following assumptions for the generator $f$ and the terminal value $\xi$.

\begin{assumption}\label{asslip2}
  \begin{itemize}
   \item [(i)] The terminal value $\xi\in \mathcal{L}^p_T$,$ \ p>1$;
      \item [(ii)] The process $\left\{f_t(0):=f(t,0,\delta_0,0,\delta_0)\right\}_{t\in[0,T]}\in \mathcal{H}^p$. There exists a constant $\lambda>0$ such that for all $t\in[0,T]$, $y, y'\in \mathbb{R}$, $\mu, \mu' \in \mathcal{P}_1(\mathbb{R})$, $z, z'\in \mathbb{R}$, $\nu, \nu'\in\mathcal{P}_1(\mathbb{R})$,
      \begin{align*}
\left|f(t,y,\mu,z,\nu)-f(t,y',\mu',z',\nu')\right|\leq \lambda\left(\left|y-y'\right|+\left|z-z'\right|+W_1(\mu,\mu')+W_1(\nu,\nu')\right).
\end{align*}
%The process $\left\{f_t(0)\right\}_{t\in[0,T]}$. is uniformly bounded by some constant {\color{red}$f_t(0)\in \mathcal{H}^p$ is not enough?}, where $f_t(0)=f\left(t, 0, 0, 0,0 \right)$  {\color{red}if $f$ depends on the distribution of the solution, $f_t(0)$ should be $f(t,0,\delta_0,0,\delta_0)$ and the Lipschitz continuity should be corrected accordingly}
    %\begin{align*}
%\left|f(t,y_1,y'_1,z_1,z'_1)-f(t,y_2,y'_2,z'_1,z'_2)\right|\leq \lambda\left(\left|y_1-y_2\right|+\left|y'_1-y'_2\right|+\left|z_1-z_2\right|+\left|z'_1-z'_2\right|\right).
%\end{align*}
  \end{itemize}  
\end{assumption}

Similar to Proposition 3.3 in \cite{LX}, we have the following a priori estimates for MFBSDEs. 
\begin{lemma}\label{proMFBSDE}
Let Assumption \ref{asslip2} hold. Assume $(Y, Z)$ is the solution of the mean-field BSDE \eqref{MFBSDE}. Then, there exists a constant $M$ depending on the constant $\lambda,\ p,\ T$ such that
\begin{align}
&\E\left[ \sup_{s \in [0,T]} |Y_s|^p \right] + \E\left[\left(\int_0^T |Z_s|^2ds\right)^{\frac{p}{2}} \right]\leq M(\lambda, p, T) \E\left[ |\xi|^p + \left( \int_0^T |f_s(0)| ds \right)^p\right].
\end{align}
\end{lemma}
 
\iffalse
\subsection{Doubly mean reflected BSDEs}
In this subsection, we introduce the doubly mean reflected BSDEs studied in \cite{Li'}. Roughly speaking, the constraints for the doubly mean reflected BSDEs are made on the law of the solution but not on the paths of the solution as for the reflected BSDEs. More specifically, the doubly mean reflected BSDE is of the following type:% The main purpose of this paper is to study the BSDE with double mean reflections of the following type
\begin{equation}\label{MRBSDE}
\begin{cases}
Y_t=\xi+\int_t^T f(s,Y_s,Z_s)ds-\int_t^T Z_s dB_s+K_T-K_t,\ t\in[0,T], \\
\E[L(t,Y_t)]\leq 0\leq \E[R(t,Y_t)], \ K_t=K^R_t-K^L_t,\ K^R,K^L\in I[0,T], \\
 \int_0^T \E[R(t,Y_t)]dK_t^R=\int_0^T \E[L(t,Y_t)]dK^L_t=0,
\end{cases}
\end{equation}
Here, the driver $f:[0,T]\times\Omega\times\mathbb{R}\times\mathbb{R}\rightarrow \mathbb{R}$ is independent from the law of the solution satisfies (H1) and (H2). The loss functions $L,R:\Omega\times [0,T]\times\mathbb{R}\rightarrow \mathbb{R}$ are measurable maps with respect to $\mathcal{F}_T\times \mathcal{B}([0,T])\times \mathcal{\mathbb{R}}$ satisfying the following conditions.

\begin{assumption}\label{ass2}
\begin{itemize}
\item[(1)] For any fixed $(\omega,x)\in \Omega\times\mathbb{R}$, $L(\omega,\cdot,x), R(\omega,\cdot,x)$ are continuous;%for any fixed $x\in\mathbb{R}$, $L(\cdot,\cdot,x), R(\cdot,\cdot,x)$ is progressively measurable and  
\item[(2)]  $\E[\sup_{t\in[0,T]}|L(t,0)|]<\infty$, $\E[\sup_{t\in[0,T]}|R(t,0)|]<\infty$;%for any fixed $(t,x)\in [0,T]\times \mathbb{R}$,
\item[(3)] For any fixed $(\omega,t)\in \Omega\times [0,T]$, $L(\omega,t,\cdot),R(\omega,t,\cdot)$ are strictly increasing and there exists two constants $0<c<C$ such that for any $x,y\in \mathbb{R}$,
\begin{align*}
&c|x-y|\leq |L(\omega,t,x)-L(\omega,t,y)|\leq C|x-y|,\\
&c|x-y|\leq |R(\omega,t,x)-R(\omega,t,y)|\leq C|x-y|;
\end{align*}
\item[(4)] $\inf_{\omega,t,x} (R(\omega,t,x)-L(\omega,t,x))>0$.
\end{itemize}
\end{assumption}
 
 \begin{theorem}[\cite{Li'}]\label{existence and uniqueness for MRBSDE}
Suppose that the driver $f[0,T]\times\Omega\times\mathbb{R}\times\mathbb{R}\rightarrow \mathbb{R}$ satisfies (H1) and (H2) and the loss functions $L,R$ satisfy Assumption \ref{ass2}. Let $\xi\in L^2(\mathcal{F}_T)$ be such that $\E[L(T,\xi)]\leq 0\leq \E[R(T,\xi)]$. Then, the doubly mean reflected BSDE \eqref{MRBSDE} admits a unique solution $(Y,Z,K)\in \mathcal{S}^2\times\mathcal{M}^2\times BV[0,T]$. 
 \end{theorem}
 
The construction for solutions to doubly mean reflected BSDE plays a key role in the investigation of mean field BSDEs with double mean reflections. In the following, we introduce the main idea of the construction for the constant driver case, i.e., $f\equiv C$ with $C\in \mathcal{H}^2$. In this case, set
\begin{align*}
&\widetilde{Y}_t=\xi+\int_t^T C_s ds-(M_T-M_t),\\
&s_t=\E[\int_0^t C_sds], \ a=\E[\xi],
\end{align*}
where $M_t=\E_t[\xi+\int_0^T C_s ds]-\E[\xi+\int_0^T C_s ds]$. For any $(t,x)\in[0,T]\times \mathbb{R}$, we define
\begin{align}\label{lr}
l(t,x):=\E[L(t,\widetilde{Y}_t-\E[\widetilde{Y}_t]+x)], \ r(t,x):=\E[R(t,\widetilde{Y}_t-\E[\widetilde{Y}_t]+x)].
\end{align}
We may check that $l,r$ satisfy Assumption \ref{ass1} and 
\begin{align*}
l(T,a)=\E[L(T,\xi)]\leq 0\leq \E[R(T,\xi)]=r(T,a).
\end{align*}
Let $(x,K)$ be the solution to the backward Skorokhod problem $\mathbb{BSP}_l^r(s,a)$. Now, we set 
\begin{align*}
Y_t=\xi+\int_t^T C_s ds-(M_T-M_t)+(K_T-K_t).%=\xi+\int_t^T C_sds-\int_t^T Z_s dB_s+K_T-K_t.
\end{align*}
Then, $Y$ is the first component of the solution to doubly mean reflected BSDE with terminal value $\xi$, constant driver $C$ and loss functions $L,R$. 

Besides, for any $t\in[0,T]$, let $\bar{r}_t$ and $\bar{l}_t$ be the unique solution to the following two equations, respectively
\begin{align*}
r(t,x)=0, \ l(t,x)=0,
\end{align*}
where $l,r$ are defined in \eqref{lr}. Then, for any $t\in[0,T]$, we have
%\begin{theorem}[\cite{FS,Li'}]\label{theorem3.6}
 %Under Assumptions (H1'), (H2'), (H3) and (H4), let $(Y,Z,K)$ be the solution to the doubly mean reflected BSDE  \eqref{nonlinearyz}. Then, for  any $t\in[0,T]$, we have
\begin{align}\label{representationEYt}
\E[Y_t]=\sup_{q\in[t,T]}\inf_{s\in[t,T]}R_t(s,q)=\inf_{s\in[t,T]}\sup_{q\in[t,T]}R_t(s,q),
\end{align}
where 
\begin{align*}
R_t(s,q)=\E\left[\int_t^{s\wedge q} C_u du+\xi I_{\{s\wedge q=T\}}\right]+\bar{l}_q I_{\{q< T,q\leq s\}}+\bar{r}_s I_{\{s<q\}}.
\end{align*}
%\end{theorem}
\fi



\section{Doubly mean reflected MFBSDEs with nonlinear resistance: the Lipschitz case}

%{\color{red}Write mean field BSDE as MFBSDE throughout this paper.} 
The main purpose of this paper is to study the MFBSDE with nonlinear resistance and double mean reflections of the following type
\begin{equation}\label{nonlinearyz}
\begin{cases}
Y_t=\xi+\int_t^T f(s,Y_s,\P_{Y_s},Z_s,\P_{Z_s},G_s(K))ds-\int_t^T Z_s dB_s+K_T-K_t, \\
\E[L(t,Y_t)]\leq 0\leq \E[R(t,Y_t)], \\
K_t=K^R_t-K^L_t,\ K^R,K^L\in I[0,T],\\
\int_0^T \E[R(t,Y_t)]dK_t^R=\int_0^T \E[L(t,Y_t)]dK^L_t=0.
\end{cases}
\end{equation}
The loss functions $L,R:\Omega\times [0,T]\times\mathbb{R}\rightarrow \mathbb{R}$ are measurable maps with respect to $\mathcal{F}_T\times \mathcal{B}([0,T])\times \mathcal{B}(\mathbb{R})$ satisfying the following conditions.

\begin{assumption}\label{assLR}
\begin{itemize}
\item[(1)] For any fixed $(\omega,x)\in \Omega\times\mathbb{R}$, $L(\omega, \cdot, x), R(\omega, \cdot, x)$ are continuous;%for any fixed $x\in\mathbb{R}$, $L(\cdot,\cdot,x), R(\cdot,\cdot,x)$ is progressively measurable and  
\item[(2)]  There exists a constant $M>0$ such that $$\mathbf{E}\left[\sup_{t\in[0,T]}\left|L(t,0)\right|+\sup_{t\in[0,T]}\left|R(t,0)\right|\right]\leq M.$$
\item[(3)] For any fixed $(\omega,t)\in \Omega\times [0,T]$, $L(\omega,t,\cdot),R(\omega,t,\cdot)$ are strictly increasing and there  exist two constants $c,C$ satisfying $0<c<C$ such that for any $x,y\in \mathbb{R}$,
\begin{align*}
&c\left|x-y\right|\leq \left|L(\omega,t,x)-L(\omega,t,y)\right|\leq C\left|x-y\right|,\\
&c\left|x-y\right|\leq \left|R(\omega,t,x)-R(\omega,t,y)\right|\leq C\left|x-y\right|;
\end{align*}
\item[(4)] $\inf_{\omega,t,x} \left(R(\omega,t,x)-L(\omega,t,x)\right)>0$.
\end{itemize}
\end{assumption}
Compared with the doubly mean reflected BSDEs studied in \cite{L}, the generator $f$ in \eqref{nonlinearyz} depends not only on the distribution of the solution, but also on the compensating term $K$, which is required to satisfy the following assumption.
\begin{assumption}\label{ass1ipterminalp}
\begin{itemize}
      \item[(i)]  The terminal value $\xi \in \mathcal{L}^{p}_T $ satisfies $\mathbf{E}[L(T,\xi)]\leq 0 \leq \mathbf{E}[R(T,\xi)]$; 
      \item[(ii)] The generator $f$: $[0,T]\times\Omega\times \mathbb{R}\times \mathcal{P}_1(\mathbb{R})\times\mathbb{R}\times \mathcal{P}_1(\mathbb{R})\times \mathbb{R}$ $\rightarrow$ $\mathbb{R}$ is a $\mathcal{P}\times\mathcal{B}(\mathcal{P}_1(\mathbb{R}))\times\mathcal{B}(\mathbb{R})\times\mathcal{B}(\mathcal{P}_1(\mathbb{R}))\times\mathcal{B}(\mathbb{R})$-measurable map, and there exists a constant $\lambda>0$ such that for all $t\in[0,T]$, $y, y', k, k' \in \mathbb{R}$, $\mu, \mu' \in \mathcal{P}_1(\mathbb{R})$, $z, z'\in \mathbb{R}$, $\nu, \nu'\in\mathcal{P}_1(\mathbb{R})$, 
\begin{align*}
\left|f(t,y,\mu,z,\nu,k)-f(t,y',\mu',z',\nu',k')\right|\leq \lambda\left(\left|y-y'\right|+\left|z-z'\right|+W_1(\mu,\mu')+W_1(\nu,\nu')+\left|k-k'\right|\right).
\end{align*}
   \end{itemize}
\end{assumption}

For each $t\in[0,T]$, the resistance function $G_t:  BV[0,T]\rightarrow \mathbb{R}$ satisfies the following conditions.
\begin{assumption}\label{assgk}
For any $t\in[0,T]$, given $y\in BV[0,T]$,  we have $G_{t}(y)=G_{t}\left ( \left \{ y_{s\wedge t} \right \}_{0\le s\le t\le T}  \right ) $ and  $G_t(0)=0$. Moreover, there exists a constant $\lambda>0$, such that for any $y,y'\in BV[0,T]$,   $$\left | G_{t}(y)-G_t(y') \right | \leq \lambda \sup _{u\in[0,t]}\left | y_u-y'_{u} \right |.$$
\end{assumption}

\begin{example}
  %Assumption $\ref{ass-2}$ is reasonable. For instance, we can consider $G_{t}(y)=\int_{0}^{t}y_sds$.  
  We give some simple  examples of  $G_t(\cdot)$ satisfying Assumption \ref{assgk}.

  \begin{itemize}
      \item[(i)] $G_t(y)=y_t$.
      \item[(ii)] $G_t(y)=\sup_{0\leq s\leq t}y_s$.
      \item[(iii)] $G_{t}(y)=\int_{0}^{t}y_sds$.
  \end{itemize}
\end{example}

The construction of the solution to doubly mean reflected MFBSDE with nonlinear resistance \eqref{nonlinearyz} relies on the one of the solution to doubly mean reflected BSDE with constant generator as listed below. %Focus on the $\text{BSDE}$ with double mean reflections whose coefficient does not depend on $(Y,Z)\in \mathcal{S}^p\times \mathcal{H}^{p,d}$, we have the following proposition.
\begin{proposition}\label{pro-1}
  Let  Assumption \ref{ass1ipterminalp} and \ref{assLR} hold. Given $C\in \mathcal{H}^p$, the BSDE with double mean reflections 
   \begin{equation}\label{MFBSDEC}
   \begin{cases}
Y_t=\xi+\int_t^T C_s ds-\int_t^T Z_s dB_s+K_T-K_t, \\
\E[L(t,Y_t)]\leq 0\leq \E[R(t,Y_t)], \\
K_t=K^{R}_t-K^{L}_t, \int_0^T \E[R(t,Y_t)]dK_t^{R}=\int_0^T \E[L(t,Y_t)]dK^{L}_t=0,
\end{cases}
\end{equation}
has a unique solution $\left(Y,Z,K\right)\in \mathcal{S}^p\times\mathcal{H}^{p}\times BV[0,T]$.
\end{proposition}
\begin{proof}
 The proof is similar to the one for Proposition 3.4 in \cite{LS}. For the purpose of the remaining proof, we briefly describe the construction here.
 
Let $(\tilde{Y},Z)\in \mathcal{S}^p\times \mathcal{H}^p$ be the solution to the BSDE with terminal value $\xi$ and constant generator $C$. For any $t\in[0,T]$, set 
 $$\begin{aligned}
%  &\tilde{Y}_t=\xi+\int_{t}^{T}C_s ds-(M_T-M_t)\\
  & s_t=\E\left[\int_{0}^{t}C_s ds\right], a=\E\left[\xi\right]
 \end{aligned}$$
% where $M_t=\E_t\left[\xi+\int_{0}^{T}C_s ds\right]-\E\left[\xi+\int_{0}^{T}C_s ds\right]$. Clearly, we have $s\in C$ and $M$ is a p-integrable martingale. Therefore, there exits some $Z\in \mathcal{H}^{p,1}$, such that  $M_t=\int_{0}^{t}Z_s d B_s$.\\
and for any $(t,x)\in[0,T]\times \mathbb{R}$, we define 
$$l(t,x):=\E\left[L\left(t,\tilde{Y}_t-\E[\tilde{Y}_t]+x\right)\right],r(t,x):=\E\left[R\left(t,\tilde{Y}_t-\E[\tilde{Y}_t]+x\right)\right].$$
It is easy to check that $s$ is a continuous function  and $l,r$ satisfy Assumption \ref{asslr}. By Theorem \ref{BSP}, the backward Skorokhod problem $\mathbb{BSP}_l^r(s,a)$ admits a unique solution $(x,K)$. Set$$Y_t=\tilde{Y}_t+K_T-K_t=\xi+\int_{t}^{T}C_s ds-\int_{t}^{T}Z_s d B_s+K_T-K_t.$$ 
Then, $(Y,Z,K)$ is the solution to \eqref{MFBSDEC}.
\end{proof}



Now, we are ready to present the main result in this section.
\begin{theorem}\label{thlip}
 Let Assumptions $\ref{ass1ipterminalp}-\ref{assgk}$ hold. Then,  the doubly mean reflected MFBSDE with nonlinear resistance \eqref{nonlinearyz} admits a unique solution $(Y,Z,K)\in \mathcal{S}^p \times \mathcal{H}^{p} \times BV[0,T].$  
\end{theorem}

We apply a contraction mapping argument to prove Theorem \ref{thlip}. For this purpose, for each fixed $(y,z,k)\in \mathcal{S}^p \times \mathcal{H}^{p} \times BV[0,T]$, according to Proposition \ref{pro-1},   the following BSDE with double mean reflections 
\begin{equation}\label{yzk}
\begin{cases}
Y_t^{y,z,k}=\xi+\int_t^T f(s,y_s,\P_{y_s},z_s,\P_{z_s}, G_s(k))ds-\int_t^T Z_s^{y,z,k} dB_s+K_T^{y,z,k}-K_t^{y,z,k}, \\
\E[L(t,Y_t^{y,z,k})]\leq 0\leq \E[R(t,Y_t^{y,z,k})], \\
K^{y,z,k}_t=K^{y,z,k,R}_t-K^{y,z,k,L}_t, \\
\int_0^T \E[R(t,Y_t^{y,z,k})]dK_t^{y,z,k,R}=\int_0^T \E[L(t,Y_t^{y,z,k})]dK^{y,z,k,L}_t=0
\end{cases}
\end{equation}
has a unique solution $(Y^{y,z,k},Z^{y,z,k},K^{y,z,k})\in \mathcal{S}^p \times \mathcal{H}^{p} \times BV[0,T]$. Then, we define the  map
\begin{align*}
\Gamma:\mathcal{S}^p\times\mathcal{H}^{p}\times BV[0,T]&\rightarrow \mathcal{S}^p\times\mathcal{H}^{p}\times BV[0,T],  \\
\Gamma(y^i,z^i,k^i)&=(Y^i,Z^i,K^i).
\end{align*}
%To obtain the well-posedness of mean doubly reflected MFBSDE \eqref{nonlinearyz} with nonlinear resistance in the general time horizon $[0,T]$, we need to introduce the following lemma, which states that $\Gamma$ is a contraction mapping. 
We first show that $\Gamma$ is a contraction mapping  when $T$ is sufficiently small. 
\begin{lemma}\label{le3-2}
Let $\varepsilon$ satisfy $\beta(\lambda,c,C,p)\max\left(\varepsilon^{\frac{p}{2}},\varepsilon^p\right)<1$, where $\beta$ is a positive constant depending only on $\lambda, C,c,p$. For $T\in (0,\varepsilon]$, the solution map $\Gamma$ is contractive. 
\end{lemma}

\begin{proof}
In this proof,  $\beta$ always represents a positive constant depending on $\lambda,C,c,p$, which may vary from line to line. Let $(Y^i,Z^i,K^i):=(Y^{y^i,z^i,k^i},Z^{y^i,z^i,k^i},K^{y^i,z^i,k^i})$ be the solution to BSDE  with double mean reflections \eqref{yzk} associated with data $(y^i,z^i,k^i)$, $i=1,2$.

Set
\begin{align*}
&\hat{F}_t=F^1_t-F^2_t, \textrm{ where } F=Y,Z,K,y,z,k, \\ &\hat{f}_t=f(t,y^1_t,\P_{y^1_t},z^1_t,\P_{z^1_t},G_t(k^1))-f(t,y^2_t,\P_{y^2_t},z^2_t,\P_{z^2_t},G_t(k^2)).
\end{align*}
It is easy to check that 
\begin{align*}
|\hat{Y}_t|&=\left|\E_t\left[\int_t^T \hat{f}_sds\right]+\hat{K}_T-\hat{K}_t\right|\\
&\leq \lambda \E_t\left[\int_0^T \left(|\hat{y}_s|+|\hat{z}_s|+\E[|\hat{y}_s|]+\E[|\hat{z}_s|]+\sup_{r\in[0,s]}|\hat{k}_r|\right)ds \right]+|\hat{K}_T|+|\hat{K}_t|.
\end{align*}
Applying the Doob's maximal inequality, we obtain that  
\begin{equation}\begin{split}\label{differY}
 \E\left[\sup_{t\in[0,T]}|\hat{Y}_t|^p\right]
 &\leq \beta \E\left[\left(\int_{0}^{T}|\hat{y}_s|+|\hat{z}_s|+\E[|\hat{y}_s|]+\E[|\hat{z}_s|]+\sup_{r\in[0,s]}|\hat{k}_r|ds\right)^p\right]+\beta \sup_{t\in[0,T]}|\hat{K}_t|^p.\\   
\end{split}
\end{equation}
Recalling the proof of Proposition \ref{pro-1} and \eqref{diffK}, we have
\begin{align*}
\sup_{t\in[0,T]}|\hat{K}_t|\leq  \beta \left\{\sup_{t\in[0,T]}|\hat{s}_t|+\sup_{(t,x)\in[0,T]\times\mathbb{R}}|\hat{l}(t,x)|\vee \sup_{(t,x)\in[0,T]\times\mathbb{R}} |\hat{r}(t,x)|\right\},
\end{align*}
where $\hat{s}_t=s^1_t-s^2_t$, $\hat{l}(t,x)=l^1(t,x)-l^2(t,x)$, $\hat{r}(t,x)=r^1(t,x)-r^2(t,x)$ and for $i=1,2$, 
\begin{align*}
&s^i_t=\E\left[\int_0^t f\left(s,y^i_s,\P_{y^i_s},z^i_s,\P_{z^i_s},G_s(k^i)\right)ds\right], \\ 
&l^i(t,x)=\E\left[L\left(t,\widetilde{Y}^i_t-\E[\widetilde{Y}^i_t]+x\right)\right],\\  &r^i(t,x)=\E\left[R\left(t,\widetilde{Y}^i_t-\E[\widetilde{Y}^i_t]+x\right)\right].
\end{align*}
Here, $\widetilde{Y}^i$ is the first component of the solution to the $\text{BSDE}$ with terminal value $\xi$ and constant driver $f^i$, $i=1,2$, where $$f^i_s=f\left(s,y_s^i,\P_{y_s^i},z_s^i,\P_{z^i_s},G_s(k^i)\right).$$ By Assumptions \ref{assLR} and \ref{assgk}, noting that $\widetilde{Y}^i_t=\E_t\left[\xi+\int_t^T f\left(s,y_s^i,\P_{y_s^i},z_s^i,\P_{z^i_s},G_s(k^i)\right)d s\right]$, we obtain that 
\begin{align*}
\left|l^1(t,x)-l^2(t,x)\right|&\leq C\E\left[\left|\left(\widetilde{Y}^1_t-\E[\widetilde{Y}^1_t]\right)-\left(\widetilde{Y}^2_t-\E[\widetilde{Y}^2_t]\right)\right|\right]\leq 2C\E[|\widetilde{Y}^1_t-\widetilde{Y}^2_t|]\\
&\leq \beta \E\left[\int_0^T |\hat{y}_s|+|\hat{z}_s|+\E[|\hat{y}_s|]+\E[|\hat{z}_s|]+\sup_{r\in[0,s]}|\hat{k}_r|d s\right ].
\end{align*}
Similar estimates hold for $|s^1_t-s^2_t|$ and $|r^1(t,x)-r^2(t,x)|$. 
Applying H{\"o}lder inequality, the above analysis yields that 
\begin{align}\label{differK}
\sup_{t\in[0,T]}|\hat{K}_t|^p\leq \beta\E\left[\left(\int_0^T |\hat{y}_s|+|\hat{z}_s|+\E[|\hat{y}_s|]+\E[|\hat{z}_s|]+\sup_{r\in[0,s]}|\hat{k}_r|ds\right)^p \right].
\end{align}
Combining Eqs. \eqref{differY} and \eqref{differK} yields that 
\begin{align*}
\E\left[\sup_{t\in[0,T]}|\hat{Y}_t|^p\right]\leq \beta\E\left[\left(\int_0^T |\hat{y}_s|+|\hat{z}_s|+\E[|\hat{y}_s|]+\E[|\hat{z}_s|]+\sup_{r\in[0,s]}|\hat{k}_r|d s \right)^p\right].
\end{align*}
Note that we have
\begin{align*}
\int_0^t \hat{Z}_s dB_s=\int_0^t\hat{f}_s ds+\hat{Y}_t-\hat{Y}_0+\hat{K}_t.
\end{align*}
Simple calculation yields that
\begin{align*}
\E\left[\sup_{t\in[0,T]}\left|\int_0^t \hat{Z}_sd B_s\right|^p\right]&\leq \beta\E\left[\sup_{t\in[0,T]}\left|\int_0^t\hat{f}_s ds\right|^p+\sup_{t\in[0,T]}|\hat{Y}_t|^p+\sup_{t\in[0,T]}|\hat{K}_t|^p\right]\\
&\leq \beta\E\left[\left(\int_0^T |\hat{y}_s|+|\hat{z}_s|+\E[|\hat{y}_s|]+\E[|\hat{z}_s|]+\sup_{r\in[0,s]}|\hat{k}_r| ds\right)^p\right],
\end{align*}
Combining with the B-D-G inequality yields that 
\begin{align*}
   \E\left[\left(\int_0^T |\hat{Z}_s|^2 d s\right)^{\frac{p}{2}}\right]&\leq \beta\E\left[\left(\int_0^T |\hat{y}_s|+|\hat{z}_s |+\E[|\hat{y}_s|]+\E[|\hat{z}_s|]+\sup_{r\in[0,s]}|\hat{k}_r| ds\right)^p\right].
\end{align*}
By the H\"{o}lder inequality and the Fubini theorem, we finally deduce that 
\begin{align*}
&\E\left[\sup_{s\in[0,T]}|\hat{Y}_s|^p+\left(\int_0^T |\hat{Z}_s|^2 d s\right)^{\frac{p}{2}}+\sup_{t\in[0,T]}|\hat{K}_t|^p\right]\\
\leq &\beta\E\left[\left(\int_0^T \left(|\hat{y}_s|+|\hat{z}_s|+\E[|\hat{y}_s|]+\E[|\hat{z}_s|]+\sup_{r\in[0,s]}|\hat{k}_r|\right)ds\right)^p\right]\\
\leq &\beta\E\left[\left(\int_0^T |\hat{y}_s|+|\hat{z}_s|+\E[|\hat{y}_s|]+\E[|\hat{z}_s|]ds\right)^p+\left(\int_0^T\sup_{r\in[0,s]}|\hat{k}_r|ds\right)^p\right]\\
\leq &\beta\max\left(T^{\frac{p}{2}},T^p\right)\E\left[\sup_{t\in[0,T]}|\hat{y}_t|^p+\left(\int_0^T|\hat{z}_s|^2 ds\right)^{\frac{p}{2}}\right]+ T^p \sup_{r\in[0,T]}|\hat{k}_r|^p\\
\leq &\beta\max\left(T^{\frac{p}{2}},T^p\right)\E\left[\sup_{t\in[0,T]}|\hat{y}_t|^p+\left(\int_0^T|\hat{z}_t|^2 dt\right)^{\frac{p}{2}}+\left(\sup_{t\in[0,T]}|\hat{k}_t|\right)^p\right].
\end{align*}
Therefore, we can choose an appropriate  $T \in(0,\varepsilon]$ such that
$$
\begin{aligned}
  &\left\|Y^{1}-Y^{2}\right\|_{\mathcal{S}^{p}}+\left\|Z^{1}-Z^{2}\right\|_{\mathcal{H}^{p}}+\sup_{t\in[0,T]}\left|K_{t}^{1}-K_{t}^{2}\right|\\
  &\leq \beta\max\left(T^{\frac{p}{2}},T^p\right) \left(\left\|y^1-y^2\right\|_{\mathcal{S}^{p}}+\left\|z^1-z^2\right\|_{\mathcal{H}^{p}}+\sup_{t\in[0,T]}\left|k_t^1-k_t^2\right|\right). %\beta\max\left(T^{\frac{p}{2}},T^p\right)
\end{aligned}
$$
The proof is complete.
\end{proof}



Now,  we are ready to give the proof of Theorem \ref{thlip}. 

\begin{proof}[Proof of Theorem \ref{thlip}] First, for a given $k\in BV[0,T]$, according to Theorem 3.5 in \cite{LS},  the following MFBSDE with double mean reflections
 \begin{equation}\label{MFBSDEDMRkL}
\begin{cases}
Y_t=\xi+\int_t^T f(s,Y_s,\mathbf{P}_{Y_s},Z_s,\mathbf{P}_{Z_s},G_s(k))ds-\int_t^T Z_s dB_s+K_T-K_t, \\
\mathbf{E}[L(t,Y_t)]\leq 0\leq \mathbf{E}[R(t,Y_t)], \\
K_t=K^R_t-K^L_t,\ K^R,K^L\in I[0,T],\\
\int_0^T \mathbf{E}[R(t,Y_t)]dK_t^R=\int_0^T \mathbf{E}[L(t,Y_t)]dK^L_t=0.\\
\end{cases}
\end{equation}
admits a unique solution $(Y,Z,K)\in\mathcal{S}^p \times \mathcal{H}^{p} \times BV[0,T]$. 

Now, let $N$ be the smallest integer such that $N\geq \frac{T}{\varepsilon}$, where $\varepsilon$ is be chosen as Lemma \ref{le3-2}. For $k^i\in BV[0,T],\ i=1,2$, let $(Y^i,Z^i,K^i)$ be the unique solution to \eqref{MFBSDEDMRkL} with data $k^i$. According to Lemma \ref{le3-2}, for any $1\leq i\leq N$, 
$$
\begin{aligned}
  &\left\|Y^{1}-Y^{2}\right\|_{\mathcal{S}^{p}{[(i-1)\varepsilon,i\varepsilon \wedge T]}}+\left\|Z^{1}-Z^{2}\right\|_{\mathcal{H}^{p}[(i-1)\varepsilon,i\varepsilon \wedge T]}+\sup_{t\in[(i-1)\varepsilon,i\varepsilon \wedge T]}\left|K_{t}^{1}-K_{t}^{2}\right|\\
  &\leq \beta\max\left(\varepsilon^{\frac{p}{2}},\varepsilon^p\right) \left(\left\|Y^{1}-Y^{2}\right\|_{\mathcal{S}^{p}[(i-1)\varepsilon,i\varepsilon \wedge T]}+\left\|Z^{1}-Z^{2}\right\|_{\mathcal{H}^{p}[(i-1)\varepsilon,i\varepsilon \wedge T]}+\sup_{t\in[(i-1)\varepsilon,i\varepsilon \wedge T]}\left|k_t^1-k_t^2\right|\right).
\end{aligned}
$$
Therefore, we have 
$$\sup_{t\in[(i-1)\varepsilon,i\varepsilon \wedge T]}\left|K_t^1-K_t^2\right|\leq \max_{1\leq i \leq N}\left(\sup_{t\in[(i-1)\varepsilon,i\varepsilon \wedge T]}\left|K_t^1-K_t^2\right|\right)\leq \beta\max\left(\varepsilon^{\frac{p}{2}},\varepsilon^p\right)\sup_{t\in[0,T]}\left|k_t^1-k_t^2\right|.$$
By the principle of standard contraction mapping, the doubly mean reflected MFBSDE \eqref{nonlinearyz} with resistance  admits a unique solution $(Y,Z,K)\in\mathcal{S}^p \times \mathcal{H}^{p} \times BV[0,T].$   
\end{proof}

\section{Doubly mean reflected MFBSDEs with nonlinear resistance: the quadratic case}	
In this section, we investigate the well-posedness of doubly mean reflected \text{MFBSDE} with nonlinear resistance \eqref{nonlinearyz} when its generator $f$ has quadratic growth  and the terminal value is bounded. More precisely, we propose the  following assumption.
\begin{assumption}\label{assqua}
\begin{itemize}
      \item[(i)]  The terminal value $\xi \in \mathcal{L}^{\infty}_T $ satisfies $\mathbf{E}[L(T,\xi)]\leq 0 \leq \mathbf{E}[R(T,\xi)]$, i.e. there exist some constants $H_1$ such that $\left \| \xi  \right \|_{\mathcal{L}^\infty}\leq H_1$; 
      \item[(ii)] The process $\left\{f\left(t, 0, \delta_0, 0,\delta_0,0 \right)\right\}_{t\in[0,T]}$ is uniformly bounded by some constant $H_2$, $\mathbb{P}$-a.s. There exist some positive constants $\lambda$ and $\alpha\in[0,1)$ such that, $\mathbb{P}$-a.s., for any $t \in[0, T]$, $ y_1, y_2 , k_1, k_2\in \mathbb{R}$, $\mu_1, \mu_2 \in \mathcal{P}_1(\mathbb{R})$, $z_1, z_2 \in \mathbb{R}$ and $\nu_1, \nu_2 \in \mathcal{P}_1(\mathbb{R})$, we have
	$$
    \begin{aligned}
		& \left|f\left(t, y_1, \mu_1, z_1,\nu_1, k_1\right)-f\left(t, y_2, \mu_2, z_2, \nu_2, k_2\right)\right| \\
		\leq &\lambda\left(\left|y_1-y_2\right|+\left(1+\left|z_1\right|+\left|z_2\right|\right)\left|z_1-z_2\right|+\left|k_1-k_2\right|\right)\\
        & +\lambda \mathcal{W}_1\left(\mu_1,\mu_2\right)+\lambda \left(1+\left(\mathcal{W}_1\left(\nu_1,\delta_0\right)\right)^{\alpha}+\left(\mathcal{W}_1\left(\nu_2,\delta_0\right)\right)^{\alpha}\right)\mathcal{W}_1\left(\nu_1,\nu_2\right).
	\end{aligned}
	$$ 
   \end{itemize}
 \end{assumption}
 
\begin{remark}
\begin{itemize}
    \item[(i)] We will take $H=H_1 \vee H_2$ in the following text.
    \item[(ii)]% We use the following representation for $2$-Wasserstein distance
%\begin{align*}
%W_2(\mu,\nu)=\inf_{X \sim\mu, Y\sim \nu}\E[\left|X-Y\right|^2]^{\frac{1}{2}},
%\end{align*}
 Applying the H\"{o}lder inequality, we have $$W_1(\mu,\nu)\leq W_2(\mu,\nu).$$
\end{itemize}
\end{remark}

First, we state the main result in this section.
\begin{theorem}\label{globalMFBSDE}
Let Assumptions \ref{assLR}, \ref{assgk}, \ref{assqua} hold. Suppose that for any $(t,y,\mu,\nu,k)\in [0,T]\times\mathbb{R}\times\mathcal{P}_1(\mathbb{R})\times\mathcal{P}_1(\mathbb{R})\times\mathbb{R}$, there exists a constant $\tilde{H}>0$, such that $\left|f(t,y,\mu,0,\nu,k)\right|\leq \tilde{H}$. Then, for any time horizon $[0,T]$, doubly mean reflected MFBSDE  with nonlinear resistance \eqref{nonlinearyz} has a unique solution $ (Y,Z,K)\in\mathcal{S}^{\infty}\times BMO \times BV[0,T].$     
\end{theorem}

The construction of the solution to doubly mean reflected MFBSDE with nonlinear resistance needs a contraction mapping argument. Due to the fact that $f$ has quadratic growth in $z$, the coefficient depends on the uniform estimate of $Y$ (see $\tilde{A}$ in Eq. \eqref{contraction}, where $\tilde{A}$ appears in \eqref{BA}). Therefore, we first provide the following result concerning the uniform estimate of the first component of the solution to a doubly mean reflected MFBSDE with an additional bounded assumption for the generator.
%To obtain the well-posedness of doubly mean reflected MFBSDE with nonlinear resistance for arbitrary time horizon, we need to make sure that the first component $Y$ of any local solution $(Y,Z,K)$ on $[T-h,T]$ has a uniform estimate. This property requires an additional bounded condition for the generator $f$.
\begin{lemma}\label{lem4.6}
Let Assumptions \ref{assLR}, \ref{assgk}, \ref{assqua} hold. Suppose that the generator $f$ does not depend on $k$, and for any $(t,y,\mu,\nu)\in [T-h,T]\times\mathbb{R}\times\mathcal{P}_1(\mathbb{R})\times\mathcal{P}_1(\mathbb{R})$, there exists a constant $\hat{H}$ such that $\left|f(t,y,\mu,0,\nu)\right|\leq \hat{H}$. Moreover, assume that the doubly mean reflected MFBSDE \eqref{nonlinearyz} has a local solution $(Y,Z,K)\in \mathcal{S}^{\infty}[T-h,T]\times BMO[T-h,T]\times BV[T-h,T]$ for some $0<h\leq T$. Then, there exists a constant $\bar{H}$ depending only on $C,c,H,\hat{H},\lambda, T,M$ such that 
$$\left\|Y\right\|_{\mathcal{S}^{\infty}[T-h,T]}\leq \bar{H}.$$
\end{lemma}
\begin{proof}
 First, according to Theorem 3.11 in \cite{HHTW}, the following  quadratic MFBSDE 
\begin{equation}\label{yz}
 \bar{Y}_{t}=\xi+\int_t^{T} f(s,Y_s,\P_{Y_s},\bar{Z}_{s}, \P_{\bar{Z}_{s}})d s-\int_t^{T} \bar{Z}_{s} dB_s , 
 \end{equation}
 admits a unique solution $(\bar{Y},\bar{Z})\in\mathcal{S}^{\infty}[T-h,T]\times BMO[T-h,T]$. Therefore, similar as the proof of Lemma \ref{4.1.1}, $Y$ has the following representation:
$$Y_{t}=\bar{Y}_{t}+K_{T}-K_{t}=\bar{Y}_{t}+\bar{K}_{T-t},\ \forall t\in[T-h,T],$$
and $\bar{K}$ is the second component of the solution to the Skorokhod problem $\mathbb{SP}_{\bar{l}}^{\bar{r}}(\bar{s})$, which satisfies $K_t=\bar{K}_{h}-\bar{K}_{T-t}$. 
Here for each $t\in[T-h,T]$, 
\begin{equation}\label{salry}\begin{split}
&\bar{s}_t=\E\left[\bar{Y}_{T-t}\right],\\
&\bar{l}(t,x):=\E\left[L(T-t,\bar{Y}_{T-t}-\E\left[\bar{Y}_{T-t}\right]+x)\right],\\
&\bar{r}(t,x):=\E\left[R(T-t,\bar{Y}_{T-t}-\E\left[\bar{Y}_{T-t}\right]+x)\right].\\
\end{split}\end{equation}
Furthermore, 
\begin{equation}\begin{split}
\left\|Y\right\|_{\mathcal{S}^{\infty}[T-h,T]}&\leq\left\|\bar{Y}\right\|_{\mathcal{S}^{\infty}[T-h,T]}+\sup_{t\in[T-h,T]}\left|\bar{K}_{T-t}\right|\\
&\leq\left\|\bar{Y}\right\|_{\mathcal{S}^{\infty}[T-h,T]}+\sup_{t\in[T-h,T]}\left|\bar{K}_{T-t}-\bar{K}^0_{T-t}\right|+\sup_{t\in[T-h,T]}\left|\bar{K}^0_{T-t}\right|,\\
\end{split}   
\end{equation}
where $\bar{K}^0$ is the second component of the solution to the Skorokhod problem $\mathbb{SP}^{\bar{r}^0}_{\bar{l}^0}(\bar{s}^0)$ and $\bar{s}^0$, $\bar{r}^0$, $\bar{l}^0$ are the same as in  \eqref{slr}. In view of Proposition \ref{continuity}, we derive that 
$$
\begin{aligned}
\sup_{t \in[T-h, T]}\left|\bar{K}_{T-t}-\bar{K}^0_{T-t}\right|&=\sup_{t \in[0, h]}\left|\bar{K}_{t}-\bar{K}^0_{t}\right|\\
&\leq \frac{3C}{c}\sup_{t \in[T-h, T]}\mathbf{E}\left[\left|\bar{Y}_{t}\right|\right]+\frac{C}{c}\mathbf{E}\left[\left|\xi\right|\right]\\
&\leq \frac{3C}{c}\left\|\bar{Y}\right\|_{\mathcal{S}^{\infty}[T-h, T]}+\frac{C}{c}H.\\  
\end{aligned}		
$$
Recalling \eqref{barK0}, we have
\begin{equation*}
 \sup_{t\in[T-h,T]}\left|\bar{K}^0_{T-t}\right|\leq \sup_{t\in[0,T]}\left|\bar{K}^0_{t}\right|\leq \frac{2C}{c}\mathbf{E}[\left|\xi\right|]+\frac{M}{c}\leq \frac{2C}{c}H+\frac{M}{c}.   
\end{equation*}
According to Proposition 2.2 in \cite{BE}, it is easy to check that for each $t\in[T-h,T]$,
$$\left|\bar{Y}_t\right|\leq(H+\hat{H}T)e^{\lambda T}:=\bar{H}^1.$$
%and
%$$\left\|Z\right\|_{BMO[T-h,T]}^2 \leq \frac{1\vee T}{3}\left(1+\frac{4\hat{H}}{\lambda}+2\bar{H}^1\right)e^{3\lambda\bar{H}^1}:=\bar{H}^2.$$
Therefore,
\begin{equation}\label{barH}
\left\|Y\right\|_{\mathcal{S}^{\infty}[T-h,T]}\leq(1+\frac{3C}{c})\bar{H}^1+\frac{3C}{c}H+\frac{M}{c}:=\bar{H}, \ \forall t\in[T-h,T],   
\end{equation}
which is the desired result.
\end{proof}

Next, when $T$ is sufficiently small, we show that the quadratic MFBSDE with double mean reflections and nonlinear resistance \eqref{nonlinearyz} admits a unique solution  by utilizing a contraction mapping method. For this purpose, given $
\left(y, z, k\right)\in \mathcal{S}^{\infty} \times BMO\times BV[0,T]$, consider the following quadratic MFBSDE with double mean reflections  %Under Assumptions \ref{assLR}, \ref{assgk}, \ref{assqua}, for each $\left(y, z, k\right)\in \mathcal{S}^{\infty} \times \text{BMO}\times BV[0,T]$, it follows from Theorem 4.2 in \cite{LS} that following quadratic MFBSDE with double mean reflections 
\begin{equation}\label{nonlinearyz-1}
\begin{cases}
Y_t^{y,z,k}=\xi+\int_t^T f(s,y_s,\P_{y_s},Z_s^{y,z,k},\P_{z_s},G_s(k))ds-\int_t^T Z_s^{y,z,k} dB_s+K_T^{y,z,k}-K_t^{y,z,k},\\
\E[L(t,Y_t^{y,z,k})]\leq 0\leq \E[R(t,Y_t^{y,z,k})], \\
K_t^{y,z,k}=K^{y,z,k,R}_t-K^{y,z,k,L}_t,\ K^{y,z,k,R},K^{y,z,k,L}\in I[0,T],\\
\int_0^T \E[R(t,Y_t)]dK^{y,z,k,R}_t=\int_0^T \E[L(t,Y_t)]dK^{y,z,k,L}_t=0.
\end{cases}
\end{equation}  
Then, under Assumptions \ref{assLR}, \ref{assgk}, \ref{assqua}, it follows from Theorem 4.2 in \cite{LS} that the above equation has a unique solution $\left(Y^{y,z,k}, Z^{y,z,k}, K^{y,z,k}\right)\in \mathcal{S}^{\infty} \times BMO\times BV[0,T]$.     
Now, we define the map 
\begin{align*}
\Gamma:\mathcal{S}^{\infty} \times BMO\times BV[0,T]&\rightarrow \mathcal{S}^{\infty} \times BMO \times BV[0,T],  \\
\Gamma\left(y,z,k\right)&:=\left(Y^{y,z,k}, Z^{y,z,k}, K^{y,z,k}\right).
\end{align*}
We set 
\begin{equation}\label{BA}
    \mathscr{B}_{\tilde{A}} := \left\{ \left(y, z, k\right)\in \mathcal{S}^{\infty} \times BMO \times BV[0,T]:\left \| y \right \|_{\mathcal{S}^{\infty}}\leq \tilde{A},  \left \| z \right \|_{BMO}\leq \tilde{A}, \sup_{0 \leq t \leq T}\left|k_t\right|\leq \tilde{A}\right\},
\end{equation}
where $\tilde{A} \geq \tilde{A}_0$ and
\begin{equation}\label{A0}
 \tilde{A}_0=(3+\frac{24C}{c})H+\frac{2M}{c}+(\frac{2H}{\lambda}+\frac{1}{3})e^{9\lambda H}.   
\end{equation}


	\begin{theorem}\label{thmbddterminal}
	  Let Assumptions \ref{assLR}, \ref{assgk}, \ref{assqua} hold. Then, there exists a constant $\hat{\delta}^{\tilde{A}}:=\hat{\delta}^{\tilde{A}}(\tilde{A}, c, C, H, \lambda)>0$ such that for any $T\in(0,\hat{\delta}^{\tilde{A}}]$, the quadratic MFBSDE \eqref{nonlinearyz} with double mean reflections and nonlinear resistance  admits a unique solution $(Y, Z, K) \in \mathcal{S}^{\infty} \times BMO \times BV[0,T]$. Moreover, we have  
     $$\left \| Y \right \|_{\mathcal{S}^{\infty}}\leq \tilde{A}, \  \left \| Z \right \|_{BMO}\leq \tilde{A} \text{ and } \sup_{0 \leq t \leq T}\left|K_t\right|\leq \tilde{A}.$$
	\end{theorem}
    %an appropriate  constant $\tilde{A}>0$ {\color{red}what do you mean by saying appropriate? Does $\tilde{A}\geq \tilde{A}_0$ enough?} and 
%, \forall \left(y, z, k\right)\in \mathcal{S}^{\infty} \times \text{BMO}\times BV[0,T].$$

The proof of Theorem \ref{thmbddterminal} needs the following lemma. %To this purpose, we need to show that $\Gamma$ is a contraction mapping. 
%{\color{red}and $P$ is a constant, which will be explained in the following text.} 
\begin{lemma}\label{4.1.1}
Let Assumptions \ref{assLR}, \ref{assgk}, \ref{assqua} hold. Then, there exists a constant $\delta^{\tilde{A}}>0$ depending only on $\tilde{A}, H, \lambda, C, c$ such that for any $T\in(0,\delta^{\tilde{A}}]$, $\Gamma(\mathscr{B}_{\tilde{A}})\subset \mathscr{B}_{\tilde{A}}$. Furthermore, for any $T\in(0,\hat{\delta}^{\tilde{A
}}]$, $\Gamma$ is contractive.
\end{lemma}

\begin{proof}
\textbf{Step 1.} We first show that $\Gamma(\mathscr{B}_{\tilde{A}})\subset \mathscr{B}_{\tilde{A}}$. For notional simplification, set
$$Y_t:=Y^{y,z,k}_t, Z_t:=Z^{y,z,k}_t,K:=K^{y,z,k}_t,f^{y,z,k}(t,z):=f(t,y_t,\P_{y_t},z,\P_{z_{t}},G_{t}(k)).$$
Similar with the proof of Proposition \ref{pro-1}, we have 
 \begin{equation}\label{YmZm}
   Y_{t}=\bar{Y}_{t}+K_{T}-K_{t},\ Z_t=\bar{Z}_{t},
 \end{equation}
 where $(\bar{Y},\bar{Z})\in \mathcal{S}^{\infty}\times \text{BMO}$ is the unique solution to quadratic BSDE with terminal condition $\xi$ and the generator $f^{y,z,k}$, and $K$ is the second component of the solution to the backward Skorokhod problem $\mathbb{BSP}_{l}^{r}(s,a)$. Here, for each $t\in[0,T]$,  
\begin{equation}\begin{split}
&s_t=\E\left[\bar{Y}_{0}-\bar{Y}_{t}\right],\ a=\mathbf{E}[\xi],\\
&l(t,x)=\E\left[L(t,\bar{Y}_{t}-\E\left[\bar{Y}_{t}\right]+x)\right],\\
&r(t,x)=\E\left[R(t,\bar{Y}_{t}-\E\left[\bar{Y}_{t}\right]+x)\right].\\
\end{split}\end{equation}
In addition, according to Proposition 4.2 in \cite{L} and Remark \ref{rem2.7}, we know that the solution $K$ of \eqref{nonlinearyz-1} has the following representation 
\begin{equation}\label{kbark}
K_t=\bar{K}_{T}-\bar{K}_{T-t},
\end{equation}
where $\bar{K}$ is the second component of the solution to the Skorokhod problem $\mathbb{SP}_{\bar{l}}^{\bar{r}}(\bar{s})$.
%It is easy to verify that for each $t\in[0,T]$,$$(Y_t,Z_t)=(\bar{Y}_t+K_T-K_t,\bar{Z}_t),$$
%where $(\bar{Y}_t,\bar{Z}_t)\in\mathcal{S}^{\infty}\times \text{BMO}$ is the solution to the following standard BSDE on the time interval $[0,T]$
%\begin{equation}\label{Y}
 %\bar{Y}_t=\xi+\int^{T}_{t}f(s,y_s,\P_{y_s},\bar{Z}_s, \P_{z_s},G_s(k))ds-\int^{T}_{t}\bar{Z}_sdB_s,   
%\end{equation}
Here for each $t\in[0,T]$,  
\begin{equation}\label{salr}\begin{split}
\bar{s}_t&=\E\left[\bar{Y}_{T-t}\right],\\
\bar{l}(t,x)&=\E\left[L(T-t,\bar{Y}_{T-t}-\E\left[\bar{Y}_{T-t}\right]+x)\right],\\
\bar{r}(t,x)&=\E\left[R({T-t},\bar{Y}_{T-t}-\E\left[\bar{Y}_{T-t}\right]+x)\right].
\end{split}\end{equation}
%It is easy to check that $l,r$ satisfy Assumption \ref{asslr} and 
%\begin{align*}
%l(T,a)=\E\left[L(T,\xi)\right]\leq 0\leq \E\left[R(T,\xi)\right]=r(T,a).
%\end{align*}
%Therefore, by Theorem \ref{BSP}, the backward Skorokhod problem $\mathbb{BSP}_{l}^{r}(s,a)$ admits a unique solution $(x,K)$ and for any $t\in[0,T]$, $K_t=\bar{K}_T-\bar{K}_{T-t}$, where $\bar{K}$ is the second component of the solution to the Skorokhod problem $\mathbb{SP}_{\bar{l}}^{\bar{r}}(\bar{s})$. Here, for any $t\in[0,T]$, 
%\begin{equation}\label{barslr}\bar{s}_{t}=a+s_{T}-s_{T-t}=\mathbf{E}[\bar{Y}_{T-t}],\ \bar{l}(t,x)=l(T-t,x), \ \bar{r}(t,x)=r(T-t,x).
%\end{equation}
Combining Eqs. \eqref{YmZm} and \eqref{kbark}, we obtain that  
$$
Y_t=\bar{Y}_t+K_T-K_t=\bar{Y}_t+\bar{K}_{T-t},  \ t\in[0,T].
$$
Consequently, we have
\begin{equation}\label{re-1}
\begin{split}
 \left \| Y \right \|_{\mathcal{S}^{\infty}}&\leq \left \| \bar{Y} \right \|_{\mathcal{S}^{\infty}}+\sup_{t\in[0,T]}\left|\bar{K}_t\right|\\
 &\leq \left \| \bar{Y} \right \|_{\mathcal{S}^{\infty}}+\sup_{t\in[0,T]}\left|\bar{K}_t-\bar{K}^0_t\right|+\sup_{t\in[0,T]}\left|\bar{K}^0_t\right|,   
\end{split}   
\end{equation}
where $\bar{K}^0$ is the second component of the solution to the Skorokhod problem $\mathbb{SP}^{\bar{r}^0}_{\bar{l}^0}(\bar{s}^0)$. Here, for each $t\in[0,T]$, 
\begin{equation}\label{slr}
\bar{s}^0_t=\E[\xi],\ \bar{l}^0(t,x)=\E\left[L(T-t,x)\right], \ \bar{r}^0(t,x)=\E\left[R(T-t,x)\right].
\end{equation}
Let $\phi_t,\psi_t$ be two constants satisfying the following equations, respectively,
\begin{equation*}\begin{split}
&\bar{l}^0(t,\bar{s}_t^0+ \phi_t)=0, \ \bar{r}^0(t,\bar{s}_t^0+ \psi_t)=0.\\
\end{split}\end{equation*} 
By Proposition \ref{SP}, we have 
\begin{equation}\label{barK0}\begin{split}
\sup_{t\in[0,T]}\left|\bar{K}^{0}_{t}\right|%&=\sup_{t\in[0,T]}\left|-\max\left((-\phi_0)^+\vee\inf_{r\in[0,t]}(-\psi_r),\sup_{s\in[0,t]}[(-\phi_s)]\vee\inf_{r\in[s,t]}(-\psi_r)\right)\right|\\
&\leq \sup_{t\in[0,T]}\left|\phi_{t}\right|+\sup_{t\in[0,T]}\left|\psi_{t}\right|\\
&\leq \frac{1}{c} \sup_{t\in[0,T]}\left|\bar{l}^0(t,\E[\xi]+ \phi_t)-\bar{l}^0(t,\E[\xi])\right|+ \frac{1}{c} \sup_{t\in[0,T]}\left|\bar{r}^0(t,\E[\xi]+ \psi_t)-\bar{r}^0(t,\E[\xi])\right|\\
&\leq \frac{1}{c} \sup_{t\in[0,T]}\left|\E[L(T-t, \E[\xi])-L(T-t,0)]\right|+ \frac{1}{c} \sup_{t\in[0,T]}\left|\E[L(T-t,0)]\right|\\
&\quad+\frac{1}{c} \sup_{t\in[0,T]}\left|\E[R(T-t, \E[\xi])-L(T-t,0)]\right|+ \frac{1}{c} \sup_{t\in[0,T]}\left|\E[R(T-t,0)]\right|\\
&\leq \frac{2C}{c}\mathbf{E}[\left|\xi\right|]+\frac{M}{c}.\\
\end{split}\end{equation}
 In view of Proposition \ref{continuity}, it is easy to check that % and the analysis of equation $(4.14)$ in \cite{LS} , we know
%$$
%	\sup_{t \in[0, T]}\left|\bar{K}_{t}-\bar{K}^0_{t}\right|\leq \frac{C}{c}\sup_{t\in[0,T]}\left|\bar{s}_t-\bar{s}^0_t\right|+\frac{1}{c}(\bar{L}^0_T \vee \bar{R}^0_T),	
%$$
%where $\bar{L}_T=\sup_{(t,x)\in[0,T]\times \mathbb{R}}|\bar{l}(t,x)-\b%ar{l}^0(t,x)|$, $
%\bar{R}_T=\sup_{(t,x)\in[0,T]\times \mathbb{R}}|\bar{r}(t,x)-\bar{r}^0(t,x)|$. Recalling that  $\bar{s}^0_t=\E[\xi],\bar{s}_{t}=\mathbf{E}[\bar{Y}_{T-t}]$, we have 
%\begin{align*}
%	\sup_{t\in[0,T]}|\bar{s}_t-\bar{s}^0_t|
 %   \leq \sup_{t\in[0,T]}\mathbf{E}[|\bar{Y}_{t}|]+\E[|\xi|]
	%\end{align*}
%Applying Assumption \ref{ass2}, we obtain that 
%$$
%\begin{aligned}
%	\bar{L}^0_T
%	&= \sup_{(t,x)\in[0,T]\times \mathbb{R}} \left|\mathbf{E}\left[L(t,\bar{Y}_{T-t}-\E[\bar{Y}_{T-t}]+x))\right]-\E\left[L(T-t,x)\right]\right|\\
%	&\leq \sup_{t\in[0,T]}C\mathbf{E}\left[\left|\bar{Y}_{T-t}-\E[\bar{Y}_{T-t}]\right|\right]\\
%	&\leq 2C\sup_{t \in[0, T]}\mathbf{E}[|\bar{Y}_t|].
%\end{aligned}
%$$
%Similarly, we have $\bar{R}_T\leq 2C\sup_{t \in[0, T]}\mathbf{E}[|\bar{Y}_t|]$. The above analysis implies that 
\begin{equation}\label{barKK0}
	\sup_{t \in[0, T]}\left|\bar{K}_{t}-\bar{K}^0_{t}\right|
	\leq \frac{3C}{c}\sup_{t \in[0, T]}\mathbf{E}[|\bar{Y}_t|]+\frac{C}{c}\E[|\xi|].
\end{equation}
Let us recall the following estimates, which can be found in the proof of Lemma 2.10 in \cite{LP} (see Eqs. (13) and (14) in \cite{LP})  
\begin{equation}\label{barYZ}\begin{split}
\left\|\bar{Y}\right\|_{\mathcal{S}^{\infty}}%&\leq (1+T)L+\lambda T(2\left\|y\right\|_{\mathcal{S}^{\infty}}+\sup_{t\in[0,T]}\left|k_t\right|)+\lambda \sqrt{T}\left\|z\right\|_{BMO}+\lambda T^{\frac{1-\alpha}{2}}\left\|z\right\|_{BMO}^{1+\alpha}\\
&\leq (1+T)H+3\lambda T \tilde{A}+\lambda\sqrt{T} \tilde{A}+\lambda T^{\frac{1-\alpha}{2}} \tilde{A}^{1+\alpha},\\
\left\|Z\right\|_{BMO}^2%&\leq\frac{2 e^{3\lambda\left\|\bar{Y}\right\|_{\mathcal{S}^{\infty}}}}{3} \left(\frac{(1+T)L}{\lambda}+ T(2\left\|y\right\|_{\mathcal{S}^{\infty}}+\sup_{0\leq t\leq T}\left|k_t\right|+\frac{1}{2})+ \sqrt{T}\left\|z\right\|_{BMO}+ T^{\frac{1-\alpha}{2}}\left\|z\right\|_{BMO}^{1+\alpha}\right)\\
&\leq \frac{2 e^{3\lambda\left\|\bar{Y}\right\|_{\mathcal{S}^{\infty}}}}{3}\left(\frac{(1+T)H}{\lambda}+3 T \tilde{A}+\frac{T}{2} +\sqrt{T}\tilde{A}+T^{\frac{1-\alpha}{2}}\tilde{A}^{1+\alpha}\right).
\end{split}\end{equation}

Define 
$$\delta^{\tilde{A}}:=\min\left(\frac{H}{9\lambda \tilde{A}},\frac{H^2}{9\lambda^2 \tilde{A}^2}, \left(\frac{H}{3\lambda \tilde{A}^{1+\alpha}}\right)^{\frac{2}{1-\alpha}}\right).$$
Therefore, we  can derive that for each $T\in(0,\delta^{\tilde{A}}]$, 
$$
\begin{aligned}
 &\left \| \bar{Y} \right \|_{\mathcal{S}^{\infty}}\leq 3H,\ \ \ \ \left \| Z \right \|_{BMO} \leq \sqrt{\left(\frac{2H}{\lambda}+\frac{1}{3}\right)e^{9\lambda H}}\leq \tilde{A},\\
 &\left \| Y \right \|_{\mathcal{S}^{\infty}}\leq \left(1+\frac{4C}{c}\right)3H+\frac{M}{c}\leq \tilde{A}.
\end{aligned} 
$$
Meanwhile, recalling \eqref{barK0} and \eqref{barKK0}, we have 
$$
\begin{aligned}
 \sup_{t\in[0,T]}\left|K_t\right|&=\sup_{t\in[0,T]}\left|\bar{K}_T-\bar{K}_{T-t}\right|\leq2\sup_{t\in[0,T]}\left|\bar{K}_t\right|\\
 &\leq2\sup_{t \in[0, T]}\left|\bar{K}_{t}-\bar{K}^0_{t}\right|+2\sup_{t \in[0, T]}\left|\bar{K}^0_{t}\right|\\
 &\leq \frac{24C}{c}H+\frac{2M}{c}\leq \tilde{A}.\\   
\end{aligned} 
$$
%Then recalling \eqref{barK0} and \eqref{barKK0}, we have 
%$$
%\begin{aligned}
% \sup_{t\in[0,T]}\left|K_t\right|&=\sup_{t\in[0,T]}\left|\bar{K}_T-\bar{K}_{T-t}\right|\leq2\sup_{t\in[0,T]}\left|\bar{K}_t\right|\\
 %&\leq2\sup_{t \in[0, T]}\left|\bar{K}_{t}-\bar{K}^0_{t}\right|+2\sup_{t \in[0, T]}\left|\bar{K}^0_{t}\right|\\
% &\leq \frac{24C}{c}H+\frac{2M}{c}\leq \tilde{A},\\   
%\end{aligned} 
%$$


\textbf{Step 2.} Now, we show that $\Gamma$ is a contraction mapping.  For each $i=1,2$, we denote 
$$
(Y^i,Z^i,K^i)=\Gamma(y^i,z^i,k^i), \ f^i(t,z)=f(t,y^i_{t},\P_{y^i_{t}},z,\P_{z^i_{t}},G_{t}(k^i)).
$$
Similar as the analysis in Step 1, for $i=1,2$, we have     
\begin{equation}\label{re-2}
Y^{i}_t=\bar{Y}^{i}_{t}+K^{i}_{T}-K^{i}_t, \quad \forall t\in[0,T],   
\end{equation} 
where $(\bar{Y}^i,\bar{Z}^i)\in\mathcal{S}^{\infty}\times BMO$ is the solution to BSDE with terminal condition $\xi$ and the generator $f^i$, and $K^i$ is the second component of the solution to the backward Skorokhod problem $\mathbb{BSP}_{l^i}^{r^i}(s^i,a^i)$. Here, for each $t\in[0,T]$,  
\begin{equation}\label{bsalri}\begin{split}
&s^i_t=\E\left[\bar{Y}^i_{0}-\bar{Y}^i_{t}\right],\ a^i=\mathbf{E}[\xi],\\
&l^i(t,x)=\E\left[L(t,\bar{Y}^i_{t}-\E\left[\bar{Y}^i_{t}\right]+x)\right],\\
&r^i(t,x)=\E\left[R(t,\bar{Y}^i_{t}-\E\left[\bar{Y}^i_{t}\right]+x)\right].\\
\end{split}\end{equation}
Therefore, we have
\begin{equation}\begin{split}
   \left\|Y^1-Y^2\right\|_{\mathcal{S}^{\infty}}%&=\left\|\bar{Y}^1-\bar{Y}^2+\bar{K}_{T-t}^1-\bar{K}_{T-t}^2\right\|_{\mathcal{S}^{\infty}}\\
   &\leq \left\|\bar{Y}^1-\bar{Y}^2\right\|_{\mathcal{S}^{\infty}}+2\sup_{t\in[0,T]}\left|K^1_t-K_t^2\right|.
\end{split}   
\end{equation}
Recall the following estimates from the proof 
of Lemma 2.11 in \cite{LP} (see Eqs. (23) and (24) in \cite{LP})
$$
\begin{aligned}
   \left\|\bar{Y}^1-\bar{Y}^2\right\|_{\mathcal{S}^{\infty}} &\leq \lambda \left(2T\left\|y^1-y^2\right\|_{\mathcal{S}^{\infty}}+T\sup_{t\in[0,T]}\left|k_t^1-k_t^2\right|+\sqrt{3T+6T^{1-\alpha}\tilde{A}^{2\alpha}}\left\|z^1-z^2\right\|_{BMO}\right),\\
   \left\|Z^1-Z^2\right\|_{BMO}&\leq 2\lambda \sqrt{1+12\lambda^2+24 \lambda^2\tilde{A}^2}\left(2T\left\|y^1-y^2\right\|_{\mathcal{S}^{\infty}}+T\sup_{t\in[0,T]}\left|k_t^1-k_t^2\right|\right. \\
   &\quad \left.+\sqrt{3T+6T^{1-\alpha}\tilde{A}^{2\alpha}}\left\|z^1-z^2\right\|_{BMO}\right).\\
\end{aligned}
$$
In view of Proposition \ref{continuous}, we obtain that
$$
\begin{aligned}
\sup_{t \in[0, T]}\left|K^1_{t}-K^2_{t}\right|
\leq \frac{12C}{c}\sup_{t \in[0, T]}\mathbf{E}\left[\left|\bar{Y}^1_t-\bar{Y}^2_t\right|\right]\leq \frac{12C}{c}\left\|\bar{Y}^1-\bar{Y}^2\right\|_{\mathcal{S}^{\infty}}.
\end{aligned}		
$$
%{\color{blue}What is $\bar{L}_T$? Therefore, we have How to obtain the following equation?} 
Therefore, we derive 
\begin{equation}\label{contraction}\begin{split}
   &\left\|Y^1-Y^2\right\|_{\mathcal{S}^{\infty}}+\left\|Z^1-Z^2\right\|_{BMO}+\sup_{t \in[0, T]}\left|K^1_{t}-K^2_{t}\right|\\
   %&\leq \left((1+\frac{3C}{c})+2\sqrt{1+12\lambda^2+24\lambda^2\tilde{A}^2}+\frac{3C}{c}\right)\left\|\bar{Y}^1-\bar{Y}^2\right\|_{\mathcal{S}^{\infty}}\\
   \leq &\lambda \hat{A} \left(2T\left\|y^1-y^2\right\|_{\mathcal{S}^{\infty}}+\sqrt{3T+6T^{1-\alpha}\tilde{A}^{2\alpha}}\left\|z^1-z^2\right\|_{BMO}+T\sup_{t\in[0,T]}\left|k^1_t-k^2_t\right|\right),\\
\end{split}   
\end{equation}
where $$\hat{A}=\left(1+\frac{36C}{c}+2\sqrt{1+12\lambda^2+24\lambda^2\tilde{A}^2}\right).$$

Now we define
\begin{equation}\label{delta}
  \hat{\delta}^{\tilde{A}}:=\min\left(\frac{1}{4\hat{A}\lambda}, \frac{1}{12\hat{A}\lambda}, \left(\frac{1}{24\tilde{A}^{2\alpha}}\hat{A}^2\lambda^2\right)^{\frac{1}{1-\alpha}},\delta^{\tilde{A}}\right). 
\end{equation}
It is straightforward to check that for any $T\in(0,\hat{\delta}^{\tilde{A}}]$, we have
$$\begin{aligned}
  &\left\|Y^1-Y^2\right\|_{\mathcal{S}^{\infty}}+\left\|Z^1-Z^2\right\|_{BMO}+\sup_{t \in[0, T]}\left|K^1_{t}-K^2_{t}\right|\\
   \leq&\frac{1}{2}\left(\left\|y^1-y^2\right\|_{\mathcal{S}^{\infty}}+\left\|z^1-z^2\right\|_{BMO}+\sup_{t\in[0,T]}\left|k^1_t-k^2_t\right|\right).\\  
\end{aligned}$$
The proof is complete.
\end{proof}

Now, we are ready to give the proof of Theorem $\ref{thmbddterminal}$.

\begin{proof}[Proof of Theorem \ref{thmbddterminal}] We take $\tilde{A}\geq \tilde{A}_0$ and choose $\hat{\delta}^{\tilde{A}}$ as $(\ref{delta})$. For $T\in(0,\hat{\delta}^{\tilde{A}}]$, we define $(Y^0,Z^0,K^0)=(0,0,0)$ and $(Y^{i+1},Z^{i+1},K^{i+1})=(Y^{Y^{i},Z^{i},K^{i}},Z^{Y^{i},Z^{i},K^{i}},K^{Y^{i},Z^{i},K^{i}})$ by recurrence. 
For convenience, we denote 
$$
\begin{aligned}
  f^n_t&:=f(t,Y^{n-1}_{t},\mathbf{P}_{Y^{n-1}_{t}},Z_t^{n},\mathbf{P}_{Z_t^{n-1}},G_{t}(K^{n-1})),\\
  f_t&:=f(t,Y_{t},\mathbf{P}_{Y_{t}},Z_t,\mathbf{P}_{Z_t},G_{t}(K)).
\end{aligned}
$$
By the proof of Lemma \ref{4.1.1}, for each $n\geq 1$, $K^n$ is the second component of the solution to $\mathbb{BSP}^{r^n}_{l^n}(s^n,a^n)$, where for each $t\in[0,T]$,
\begin{equation}\label{sarln}\begin{split}
& s_t^{n}=\mathbf{E}\left[\int_{0}^{t} f_s^n ds\right],\ l^n(t,x)=\E\left[L(t,\widetilde{Y}^n_t-\E[\widetilde{Y}^n_t]+x)\right],\\
& a^n=\mathbf{E}[\xi], \ r^n(t,x)=\E\left[R(t,\widetilde{Y}^n_t-\E[\widetilde{Y}^n_t]+x)\right]. \end{split}
 \end{equation}
and $\widetilde{Y}^n$ is the first component of the solution to  the following BSDE:
\begin{equation*}
\widetilde{Y}^n_t=\xi+\int_{t}^{T} f_s^nds-\int_{t}^{T}\widetilde{Z}^n_s dB_s.    
\end{equation*}
It follows from Lemma \ref{4.1.1} that there exists $(Y,Z,K)\in \mathscr{B}_{\tilde{A}}$ such that %{\color{blue}what is $(Y^n,Z^n,K^n)$?}
\begin{equation}\label{YnY}
\left\|Y^n-Y\right\|_{\mathcal{S}^{\infty}}\to 0,\ \left\|Z^n-Z\right\|_{BMO}\to 0,\ \sup_{t\in[0,T]}|K_t^n-K_t|\to 0.    
\end{equation}

In order to show that the triple $(Y,Z,K)$ is indeed a solution to MFBSDE with nonlinear resistance, it suffices to prove the Skorokhod conditions hold. The proof is similar to the one for Theorem 4.4 in \cite{LS}. For readers' convenience, we give a short proof here. For this purpose, let $(Y',Z',K')$ be the solution to the doubly mean reflected BSDE on $[0,T]$ with terminal value $\xi$, constant coefficient $\{f(s,Y_s,\P_{Y_s},Z_s,\P_{Z_s},G_s(K))\}_{s\in[0,T]}$ and loss functions $L,R$. It suffices to prove that $Y'\equiv Y$ and $K'\equiv K$. Let $(\widetilde{Y},\widetilde{Z})$ be the solution to the BSDE on $[0,T]$ with terminal value $\xi$ and constant coefficient $\{f(s,Y_s,\P_{Y_s},Z_s,\P_{Z_s},G_s(K))\}_{s\in[0,T]}$. According to Proposition \ref{pro-1}, we have
\begin{equation*}
Y'_{t}=\widetilde{Y}_t+K'_{T}-K'_{t}=\xi+\int_{t}^{T} f_s ds-\int_{t}^{T}\widetilde{Z}_s dB_s+K'_{T}-K'_{t},    
\end{equation*}
where $K'$ is the second component of the solution to $\mathbb{BSP}^{r}_{l}(s,a)$. 
Here, for any $(t,x)\in[0,T]\times \mathbb{R}$,
\begin{equation}\label{lrsaL}
\begin{split}
& s_t=\E\left[\int_{0}^t f_s ds\right], \ l(t,x)=\E\left[L(t,\widetilde{Y}_t-\E[\widetilde{Y}_t]+x)\right],\\
& a=\E[\xi], \ r(t,x)=\E\left[R(t,\widetilde{Y}_t-\E[\widetilde{Y}_t]+x)\right].%,\ \forall(t,x)\in[t_0,t_0+h]\times \mathbb{R}.
\end{split}
\end{equation}
Therefore, we have
\begin{align*}
   \sup_{t\in[0,T]}|s_{t}^{n}-s_{t}|\rightarrow 0, \ n\rightarrow \infty.
\end{align*}
In addition, applying Assumptions \ref{assLR}, \ref{assqua} and B-D-G inequality , we obtain that 
$$
\begin{aligned}
	\bar{L}^n_{T}&:= \sup_{(t,x)\in[0,T]\times \mathbb{R}}\left|l^n(t,x)-l(t,x)\right|\leq 2C\sup_{t\in[0,T]}\mathbf{E}\left[\left|\widetilde{Y}^n_t-\widetilde{Y}_t\right|\right]\\
    %&=\sup_{(t,x)\in[0,T]\times \mathbb{R}} \left|\mathbf{E}\left[L(t,\widetilde{Y}^n_t-\E[\widetilde{Y}^n_t]+x)\right]-\E\left[L(t,\widetilde{Y}_t-\E[\widetilde{Y}_t]+x)\right]\right|\\
	%&\leq \sup_{t\in[0,h]}\mathbf{E}\left[C\left|y^{(m)}_{t_0+h-t}-\widetilde{Y}_{t_0+h-t}-\mathbf{E}\left[y^{(m)}_{t_0+h-t}-\widetilde{Y}_{t_0+h-t}\right]\right|\right]\\
    &\leq 2C\sup_{t\in[0,T]}\mathbf{E}\left[\int_{t}^{T}\left|f_s^n-f_s\right|ds\right]+2C\sup_{t\in[0,T]}\mathbf{E}\left[\left|\int_{t}^{T}\tilde{Z}^n_s-\tilde{Z}_sdB_s\right|\right]\\
    &\leq 2\lambda \mathbf{E}\left[\int_{0}^{T}\left|Y^n_s-Y_s\right|ds\right]+\lambda \int_{0}^{T}\left|K^{n-1}_s-K_s\right|ds \\
    &\quad + \lambda \mathbf{E}\left[\int_{0}^{T}\left(1+\left|Z^n_s\right|+\left|Z_s\right|\right)\left|Z^n_s-Z_s\right|ds\right]\\
&\quad+\lambda \mathbf{E}\left[\int_{0}^{T}\left(1+\mathbf{E}\left[\left|Z^n_s\right|\right]+\mathbf{E}\left[\left|Z_s\right|\right]\right)\left|Z^n_s-Z_s\right|ds\right]\\
&\quad+\Lambda \mathbf{E}\left[\int_{0}^{T}\left(\tilde{Z}^n_s-\tilde{Z}_s\right)^2ds\right]^{\frac{1}{2}},\\ 
\end{aligned}
$$
where $\Lambda$ is a constant. Combining with \eqref{YnY}, we have $\lim_{n\rightarrow \infty}\bar{L}^n_{T}=0.$ 

Similarly, 
$$
\begin{aligned}
	\bar{R}^n_{T}&:= \sup_{(t,x)\in[0,T]\times \mathbb{R}}\left|r^n(t,x)-r(t,x)\right|\\
	&\leq 2C\sup_{t\in[0,T]}\mathbf{E}\left[\left|\widetilde{Y}^n_t-\widetilde{Y}_t\right|\right]\rightarrow 0, \ \ n \rightarrow \infty.
\end{aligned}
$$
 In view of Proposition \ref{continuous}, we have 
$$
\begin{aligned}
 \sup_{t\in[0,T]}\left|K_t^{n}-K'_t\right|
 &\leq \frac{4C}{c}\sup_{t\in[0,T]}\left|s_t^{n}-s_t\right|+\frac{2}{c}\left(\bar{L}^n_{T}\vee\bar{R}^n_{T}\right)\to 0,\ \text{as} \ \  n\to \infty , 
\end{aligned}
$$
which, together with \eqref{YnY}, implies that $ K'\equiv K$. By the uniqueness of BSDE, we obtain that $Y'\equiv Y$. Therefore, the Skorokhod condition holds. The proof is complete.    
\end{proof}

Finally, we consider the general time horizon $[0,T]$ and give the proof of Theorem \ref{globalMFBSDE}.

\begin{proof}[Proof of Theorem \ref{globalMFBSDE}] 
% The proof is similar to the one for Theorem 2.13 in \cite{LP}. For readers' convenience, we will briefly restate the proof here.
  {\bf Step 1. } We claim that for given $k\in BV[0,T]$, the following MFBSDE with double mean reflections
 \begin{equation}\label{MFBSDEDMRk}
\begin{cases}
Y_t=\xi+\int_t^T f(s,Y_s,\mathbf{P}_{Y_s},Z_s,\mathbf{P}_{Z_s},G_s(k))ds-\int_t^T Z_s dB_s+K_T-K_t, \\
\mathbf{E}[L(t,Y_t)]\leq 0\leq \mathbf{E}[R(t,Y_t)], \\
K_t=K^R_t-K^L_t,\ K^R,K^L\in I[0,T],\\
\int_0^T \mathbf{E}[R(t,Y_t)]dK_t^R=\int_0^T \mathbf{E}[L(t,Y_t)]dK^L_t=0.\\
\end{cases}
\end{equation}
admits a unique solution $(Y,Z,K)\in\mathcal{S}^{\infty}\times BMO \times BV[0,T]$. The proof is similar with the one for Theorem 2.13 in \cite{LP}. For the reader's convenience, we restate the complete proof here.

Indeed, according to Theorem $\ref{thmbddterminal}$, we take an appropriate constant $\bar{h}>0$ depending on $\tilde{A}_0,\hat{H},C,c,M,\lambda$, where $$\tilde{A}_0=(3+\frac{24C}{c})\bar{H}+\frac{2M}{c}+(\frac{2\bar{H}}{\lambda}+\frac{1}{3})e^{9\lambda \bar{H}},$$
$\bar{H}$ is the constant defined in \eqref{barH} and $\hat{H}=\tilde{H}+\lambda\sup_{t\in[0,T]}\left|k_t\right|$. 

Then the doubly mean reflected MFBSDE \eqref{MFBSDEDMRk} admits a unique solution $(Y^1,Z^1,K^1)\in\mathcal{S}^{\infty}[T-\bar{h},T]\times BMO[T-\bar{h},T]\times BV[T-\bar{h},T]$. Moreover, according to Lemma \ref{lem4.6}, we have $\left\|Y^1\right\|_{\mathcal{S}^{\infty}[T-\bar{h},T]}\leq \bar{H}.$ 

Next,  taking $T-\bar{h}$ as the terminal time and applying Theorem $\ref{thmbddterminal}$, the doubly mean reflected MFBSDE \eqref{MFBSDEDMRk} admits a unique solution $(Y^2,Z^2,K^2)\in\mathcal{S}^{\infty}[T-2\bar{h},T-\bar{h}]\times BMO[T-2\bar{h},T-\bar{h}]\times BV[T-2\bar{h},T-\bar{h}]$ to Eq. \eqref{MFBSDEDMRk}. 

Repeating this procedure, for $i=1,...,[\frac{T}{\bar{h}}]$, we set 
\begin{equation}\begin{split}
 &Y_t=\sum_{i=1}^{[\frac{T}{\bar{h}}]} Y_t^i \mathbf{1}_{[T-i\bar{h},T-(i-1)\bar{h})}+Y_T^1 \mathbf{1}_{\left\{T\right\}}, \ Z_t=\sum_{i=1}^{[\frac{T}{\bar{h}}]} Z_t^i \mathbf{1}_{[T-i\bar{h},T-(i-1)\bar{h})}+Z_T^1 \mathbf{1}_{\left\{T\right\}},\\
\end{split}  
\end{equation}
and $K_t=\sum_{i=1}^{[\frac{T}{\bar{h}}]}(\sum_{m=1}^{i-1}K^{m}_{T-(m-1)}\bar{h}+K_t^i)\mathbf{1}_{[T-i\bar{h},T-(i-1)\bar{h})}+K_T^1 \mathbf{1}_{\left\{T\right\}}$. It is easy to check that  $(Y,Z,K)\in\mathcal{S}^{\infty}\times BMO \times BV[0,T]$ is a solution to \eqref{MFBSDEDMRk}. 

According to the uniqueness of local solution on each time interval, the uniqueness of the global solution on $[0,T]$ is obtained.

{\bf Step 2.} For $k^i\in BV[0,T],\ i=1,2$, let $(Y^i,Z^i,K^i)$ be the unique solution to \eqref{MFBSDEDMRk} with data $k^i$. According to Lemma \ref{4.1.1} , there exists a constant $\tilde{h}>0$ depending only on $H,\tilde{H},\lambda, C,c,M$ and $\sup_{t\in[0,T]}\left|k_t^1\right|, \sup_{t\in[0,T]}\left|k_t^2\right|$, such that for any $1 \leq i \leq N$, where $N$ is the smallest integer satisfying $N\geq \frac{T}{\tilde{h}}$,
$$
\begin{aligned}
  &\left\|Y^{1}-Y^{2}\right\|_{\mathcal{S}^{\infty}[(i-1)\tilde{h},i\tilde{h} \wedge T]}+\left\|Z^{1}-Z^{2}\right\|_{BMO[(i-1)\tilde{h},i\tilde{h} \wedge T]}+\sup_{t\in[(i-1)\tilde{h},i\tilde{h} \wedge T]}\left|K_{t}^{1}-K_{t}^{2}\right|\\
  &\leq \frac{1}{2} \left(\left\|Y^{1}-Y^{2}\right\|_{\mathcal{S}^{\infty}[(i-1)\tilde{h},i\tilde{h} \wedge T]}+\left\|Z^{1}-Z^{2}\right\|_{BMO[(i-1)\tilde{h},i\tilde{h} \wedge T]}+\sup_{t\in[(i-1)\tilde{h},i\tilde{h} \wedge T]}\left|k_t^1-k_t^2\right|\right).
\end{aligned}
$$
%$$\sup_{t\in[(i-1)\tilde{h},i\tilde{h} \wedge T]}\left|K_t^1-K_t^2\right|\leq \frac{1}{2}\sup_{t\in[0,T]}\left|k_t^1-k_t^2\right|.$$
Therefore, we have 
$$\sup_{t\in[0,T]}\left|K_t^1-K_t^2\right|\leq \max_{1\leq i \leq N}\left(\sup_{t\in[(i-1)\tilde{h},i\tilde{h} \wedge T]}\left|K_t^1-K_t^2\right|\right)\leq \frac{1}{2}\sup_{t\in[0,T]}\left|k_t^1-k_t^2\right|.$$
By the principle of standard contraction mapping, the doubly mean reflected MFBSDE with resistance \eqref{nonlinearyz} admits a unique solution $(Y,Z,K)\in\mathcal{S}^{\infty}\times BMO \times BV[0,T]$. 
\end{proof}
%{\color{red} Some proof of Section 3 and Section 4 are }

%\begin{remark}
    %Compared with the bounded terminal value case in \cite{LS}, we have obtained the well-posedness of doubly mean reflected MFBSDE with Z-distribution in the {\bf Step 1} of Theorem $\ref{globalMFBSDE}$. 
%\end{remark}


   
\section{Doubly mean reflected MFBSDEs with density function}
In this section, we consider the well-posedness of doubly mean reflected MFBSDE with resistance given in terms of density function. Before investigating this problem, we first consider the regularity of the last component of the solution  to MFBSDEs with double mean reflections. For this purpose, we introduce the following assumption for the loss functions.

\begin{assumption}\label{assLRbounded}
    The loss functions $L,R \in C^{1,2}_b([0,T]\times\mathbb{R})$, where $C^{1,2}_b([0,T]\times\mathbb{R})$ is the collection  of functions that are continuously differentiable in their first variable and twice continuously differentiable in their second variable, and both
derivatives are uniformly bounded.
\end{assumption}
\begin{proposition}\label{Klip}
    Suppose that Assumptions \ref{asslip2}, \ref{assLRbounded} hold.  Let $(Y,Z,K)$ denote the unique deterministic solution to MFBSDE \eqref{MFBSDEDMR} with double mean reflections. Moreover, suppose that 
    \begin{align}\label{bdd of Z}
        \sup_{t\in[0,T]}\E[|Z_t|^2]<\infty.
    \end{align}
Then the process $K$ is Lipschitz continuous. Furthermore, $K$ has the density function.
\end{proposition}
\begin{proof}
    According to Proposition 4.2 in \cite{L}, solution $K$ is the second component of solution of backward Skorokhod problem $\mathbb{BSP}_{l}^{r}(s,a)$. Here, for each $t\in[0,T]$,
\begin{equation}\begin{split}
&s_t=\mathbf{E}\left[\int_{0}^{t}f(s,Y_s,\mathbf{E}[Y_s],Z_s,\mathbf{E}[Z_s])ds\right],\ a=\mathbf{E}[\xi],\\
&l(t,x):=\mathbf{E}\left[L(t,y_t-\E[y_t]+x)\right], \ r(t,x):=\mathbf{E}\left[R(t,y_t-\E[y_t]+x)\right], 
\end{split}
 \end{equation}
where $(y,z)$ is the solution to  BSDE with terminal $\xi$ and $\left\{f(t,Y_t,\mathbf{E}[Y_t],Z_t,\mathbf{E}[Z_t])\right\}_{t\in[0,T]}$. 
For any $t\in[0,T]$, $K_t=\bar{K}_T-\bar{K}_{T-t}$, where $\bar{K}$ is the second component of the solution to the Skorokhod problem $\mathbb{SP}_{\bar{l}}^{\bar{r}}(\bar{s})$. Here, for any $t\in[0,T]$, 
\begin{equation}
\bar{s}_{t}=a+s_{T}-s_{T-t}=\mathbf{E}[y_{T-t}],\ \bar{l}(t,x)=l (T-t,x), \ \bar{r}(t,x)=r(T-t,x).
\end{equation}
Let {$\bar{\phi}_{t}, \bar{\psi}_{t}$}  be the unique solutions to the following equations, respectively
$$\bar{l}(t,\bar{s}_t+x)=0, \ \bar{r}(t,\bar{s}_t+x)=0.$$
Applying Remark 2.7 in \cite{Li}, we have 
\begin{equation}
\begin{split}
   \left|K_t-K_s\right|&\leq \sup_{u_1,u_2\in[s,t]}\left|K_{u_1}-K_{u_2}\right|\\
   &=\sup_{u_1,u_2\in[s,t]}\left|\bar{K}_T-\bar{K}_{T-u_1}-(\bar{K}_T-\bar{K}_{T-u_2})\right|\\
   &\leq \sup_{u_1,u_2\in[s,t]}\left|\bar{K}_{T-u_1}-\bar{K}_{T-u_2}\right|=\sup_{t_1,t_2\in[T-t,T-s]}\left|\bar{K}_{t_1}-\bar{K}_{t_2}\right|\\
   &\leq \sup_{p_1,p_2\in[T-t,T-s]}\left|\bar{\phi}_{p_1}-\bar{\phi}_{p_2}\right|+\sup_{p_1,p_2\in[T-t,T-s]}\left|\bar{\psi}_{p_1}-\bar{\psi}_{p_2}\right|\\
   &{= \sup_{p_1,p_2\in[s,t]}\left|{\phi}_{p_1}-{\phi}_{p_2}\right|+\sup_{p_1,p_2\in[s,t]}\left|{\psi}_{p_1}-{\psi}_{p_2}\right|},
\end{split}    
\end{equation}
where $\bar{\phi}_t=\phi_{T-t}$ and ${\psi}_t=\bar{\psi}_{T-t}$. It remains to prove that $\phi_t,\psi_t$ are Lipschitz continuous. We only prove this fact for $\phi$. First, we have
\begin{equation}
\begin{split}
    \left|\phi_{t}-\phi_{s}\right|&=\left|\bar{\phi}_{T-t}-\bar{\phi}_{T-s}\right|\\
    &\leq  \frac{1}{c}\left|\bar{l}(T-t,\bar{s}_{T-t}+\bar{\phi}_{T-t})-\bar{l}(T-t,\bar{s}_{T-t}+\bar{\phi}_{T-s})\right|\\
    &=\frac{1}{c}\left|\bar{l}(T-t,\bar{s}_{T-t}+\bar{\phi}_{T-s})\right|\\
    % &=\frac{1}{c}\left|l(p_1,\bar{s}_{T-p_1}+\bar{\phi}_{T-p_2})\right|\\
    &=\frac{1}{c}\left|\mathbf{E}\left[L(t,y_{t}+\phi_{s})\right]\right|.
\end{split}    
\end{equation}
Applying It\^o's  formula to $L(\cdot,y_{\cdot}+\phi_{s})$ on the interval $[s,t]$ and then taking expectations on both sides of the equation, combining with Assumptions \ref{asslip2},\ \ref{assLRbounded} and \eqref{bdd of Z}, we have
\begin{equation}
\begin{split}
  \left|\mathbf{E}\left[L(t,y_{t}+\phi_{s})\right]\right|&\leq  \mathbf{E}\left[\int_{s}^{t}\left|\frac{\partial L}{\partial r}(r,y_{r}+\phi_{s})\right|dr\right]+\frac{1}{2}\mathbf{E}\left[\int_{s}^{t}\left|\frac{\partial^2 L}{\partial x^2}(r,y_{r}+\phi_{s})Z_r^2\right|dr\right]\\ 
  &\quad+\mathbf{E}\left[\int_{s}^{t}\left|\frac{\partial L}{\partial x}(r,y_{r}+\phi_{s})\left(|Y_r|+\mathbf{E}[|Y_r|]+|Z_r|+\mathbf{E}[|Z_r|]\right)\right|dr\right]\\
  &\quad+\mathbf{E}\left[\int_{s}^{t}\left|\frac{\partial L}{\partial x}(r,y_{r}+\phi_{s})f(r,0,0,0,0)\right|dr\right]\\
  &\leq \Lambda |t-s|.
\end{split}
\end{equation}
where $\Lambda$ is an appropriate constant.
%Therefore, we have 
%\begin{equation*}
%\left|\phi_{p_1}-\phi_{p_2}\right| \leq M \left|p_1-p_2\right| 
%\end{equation*}}
%Similarly, we have 
%\begin{equation*}
%    \left|\psi_{p_1}-\psi_{p_2}\right|\leq M \left|p_1-p_2\right|.
%\end{equation*}
The proof is complete. %Thus, we can obtain the Lipschitz continuous for $K$.
\end{proof}
\begin{lemma}\label{estimateDYDZ}
 Consider the following MFBSDE:
\begin{equation}\label{MFBSDEm}
  Y_t=\xi+\int_t^T f(s,Y_s,\mathbf{E}[Y_s],Z_s,\mathbf{E}[Z_s])ds-\int_t^T Z_s dB_s.  
\end{equation}
Let Assumption \ref{asslip2} hold. Suppose that $f$ is continuously differentiable in $y$ and $z$ with uniformly bounded derivative. Moreover, assume that $\xi$ and $f(\cdot, y,y',z,z')$ are Malliavin differentiable for each $y,y',z,z'$ with 
    \begin{itemize}
        \item [(i)] $\sup_{\theta \in[0,T]}\mathbf{E}\left[|D_{\theta}\xi|^p\right]<\infty$;
        \item [(ii)] $\sup_{\theta \in[0,T]}\mathbf{E}\left[\left(\int_{0}^{T}|D_{\theta}f(t,Y_t,\mathbf{E}[Y_t],Z_t,\mathbf{E}[Z_t])|dt\right)^{p}\right]<\infty$.
    \end{itemize}  
    Then, we have $\mathbf{E}\left[\sup_{t\in[0,T]}\left|Z_t\right|^p\right]<\infty$. %Therefore,  $\left\{D_t Y_t\right\}_{t\in[0,T]}$ is a version of $\left\{Z_t\right\}_{t\in[0,T]}$ and    
\end{lemma}

\begin{proof}
      First, similar to Proposition 5.3 in \cite{EPQ}, it is easily check that the Malliavin derivatives $\left\{(D_{\theta}Y_t,D_{\theta}Z_t)\right\}_{0\leq \theta,t\leq T}$ satisfies the following linear BSDE:  
\begin{equation}\label{eqDYDZ}
\begin{split}
D_\theta Y_t &= D_\theta \xi + \int_t^T \Big[ \partial_y f(r,Y_r,\mathbf{E}[Y_r],Z_r,\mathbf{E}[Z_r])D_\theta Y_r+\partial_z f(r,Y_r,\mathbf{E}[Y_r],Z_r,\mathbf{E}[Z_r]) D_\theta Z_r \\
&\quad+ D_\theta f(r,Y_r,\mathbf{E}[Y_r],Z_r,\mathbf{E}[Z_r]) \Big] dr-\int_t^T D_\theta Z_r dB_r, \qquad 0 \leq \theta \leq t \leq T;
\end{split}
\end{equation}
\begin{equation}\label{eqDYDZ0}
D_\theta Y_t = 0, \qquad D_\theta Z_t = 0, \qquad 0 \leq t < \theta \leq T.
\end{equation}
According to Lemma \ref{proMFBSDE}, we have 
\begin{equation}\label{eqDYDZ1}
\begin{split}
&\mathbf{E}\left[\sup_{t\in[0,T]}\left|D_{\theta}Y_t\right|^p\right]+\mathbf{E}\left[\left(\int_{0}^{T}\left|D_{\theta}Z_t\right|^2dt\right)^{\frac{p}{2}}\right]\\
&\leq C_p\left(\mathbf{E}\left[\left|D_{\theta}\xi\right|^p\right]+\mathbf{E}\left[\left(\int_{0}^{T}|D_{\theta}f(r,Y_r,\mathbf{E}[Y_r],Z_r,\mathbf{E}[Z_r])|dt\right)^{p}\right]\right)<\infty, \  \forall \theta\in[0,T].\\
\end{split}
\end{equation}
For $t\leq s$, we have
$$Y_s=Y_t-\int_{t}^{s}f(r,Y_r,\mathbf{E}[Y_r],Z_r,\mathbf{E}[Z_r])dr+\int_{t}^{s}Z_rdB_r.$$
Therefore, for $t< \theta \leq s$,
\begin{equation}
\begin{split}
D_\theta Y_s &= Z_{\theta}- \int_{\theta}^{s} \Big[ \partial_y f(r,Y_r,\mathbf{E}[Y_r],Z_r,\mathbf{E}[Z_r])D_\theta Y_r+\partial_z f(r,Y_r,\mathbf{E}[Y_r],Z_r,\mathbf{E}[Z_r]) D_\theta Z_r \\
&\quad+ D_\theta f(r,Y_r,\mathbf{E}[Y_r],Z_r,\mathbf{E}[Z_r]) \Big] dr+\int_{\theta}^{s} D_\theta Z_r dB_r;
\end{split}
\end{equation}
Then, taking $\theta=s$, we have $D_sY_s=Z_s, \ a.s.$  Therefore,  $\mathbf{E}\left[\sup_{t\in[0,T]}\left|Z_t\right|^p\right]<\infty$.
\end{proof}
 

\begin{remark}
 For example, let us consider MFBSDE with the terminal condition $\xi=g(X_T)$ and the   generator $f(t,y,y',z,z')=F(X_t,y,y',z,z')$, where $X_t$ satisfies the following SDE:
    \begin{equation*}
        X_t=x_0+\int_{0}^{t}b(X_s)ds+\int_{0}^{t}\sigma(X_s)dB_s,
    \end{equation*}
 and the functions $b,\sigma,F,g$ are deterministic and continuously differentiable such that $\partial_xb$, $\partial_{x}\sigma$,  $\partial_{x}g$, $\partial_{x}F$,  $\partial_{y}F$, $\partial_{z}F$ are bounded. Then, the assumptions in Lemma \ref{estimateDYDZ} are satisfied. Therefore, we can derive 
  $\mathbf{E}\left[\sup_{t\in[0,T]}\left|Z_t\right|^p\right]<\infty.$ 
  
In particular, adding a deterministic function $A$ into Eq. \eqref{MFBSDEm}, i.e., $(Y,Z)$ satisfies the following equation
\begin{align*}
    Y_t=\xi+\int_t^T f(s,Y_s,\mathbf{E}[Y_s],Z_s,\mathbf{E}[Z_s])ds-\int_t^T Z_s dB_s+A_T-A_t,
\end{align*}
the conclusion of Lemma \ref{estimateDYDZ} remains valid. Therefore, consider the MFBSDE with double mean reflections:
\begin{equation}\label{EMFBSDE}
\begin{cases}
Y_t=\xi+\int_t^T f(s,Y_s,\mathbf{E}[Y_s],Z_s,\mathbf{E}[Z_s])ds-\int_t^T Z_s dB_s+K_T-K_t, \\
\mathbf{E}[L(t,Y_t)]\leq 0\leq \mathbf{E}[R(t,Y_t)],\\ 
{K_t=K^+_t-K^-_t, \ t\in[0,T]}\\
\int_0^T \mathbf{E}[R(t,Y_t)]dK^+_t=\int_0^T \mathbf{E}[L(t,Y_t)]dK^-_t=0. 
\end{cases}
\end{equation}
Under the same assumptions for $\xi$ and $f$ as in Lemma \ref{estimateDYDZ}, we still have $\mathbf{E}\left[\sup_{t\in[0,T]}\left|Z_t\right|^p\right]<\infty$.. %In addition, if $L\equiv -\infty$, Eq.\eqref{EMFBSDE} degenerates into the general mean reflected MFBSDE, whose solution still satisfies $\mathbf{E}\left[\sup_{t\in[0,T]}\left|Z_t\right|^p\right]<\infty$.  
 \end{remark}
   
Suppose that the compensating term $K$ of the doubly mean reflected MFBSDE with nonlinear resistance is absolutely continuous w.r.t. Lebesgue measure, i.e., there exists a density process $k$ such that $K_t=\int_0^t k_sds$. In this section, we consider doubly mean reflected MFBSDE \eqref{nonlinearyz} when the nonlinear resistance is given as $G_s(K)=k_s$ and we focus on the case $p=2$. In this case, doubly mean reflected MFBSDE \eqref{nonlinearyz} can be rewritten as follows:
\begin{equation}\label{mfbsdefk}
\begin{cases}
Y_t=\xi+\int_t^T \left(f(s,Y_s,\mathbf{E}[Y_s],Z_s,\mathbf{E}[Z_s],k_s)+k_s\right)ds-\int_t^T Z_s dB_s, \\
\mathbf{E}[L(t,Y_t)]\leq 0\leq \mathbf{E}[R(t,Y_t)],\\ 
{k_t=k^+_t-k^-_t, \ t\in[0,T] \textrm{ and }k^+,k^-\geq 0,}\\
\int_0^T \mathbf{E}[R(t,Y_t)]k_t^+dt=\int_0^T \mathbf{E}[L(t,Y_t)]k^-_tdt=0.
\end{cases}
\end{equation}
Unfortunately, the nonlinear resistance function $G$ does not satisfy Assumption \ref{assgk}. To obtain well-posedness of Eq.\eqref{mfbsdefk}, we will construct the solution via the following auxiliary doubly mean reflected MFBSDE without resistance
\begin{equation}\label{aumdrmfbsdewr}
\begin{cases}
\widetilde{Y}_t=\xi+\int_t^T f(s,\widetilde{Y}_s,\mathbf{E}[\widetilde{Y}_s],\widetilde{Z}_s,\mathbf{E}[\widetilde{Z}_s],0)ds-\int_t^T\widetilde{Z}_s dB_s+\widetilde{K}_T-\widetilde{K}_t, \\
\mathbf{E}[L(t,\widetilde{Y}_t)]\leq 0\leq \mathbf{E}[R(t,\widetilde{Y}_t)],\\ 
\widetilde{K}_t=\widetilde{K}^{+}_t-\widetilde{K}^{-}_t,\ \widetilde{K}^+,\widetilde{K}^{-} \in I[0,T],\\
\int_0^T \mathbf{E}[R(t,\widetilde{Y}_t)]d\widetilde{K}_t^+=\int_0^T \mathbf{E}[L(t,\widetilde{Y}_t)]d\widetilde{K}^-_t=0.
\end{cases}
\end{equation}

According to Theorem 3.5 in \cite{LS}, the doubly mean reflected MFBSDE \eqref{aumdrmfbsdewr} admits  a unique solution $(\widetilde{Y},\widetilde{Z},\widetilde{K})\in \mathcal{S}^2 \times \mathcal{H}^2 \times BV[0,T]$. Suppose additionally that the loss functions $L,R\in C^{1,2}_b([0,T]\times\mathbb{R})$ and 
\begin{align}\label{bdd of tilde Z}
    \sup_{t\in[0,T]}\E[|\widetilde{Z}_t|^2]<\infty.
\end{align}
According to Proposition \ref{Klip},  there exist three measurable functions  $\widetilde{k},\widetilde{k}^+,\widetilde{k}^-$ with $\widetilde{k}^+,\widetilde{k}^-\geq0$, such that
\begin{align*}
\widetilde{K}_t=\int_{0}^{t}\widetilde{k}_s ds, \ \widetilde{K}^+_t=\int_{0}^{t}\widetilde{k}^+_s ds, \ \widetilde{K}^-_t=\int_{0}^{t}\widetilde{k}^-_s ds.
\end{align*}
Thus, we rewrite Eq. \eqref{aumdrmfbsdewr} as follows:
\begin{equation}
\begin{cases}
\widetilde{Y}_t=\xi+\int_t^T \left(f(s,\widetilde{Y}_s,\mathbf{E}[\widetilde{Y}_s],\widetilde{Z}_s,\mathbf{E}[\widetilde{Z}_s],0)+\widetilde{k}_s\right)ds-\int_t^T\widetilde{Z}_s dB_s, \\
\mathbf{E}[L(t,\widetilde{Y}_t)]\leq 0\leq \mathbf{E}[R(t,\widetilde{Y}_t)],\\ 
{\widetilde{k}_t=\widetilde{k}^+_t-\widetilde{k}^-_t, \ \widetilde{k}^+,\widetilde{k}^-\geq 0,} \\
\int_0^T \mathbf{E}[R(t,\widetilde{Y}_t)]\widetilde{k}^+_tdt=\int_0^T \mathbf{E}[L(t,\widetilde{Y}_t)]\widetilde{k}^-_tdt=0.
\end{cases}
\end{equation}

\begin{theorem}
Let Assumptions \ref{asslip2}, \ref{assLRbounded} and \eqref{bdd of tilde Z} hold. Moreover, suppose that for any $s\in[0,T]$, each of the following two equations admits a unique deterministic nonnegative solution
\begin{equation}\label{k+k-}
    \begin{split}
        &f(s,\widetilde{Y}_s,\mathbf{E}[\widetilde{Y}_s],\widetilde{Z}_s,\mathbf{E}[\widetilde{Z}_s],k)-f(s,\widetilde{Y}_s,\mathbf{E}[\widetilde{Y}_s],\widetilde{Z}_s,\mathbf{E}[\widetilde{Z}_s],0)+k=\widetilde{k}_s^+,\\
        &f(s,\widetilde{Y}_s,\mathbf{E}[\widetilde{Y}_s],\widetilde{Z}_s,\mathbf{E}[\widetilde{Z}_s],-k)-f(s,\widetilde{Y}_s,\mathbf{E}[\widetilde{Y}_s],\widetilde{Z}_s,\mathbf{E}[\widetilde{Z}_s],0)-k=-\widetilde{k}_s^-.
    \end{split}
\end{equation}
Then the mean doubly reflected MFBSDE \eqref{mfbsdefk} with nonlinear resistance has at least one solution. 
\end{theorem}
\begin{proof}
%Let $Y:=\widetilde{Y}, \ Z:=\widetilde{Z}$, we just need to find solutions $k^+_t$ and $k^-_t$  to the following Skorokhod conditions:
%$$\int_0^T \mathbf{E}[R(t,\widetilde{Y}_t)]k^+_tdt=\int_0^T \mathbf{E}[L(t,\widetilde{Y}_t)]k^-_tdt=0.$$
Let $R_t=\{\E[R(t,Y_t)]=0\}$, $L_t=\{\E[L(t,Y_t)]=0\}$ and $A_t=\{\E[L(t,Y_t)]<0<\E[R(t,Y_t)]\}$. It is easy to check that for any $t$, $\mathbf{1}_{R_t}+\mathbf{1}_{L_t}+\mathbf{1}_{A_t}=1$ and 
\begin{align*}
    k_t=\begin{cases}
        k^+_t, & t\in R_t,\\
        0, &t\in A_t,\\
        -k^-_t, &t\in L_t.
    \end{cases}
\end{align*}
Therefore, we may rewrite the first equation to mean doubly reflected MFBSDE \eqref{mfbsdefk} with nonlinear resistance as follows:
    \begin{equation}
        \begin{split}
            Y_t&=\xi+\int_{t}^{T}\left(f(s,Y_s,\mathbf{E}[Y_s],Z_s,\mathbf{E}[Z_s],k_s)+k_s\right) ds-\int_{t}^{T}Z_sdB_s\\
           % &=\xi+\int_{t}^{T}\left(f(s,Y_s,\mathbf{E}[Y_s],Z_s,\mathbf{E}[Z_s],0)+k_s\right) ds-\int_{t}^{T}Z_sdB_s\\
            %&\quad +\int_{t}^{T}\mathbf{1}_{\left\{\mathbf{E}[L(t,Y_t)]=0\right\}}(b(s,\mathbf{E}[\widetilde{Y}_s],\mathbf{E}[\widetilde{Z}_s],k_s)-b(s,\mathbf{E}[\widetilde{Y}_s],\mathbf{E}[\widetilde{Z}_s],0)+k_s)ds\\
            %&\quad +\int_{t}^{T}\mathbf{1}_{\left\{\mathbf{E}[R(t,Y_t)]=0\right\}}(b(s,\mathbf{E}[\widetilde{Y}_s],\mathbf{E}[\widetilde{Z}_s],k_s)-b(s,\mathbf{E}[\widetilde{Y}_s],\mathbf{E}[\widetilde{Z}_s],0)+k_s)ds\\
           &=\xi+\int_{t}^{T}f(s,Y_s,\mathbf{E}[Y_s],Z_s,\mathbf{E}[Z_s],0) ds-\int_{t}^{T}Z_sdB_s\\
           &\quad +\int_{t}^{T}\mathbf{1}_{R_t}\left(f(s,Y_s,\mathbf{E}[Y_s],Z_s,\mathbf{E}[Z_s],k^+_s)-f(s,Y_s,\mathbf{E}[Y_s],Z_s,\mathbf{E}[Z_s],0)+k_s^+\right)ds\\
            &\quad +\int_{t}^{T}\mathbf{1}_{L_t}\left(f(s,Y_s,\mathbf{E}[Y_s],Z_s,\mathbf{E}[Z_s],{-k^-_s})-f(s,Y_s,\mathbf{E}[Y_s],Z_s,\mathbf{E}[Z_s],0)-k_s^-\right)ds.
        \end{split}
    \end{equation} 
For any $s\in[0,T]$, let $Y_s=\widetilde{Y}_s, Z_s=\widetilde{Z}_s$ and $k_s=k^+_s-k^-_s$, where $k^+_s,k^-_s$ are the solutions to equations in \eqref{k+k-}, respectively. We claim that $(Y,Z,k)$ is a solution to \eqref{mfbsdefk}. It remains to prove that $k^+,k^-$ satisfy the flat-off condition. In fact, for any $s\in [0,T]$, it is easy to check that $k^+_s>0$ (resp., $k^-_s>0$) if and only if $\widetilde{k}^+_s>0$ (resp., $\widetilde{k}^-_s>0$). Therefore, the flat-off conditions hold. 
    \iffalse
     \begin{equation}
        \begin{split}
            Y_t&=\xi+\int_{t}^{T}\left(f(s,Y_s,\mathbf{E}[Y_s],Z_s,\mathbf{E}[Z_s],k_s)+k_s\right) ds-\int_{t}^{T}Z_sdB_s\\
            &=\xi+\int_{t}^{T}\left(f(s,Y_s,\mathbf{E}[Y_s],Z_s,\mathbf{E}[Z_s],0)+k_s\right) ds-\int_{t}^{T}Z_sdB_s\\
            &\quad +\int_{t}^{T}\mathbf{1}_{\left\{\mathbf{E}[L(t,Y_t)]=0\right\}}(b(s,\mathbf{E}[\widetilde{Y}_s],\mathbf{E}[\widetilde{Z}_s],k_s)-b(s,\mathbf{E}[\widetilde{Y}_s],\mathbf{E}[\widetilde{Z}_s],0)+k_s)ds\\
            &\quad +\int_{t}^{T}\mathbf{1}_{\left\{\mathbf{E}[R(t,Y_t)]=0\right\}}(b(s,\mathbf{E}[\widetilde{Y}_s],\mathbf{E}[\widetilde{Z}_s],k_s)-b(s,\mathbf{E}[\widetilde{Y}_s],\mathbf{E}[\widetilde{Z}_s],0)+k_s)ds\\
           &=\xi+\int_{t}^{T}\left(f(s,Y_s,\mathbf{E}[Y_s],Z_s,\mathbf{E}[Z_s],0)+k_s\right) ds-\int_{t}^{T}Z_sdB_s\\
           &\quad +\int_{t}^{T}\mathbf{1}_{\left\{\mathbf{E}[R(t,Y_t)]=0\right\}}(b(s,\mathbf{E}[\widetilde{Y}_s],\mathbf{E}[\widetilde{Z}_s],k_s^+)-b(s,\mathbf{E}[\widetilde{Y}_s],\mathbf{E}[\widetilde{Z}_s],0)+k_s^+)ds\\
            &\quad +\int_{t}^{T}\mathbf{1}_{\left\{\mathbf{E}[L(t,Y_t)]=0\right\}}(b(s,\mathbf{E}[\widetilde{Y}_s],\mathbf{E}[\widetilde{Z}_s],-k_s^-)-b(s,\mathbf{E}[\widetilde{Y}_s],\mathbf{E}[\widetilde{Z}_s],0)-k_s^-)ds\\
        \end{split}
    \end{equation} 
    \fi
\end{proof}
\begin{remark}
  Suppose that $f$ satisfying Assumption \ref{asslip2} has the following representation 
  \begin{align}\label{repre of f}f(t,\omega,y,y',z,z',k)=a(t,\omega,y,y',z,z')+b(t,y',z',k).
  \end{align}
  For each $t\in[0,T]$, set $\bar{b}_t(k)=b(t,\mathbf{E}[\widetilde{Y}_t],\mathbf{E}[\widetilde{Z}_t],k)$. For simplicity, we omit $t$ in the subscript. Then, equations in  \eqref{k+k-} are solvable if one of the following  conditions holds:
\begin{itemize}
    \item [(i)] $\bar{b}$ is increasing in $k$;
    \item[(ii)] there exists a constant $C_k<1$, such that $\left|\bar{b}(k^1)-\bar{b}(k^2)\right|\leq C_k\left|k^1-k^2\right|$;
     %\item[(iii)] $\bar{b}$ is decreasing in $k$ and $\left|\bar{b}(k^1)-\bar{b}(k^2)\right|\leq C_k\left|k^1-k^2\right|$ for $k^1,k^2>0$, with $C_k>1$.
\end{itemize} 

For $k\geq 0$, we set $ \widetilde{b}(k)=\bar{b}(k)-\bar{b}(0)+k$ and $\widehat{b}(k)=\bar{b}(-k)-\bar{b}(0)-k$. Given a constant $c>0$, it remains to show that each of the following equations $\widetilde{b}(k)=c$ and $\widehat{b}(k)=-c$ admits a strictly positive solution. We only prove this fact for $\widetilde{b}$. It suffices to show that the mapping $\widetilde{b}:[0,\infty)\rightarrow[0,\infty)$ is a one-to-one correspondence. First, $\widetilde{b}(0)=0$ clearly holds for all cases. 

  Suppose that (i) holds. Then, it is easy to check that $\widetilde{b}$ is strictly increasing and $\lim_{k\rightarrow\infty}\widetilde{b}(k)=\infty$. % in this case that for $k\geq 0$, $\bar{b}$ is increasing in $k$, then $\widetilde{b}$ is still increasing in $k$.  

Suppose that (ii) holds. For $0\leq k < k'$, we have 
\begin{equation*}
    \begin{split}
       \widetilde{b}(k')-\widetilde{b}(k)&= \bar{b}(k')-\bar{b}(k)+k'-k\\
   &\geq -C_k\left|k'-k\right|+k'-k\\
   &=(1-C_k)(k'-k)> 0. 
    \end{split}
\end{equation*}
Then $\widetilde{b}$ is strictly increasing in $k$ and $\lim_{k\rightarrow\infty}\widetilde{b}(k)=\infty$.
%Thus $\widetilde{b}(k)$ is also increasing in $k$. Combining with $\widetilde{b}(0)=0$, there exists a positive solution $k_s^{+}$ such that $k_s^+=(\widetilde{b})^{-1}(\tilde{k}_s^+)\geq 0.$
%Since $\widetilde{b}(0)=0$ and $\widetilde{k}^+_t=0$ for each $t\in[0,T]$, we have 
%$k^+_t=0$. Combining with
%$\int_0^T \mathbf{E}[R(t,\widetilde{Y}_t)]\widetilde{k}^+_tdt=0$, we have $\int_0^T \mathbf{E}[R(t,\widetilde{Y}_t)]k^+_tdt=0$.

%If (iii) holds, in this case for $0\leq k \leq k'$, we have 
%\begin{equation*}
    %\begin{split}
       %\widetilde{b}(k')-\widetilde{b}(k)&= \bar{b}(k')-\bar{b}(k)+k'-k\\
   %&\leq -C_k(k'-k)+k'-k\\
   %&=(1-C_k)(k'-k)\leq 0. 
   % \end{split}
%\end{equation*}
%Thus $\widetilde{b}(k)$ is decreasing in $k$. 

%Combining $\widetilde{b}(0)=0$ with $\widetilde{b}(k)$ is increasing in $k$, there exists a positive solution $k_s^{+}$ such that $k_s^+=(\widetilde{b})^{-1}(\tilde{k}_s^+)\geq 0.$
%Since $\widetilde{b}(0)=0$ and $\widetilde{k}^+_t=0$ for each $t\in[0,T]$, we have 
%$k^+_t=0$.
\end{remark}

For uniqueness, we only deal with the case where the loss functions are linear. More precisely, consider the following doubly mean reflected MFBSDE with resistance:
\begin{equation}\label{drmfbsde}
\begin{cases}
Y_t=\xi+\int_t^T \left(f(s,Y_s,\mathbf{E}[Y_s],Z_s,\mathbf{E}[Z_s],k_s)+k_s \right)ds-\int_t^T Z_s dB_s, \\
l_t\leq \mathbf{E}[Y_t]\leq r_t, \ t\in[0,T],\\ 
{k_t=k^+_t-k^-_t, \ t\in[0,T] \textrm{ and }k^+,k^-\geq 0,}\\
\int_0^T (\mathbf{E}[Y_t]-l_t)k_t^{+}dt=\int_0^T (r_t-\mathbf{E}[Y_t])k_t^{-}dt=0.
\end{cases}
\end{equation}

\begin{proposition}\label{prounik}
Let Assumption \ref{asslip2} hold. Moreover, suppose that $f$ has representation \eqref{repre of f} and it satisfies $\frac{df}{dk}\geq-1$. % and the generator satisfies $\frac{df}{dk}\geq-1$.{\color{blue}Suppose 
%  \begin{equation*}
 %    \begin{split}
         %\alpha^{i,j}_s&:=\frac{f(s,Y^j_s,\mathbf{E}[Y^j_s],Z^j_s,\mathbf{E}[Z^j_s],k_s^i)-f(s,Y^j_s,\mathbf{E}[Y^j_s],Z^j_s,\mathbf{E}[Z^j_s],k_s^j)}{k_s^i-k_s^j}I_{\left\{k_s^i-k_s^j\ne 0\right\}}, \ i,j \in \mathbb{R}\\
 %        \end{split}   
  %  \end{equation*}
   % is deterministic, where $k^i,\ k^j$ are two density processes.} 
   Then the doubly mean reflected MFBSDE \eqref{drmfbsde} with resistance admits at most one unique solution.
\end{proposition}


\begin{proof}
    Suppose that $(Y^1,Z^1,k^1)$ and $(Y^2,Z^2,k^2)$ are two solutions to doubly mean reflected MFBSDE  \eqref{drmfbsde} with resistance. In order to simplify notations, denote
    $$\Delta K:=K^1-K^2,  \  f^{i,j}_s:=f(s,Y_s^i,\mathbf{E}[Y_s^i],Z_s^i,\mathbf{E}[Z_s^i],k_s^j),\text{ where }\ K= Y,\ Z,\ k,\ f;\ i,j=1,2$$
    and 
    \begin{equation*}
     \begin{split}
         \alpha^{1,2}_s:=\frac{f^{2,1}_s-f^{2,2}_s}{\Delta k_s} I_{\{\Delta k_s\neq 0\}}. %\frac{f(s,Y^j_s,\mathbf{E}[Y^j_s],Z^j_s,\mathbf{E}[Z^j_s],k_s^i)-f(s,Y^j_s,\mathbf{E}[Y^j_s],Z^j_s,\mathbf{E}[Z^j_s],k_s^j)}{k_s^i-k_s^j}I_{\left\{k_s^i-k_s^j\ne 0\right\}}.
         \end{split}   
    \end{equation*}
    
    Applying It\^o formula to $\left|\Delta Y \right|^2$ and taking expectation, we obtain 
    \begin{equation}
    \begin{split}
      & \mathbf{E}\left[\left|\Delta Y_t\right|^2\right]+\mathbf{E}\left[\int_{t}^{T}\left|\Delta Z_s\right|^2ds\right]\\
      =&2\mathbf{E}\left[\int_{t}^{T}\Delta Y_s \Delta f_s ds\right] +2\mathbf{E}\left[\int_{t}^{T}\Delta Y_s \Delta k_s ds\right]\\
      =&2\mathbf{E}\left[\int_{t}^{T}\Delta Y_s(f_s^{1,2}-f_{s}^{2,2})ds\right]+2\mathbf{E}\left[\int_{t}^{T}(1+\alpha^{1,2}_s)\Delta Y_s \Delta k_s ds\right]\\
      \leq & 2\lambda\mathbf{E}\left[\int_{t}^{T}\Delta Y_s(\left|\Delta Y_s \right|+\mathbf{E}[\left|\Delta Y_s\right|])ds\right]+2\lambda\mathbf{E}\left[\int_{t}^{T}\Delta Y_s(\left|\Delta Z_s\right|+\mathbf{E}[\left|\Delta Z_s\right|])ds\right]\\
      &\quad +\left[2\mathbf{E}\int_{t}^{T}(1+\alpha^{1,2}_s)\Delta Y_s \Delta k_s ds\right].\\
    \end{split}
    \end{equation}
    By the assumption $f$, $\alpha^{1,2}$ is a deterministic function and we have $1+\alpha^{1,2}_s\geq 0$. By the Skorokhod conditions, we have $k^{i,+}_s(\mathbf{E}[Y^i_s]-l_s)=0, \ a.e.$ and $k^{i,-}_s(r_s-\mathbf{E}[Y^i_s])=0, \ a.e.$, $i=1,2$.   Therefore,
  \begin{equation*}
      \begin{split}
          &\E\left[\int_t^T(1+\alpha_s^{1,2})\Delta Y_s \Delta k_s ds\right] =\int_t^T(1+\alpha_s^{1,2})\E\left[\Delta Y_s\right] \Delta k_sds\\
          =&\int_{t}^{T}(1+\alpha_s^{1,2})(k_s^{1,+}-k_s^{2,+})(\mathbf{E}[Y^1_s]-l_s+l_s-\mathbf{E}[Y^2_s])ds \\
           &+\int_{t}^{T}(1+\alpha_s^{1,2})(k_s^{2,-}-k_s^{1,-})(\mathbf{E}[Y^2_s]-r_s+r_s-\mathbf{E}[Y^1_s])ds \\
           =&-\int_{t}^{T}(1+\alpha_s^{1,2})(\mathbf{E}[Y^2_t]-l_s)k_s^{1,+}ds-\int_{t}^{T}(1+\alpha_s^{1,2})(\mathbf{E}[Y^2_t]-l_s)k_s^{2,+}ds\\
           &-\int_{t}^{T}(1+\alpha_s^{1,2})(r_s-\mathbf{E}[Y^2_t])k_s^{1,-}ds-\int_{t}^{T}(1+\alpha_s^{1,2})(r_s-\mathbf{E}[Y^2_t])k_s^{2,-}ds\leq 0.
      \end{split}
  \end{equation*}
   All the above analysis indicates that 
 \begin{equation}
    \begin{split}
       \mathbf{E}\left[\left|\Delta Y_t\right|^2\right]+\mathbf{E}\left[\int_{t}^{T}\left|\Delta Z_s\right|^2ds\right]
      \leq  (4\lambda+4\lambda^2)\mathbf{E}\left[\int_{t}^{T}\left|\Delta Y_s \right|^2 ds\right]+\mathbf{E}\left[\int_{t}^{T}\left|\Delta Z_s\right|^2ds\right].%+2\lambda^2\mathbf{E}\left[\int_{t}^{T}\left|\Delta Y_s \right|^2 ds \right]+ \frac{1}{2}\mathbf{E}\left[\int_{t}^{T}\left|\Delta Z_s\right|^2ds\right]\\
     % &\quad +2\lambda^2\mathbf{E}\left[\int_{t}^{T}\left|\Delta Y_s \right|^2 ds\right] + \frac{1}{2}\mathbf{E}\left[\int_{t}^{T}\left|\Delta Z_s\right|^2ds\right].\\
    \end{split}
    \end{equation}
%where we use H\"older inequality and Young inequality. 
It follows that 
\begin{equation*}  
\mathbf{E}\left[|\Delta Y_t|^2\right]\leq (4\lambda+4\lambda^2)\int_{t}^{T}\mathbf{E}\left[|\Delta Y_t|^2\right]ds.
\end{equation*}
By Gronwall inequality, $Y^1_t=Y^2_t$, on $[0,T]$. Applying It\^{o}'s formula to $(Y^1_t-Y^2_t)^2$, we obtain that $Z^1_t=Z^2_t$ on $t\in[0,T]$ and finally, we have $k^1_t=k^2_t$ on $[0,T]$.
\end{proof}
%{\color{blue} \begin{remark}
%If the generator $$f(t,\omega, y,y',z,z')=a(t,\omega,y,y',z,z')+b(t,y',z',k),$$
%in this case, $\left\{\alpha^{i,j}_s\right\}_{s\in[0,T]}$ is deterministic. Thus, we can easily check that Eq.\eqref{drmfbsde} has unique solution.
%\end{remark}}

\iffalse
\begin{proof}
    Suppose that $(Y^1,Z^1,k^1)$ and $(Y^2,Z^2,k^2)$ are two solutions to doubly mean reflected MFBSDE  \eqref{drmfbsde} {\color{red}with resistance}.
    Applying It\^o formula to $\left|Y^1-Y^2\right|^2$ and taking expectation, we can obtain 
    \begin{equation}
    \begin{split}
      & \mathbf{E}\left|Y_t^1-Y_t^2\right|^2+\mathbf{E}\int_{t}^{T}\left|Z_s^1-Z_s^2\right|^2ds\\
      =&2\mathbf{E}\int_{t}^{T}(Y^1_s-Y^2_s)(f(s,Y_s^1,\mathbf{E}[Y_s^1],Z_s^1,\mathbf{E}[Z_s^1],k_s^1)-f(s,Y_s^2,\mathbf{E}[Y_s^2],Z_s^2,\mathbf{E}[Z_s^2],k_s^2))ds\\
      &\quad+2\mathbf{E}\int_{t}^{T}(Y^1_s-Y^2_s)(k_s^1-k_s^2)ds\\
      =&2\mathbf{E}\int_{t}^{T}(Y^1_s-Y^2_s)(\alpha_s(Y_s^1-Y_s^2)+\beta_s(Z_s^1-Z_s^2)+\gamma_s(k_s^1-k_s^2))ds.\\
      &\quad+2\mathbf{E}\int_{t}^{T}(Y^1_s-Y^2_s)(k_s^1-k_s^2)ds\\
    \end{split}
    \end{equation}
    where we denote that 
    \begin{equation*}
     \begin{split}
         \alpha_s&=\frac{f(s,Y_s^1,\mathbf{E}[Y_s^1],Z_s^1,\mathbf{E}[Z_s^1],k_s^1)-f(s,Y_s^2,\mathbf{E}[Y_s^2],Z_s^1,\mathbf{E}[Z_s^1],k_s^1)}{Y_s^1-Y_s^2}I_{\left\{Y_s^1-Y_s^2\ne 0\right\}} \\
         \beta_s&=\frac{f(s,Y_s^2,\mathbf{E}[Y_s^2],Z_s^1,\mathbf{E}[Z_s^1],k_s^1)-f(s,Y_s^2,\mathbf{E}[Y_s^2],Z_s^2,\mathbf{E}[Z_s^2],k_s^1)}{Z_s^1-Z_s^2}I_{\left\{Z_s^1-Z_s^2\ne 0\right\}} \\
         \gamma_s&=\frac{f(s,Y_s^2,\mathbf{E}[Y_s^2],Z_s^2,\mathbf{E}[Z_s^2],k_s^1)-f(s,Y_s^2,\mathbf{E}[Y_s^2],Z_s^2,\mathbf{E}[Z_s^2],k_s^2)}{k_s^1-k_s^2}I_{\left\{k_s^1-k_s^2\ne 0\right\}} \\
         \end{split}   
    \end{equation*}
    By the Skorokhod conditions, we have
  \begin{equation*}
      \begin{split}
          \E\left[\int_t^T(Y^1_s-Y^2_s)(k^1_s-k^2_s) ds\right]
          &=\int_t^T\E\left[Y^1_s-Y^2_s\right](k^1_s-k^2_s)ds\\
          &=\int_{t}^{T}(k_s^{1,+}-k_s^{2,+})(\mathbf{E}[Y^1_s]-l_s+l_s-\mathbf{E}[Y^2_s])ds \\
           &\quad+\int_{t}^{T}(k_s^{2,-}-k_s^{1,-})(\mathbf{E}[Y^2_s]-r_s+r_s-\mathbf{E}[Y^1_s])ds \leq 0.
      \end{split}
  \end{equation*}
       Combining with Assumption \ref{asslip2}, we have 
\begin{equation}
    \begin{split}
      & \mathbf{E}\left|Y_t^1-Y_t^2\right|^2+\mathbf{E}\int_{t}^{T}\left|Z_s^1-Z_s^2\right|^2ds\\
      &\leq 2\int_{t}^{T} \mathbf{E}[(Y^1_s-Y^2_s)(\lambda(Y^1_s-Y^2_s)+\lambda\mathbf{E}[\left|Y_s^1-Y_s^2\right|]\frac{Y_s^1-Y_s^2}{\left|Y_s^1-Y_s^2\right|})]ds\\
      &\quad +2 \int_{t}^{T}\mathbf{E}[(Y^1_s-Y^2_s)(\lambda(Z^1_s-Z^2_s)+\lambda\mathbf{E}[\left|Z_s^1-Z_s^2\right|]\frac{Z_s^1-Z_s^2}{\left|Z_s^1-Z_s^2\right|})]ds\\
      &\leq (4\lambda+4\lambda^2) \int_{t}^{T} \mathbf{E}[\left|Y_s^1-Y_s^2\right|^2]ds+\int_{t}^{T} \mathbf{E}[\left|Z_s^1-Z_s^2\right|^2]ds
    \end{split}
\end{equation} 
  
All the above analysis indicates that 
\begin{equation*}  
\mathbf{E}\left[|Y_t^1-Y_t^2|^2\right]\leq (4\lambda+4\lambda^2)\int_{t}^{T}\mathbf{E}\left[|Y_s^1-Y_s^2|^2\right]ds.
\end{equation*}
By Gronwall inequality, $Y^1_t=Y^2_t$, on $[0,T]$, thus $Z^1_t=Z^2_t$ and $k^1_t=k^2_t$ on $[0,T]$.
\end{proof}
 Suppose that $(Y^1,Z^1,k^1)$ and $(Y^2,Z^2,k^2)$ are two solutions to doubly mean reflected MFBSDE  \eqref{drmfbsde} {\color{red}with resistance}.
    Applying It\^o formula to $\left|Y^1-Y^2\right|^2$ and taking expectation, we can obtain 
    \begin{equation}
    \begin{split}
      & \mathbf{E}\left|Y_t^1-Y_t^2\right|^2+\mathbf{E}\int_{t}^{T}\left|Z_t^1-Z_t^2\right|^2ds\\
      =&2\mathbf{E}\int_{t}^{T}(Y^1_s-Y^2_s)(f(s,Y_s^1,\mathbf{E}[Y_s^1],Z_s^1,\mathbf{E}[Z_s^1],k_s^1)-f(s,Y_s^2,\mathbf{E}[Y_s^2],Z_s^2,\mathbf{E}[Z_s^1],k_s^2))ds\\
      &\quad+2\mathbf{E}\int_{t}^{T}(Y^1_s-Y^2_s)(k_s^1-k_s^2)ds\\
      =&2\mathbf{E}\int_{t}^{T}(Y^1_s-Y^2_s)(\alpha_s(Y_s^1-Y_s^2)+\beta_s(Z_s^1-Z_s^2)+\gamma_s(k_s^1-k_s^2))ds.\\
      &\quad+2\mathbf{E}\int_{t}^{T}(Y^1_s-Y^2_s)(k_s^1-k_s^2)ds\\
    \end{split}
    \end{equation}
    where we denote that 
    \begin{equation*}
     \begin{split}
         \alpha_s&=\frac{f(s,Y_s^1,\mathbf{E}[Y_s^1],Z_s^1,\mathbf{E}[Z_s^1],k_s^1)-f(s,Y_s^2,\mathbf{E}[Y_s^2],Z_s^1,\mathbf{E}[Z_s^1],k_s^1)}{Y_s^1-Y_s^2}I_{\left\{Y_s^1-Y_s^2\ne 0\right\}} \\
         \beta_s&=\frac{f(s,Y_s^2,\mathbf{E}[Y_s^2],Z_s^1,\mathbf{E}[Z_s^1],k_s^1)-f(s,Y_s^2,\mathbf{E}[Y_s^2],Z_s^2,\mathbf{E}[Z_s^2],k_s^1)}{Z_s^1-Z_s^2}I_{\left\{Z_s^1-Z_s^2\ne 0\right\}} \\
         \gamma_s&=\frac{f(s,Y_s^2,\mathbf{E}[Y_s^2],Z_s^2,\mathbf{E}[Z_s^2],k_s^1)-f(s,Y_s^2,\mathbf{E}[Y_s^2],Z_s^2,\mathbf{E}[Z_s^2],k_s^2)}{k_s^1-k_s^2}I_{\left\{k_s^1-k_s^2\ne 0\right\}} \\
         \end{split}   
    \end{equation*}
    By the Skorokhod conditions, we have
  \begin{equation*}
      \begin{split}
          \E\left[\int_t^T(Y^1_s-Y^2_s)(k^1_s-k^2_s) ds\right]
          &=\int_t^T\E\left[Y^1_s-Y^2_s\right](k^1_s-k^2_s)ds\\
          &=\int_{t}^{T}(k_s^{1,+}-k_s^{2,+})(\mathbf{E}[Y^1_s]-l_s+l_s-\mathbf{E}[Y^2_s])ds \\
           &\quad+\int_{t}^{T}(k_s^{2,-}-k_s^{1,-})(\mathbf{E}[Y^2_s]-r_s+r_s-\mathbf{E}[Y^1_s])ds \leq 0.
      \end{split}
  \end{equation*}
       Combining with Assumption \ref{asslip2}, we have 
\begin{equation}
    \begin{split}
      & \mathbf{E}\left|Y_t^1-Y_t^2\right|^2+\mathbf{E}\int_{t}^{T}\left|Z_t^1-Z_t^2\right|^2ds\\
      &\leq 2\int_{t}^{T} \mathbf{E}[(Y^1_s-Y^2_s)(\lambda(Y^1_s-Y^2_s)+\lambda\mathbf{E}[\left|Y_s^1-Y_s^2\right|]\frac{Y_s^1-Y_s^2}{\left|Y_s^1-Y_s^2\right|})]ds\\
      &\quad +2 \int_{t}^{T}\mathbf{E}[(Y^1_s-Y^2_s)(\lambda(Z^1_s-Z^2_s)+\lambda\mathbf{E}[\left|Z_s^1-Z_s^2\right|]\frac{Z_s^1-Z_s^2}{\left|Z_s^1-Z_s^2\right|})]ds\\
      &\leq (4\lambda+4\lambda^2) \int_{t}^{T} \mathbf{E}[\left|Y_s^1-Y_s^2\right|^2]ds+\int_{t}^{T} \mathbf{E}[\left|Z_s^1-Z_s^2\right|^2]ds
    \end{split}
\end{equation} 

\begin{proof}
    The proof is similar to Proposition 3.3 in \cite{LX}. For the convenience of reading, we provide a full proof below.

    For $v \in [0,T]$ and $t\in [v,T]$,
    \begin{equation}\begin{split}
       Y_t&=\mathbf{E}\left[\xi+\int_{t}^{T}f(s,Y_s, \mathbf{P}_{Y_s},Z_s,\mathbf{P}_{Z_s})ds|\mathcal{F}_t\right]\\
       &\leq \mathbf{E}\left[\left|\xi\right|+\int_{v}^{T}\left|f(s,Y_s, \mathbf{P}_{Y_s},Z_s,\mathbf{P}_{Z_s})\right|ds|\mathcal{F}_t\right].  
    \end{split}  
    \end{equation}
Therefore, we have 
\begin{equation}
    \begin{split}
\mathbf{E}\left[\sup_{t\in[v,T]}\left|Y_t\right|^p\right]&\leq \mathbf{E}\left[\sup_{t\in[v,T]}|\mathbf{E}[\left|\xi\right|+\int_{v}^{T}\left|f(s,Y_s, \mathbf{P}_{Y_s},Z_s,\mathbf{P}_{Z_s})\right|ds|\mathcal{F}_t]|^p\right]\\
&\leq \left(\frac{p}{p-1}\right)^p \mathbf{E}\left[\left(\mathbf{E}[\left|\xi\right|+\int_{v}^{T}\left|f(s,Y_s, \mathbf{P}_{Y_s},Z_s,\mathbf{P}_{Z_s})\right|ds|\mathcal{F}_T]|]\right)^p\right]\\
&\leq \left(\frac{p}{p-1}\right)^p \mathbf{E}\left[\mathbf{E}\left[\left(\left|\xi\right|+\int_{v}^{T}\left|f(s,Y_s, \mathbf{P}_{Y_s},Z_s,\mathbf{P}_{Z_s})\right|ds\right)^p|\mathcal{F}_T\right]\right]\\ 
&\leq \left(\frac{p}{p-1}\right)^p 2^{p-1} \mathbf{E}\left[|\xi|^p+\left(\int_{v}^{T}\left|f(s,Y_s, \mathbf{P}_{Y_s},Z_s,\mathbf{P}_{Z_s})\right|ds\right)^p\right]
\end{split}
\end{equation}
Since $f$ satisfies Lipschitz continuous, we have 
\begin{align}
    \left|f(s,Y_s, \mathbf{P}_{Y_s},Z_s,\mathbf{P}_{Z_s})\right|\leq \left|f(s,0, \delta_0,0,\delta_0)\right|+\lambda \left(\left|Y_s\right|+\left|Z_s\right|+\mathbf{E}[|Y_s|]+\mathbf{E}[|Z_s|]\right) 
\end{align}
Thus, 
\begin{align*}
    &\mathbf{E}\left[\left(\int_{v}^{T}\left|f(s,Y_s, \mathbf{P}_{Y_s},Z_s,\mathbf{P}_{Z_s})\right|ds\right)^p\right]\\
    &\leq 5^{p-1}\mathbf{E}\left[\left(\int_{v}^{T}\left|f(s,0, \delta_0,0,\delta_0)\right|ds\right)^p\right]+5^{p-1}\lambda^p\mathbf{E}\left[\left(\int_{v}^{T}\left|Y_s\right|ds\right)^p\right]\\
    &\quad+5^{p-1}\lambda^p\mathbf{E}\left[\left(\int_{v}^{T}\left|Z_s\right|ds\right)^p\right]+5^{p-1}\lambda^p\left(\int_{v}^{T}\mathbf{E}\left|Y_s\right|ds\right)^p+5^{p-1}\lambda^p\left(\int_{v}^{T}\mathbf{E}\left|Z_s\right|ds\right)^p\\
    &\leq 5^{p-1}\mathbf{E}\left[\left(\int_{v}^{T}\left|f(s,0, \delta_0,0,\delta_0)\right|ds\right)^p\right]+2\cdot 5^{p-1}\lambda^p(T-v)^{p-1}\mathbf{E}\left[\int_{v}^{T}\sup_{r\in[s,T]}|Y_r|^pds\right]\\
    &\quad + 2\cdot 5^{p-1}\lambda^p(T-v)^{\frac{p}{2}}\mathbf{E}\left[\left(\int_{v}^{T}|Z_s|^2ds\right)^{\frac{p}{2}}\right]
\end{align*}
In addition, according to B-D-G inequality, we have 
\begin{align*}
&\mathbf{E}\left[\left(\int_{v}^{T}|Z_s|^2ds\right)^{\frac{p}{2}}\right]\\
& \leq C_{p}\mathbf{E}\left[\sup_{t\in[v,T]}\left|\int_{v}^{t}Z_sdB_s\right|^p\right]\\
&\leq C_{p}\left(\frac{p}{p-1}\right)^p\mathbf{E}\left[\left|\int_{v}^{T}Z_sdB_s\right|^p\right]\\
&\leq C_{p}\left(\frac{p}{p-1}\right)^p\mathbf{E}\left[\left|\xi-Y_v+\int_{v}^{T}f(s,Y_s,\mathbf{P}_{Y_s},Z_s,\mathbf{P}_{Z_s})ds\right|^p\right]\\
&\leq C_{p}\left(\frac{p}{p-1}\right)^p 3^{p-1}\mathbf{E}\left[\left|\xi\right|^p+\left|Y_v\right|^p+\left(\int_{v}^{T}\left|f(s,Y_s, \mathbf{P}_{Y_s},Z_s,\mathbf{P}_{Z_s})\right|ds\right)^p\right]\\
&\leq  C_{p}\left(\frac{p}{p-1}\right)^p 3^{p-1}\mathbf{E}\left[\left|\xi\right|^p\right]+ C_{p}\left(\frac{p}{p-1}\right)^p 3^{p-1}\mathbf{E}\left[\sup_{r\in[v,T]}\left|Y_r\right|^p\right]\\
&\quad+ C_{p}\left(\frac{p}{p-1}\right)^p 3^{p-1}\mathbf{E}\left[\left(\int_{v}^{T}\left|f(s,Y_s, \mathbf{P}_{Y_s},Z_s,\mathbf{P}_{Z_s})\right|ds\right)^p\right]\\
&\leq C_{p}\left(\frac{p}{p-1}\right)^p 3^{p-1}\mathbf{E}\left[\left|\xi\right|^p\right]+C_{p}\left(\frac{p}{p-1}\right)^p 3^{p-1}5^{p-1}\mathbf{E}\left[\left(\int_{v}^{T}\left|f(s,0, \delta_0,0,\delta_0)\right|ds\right)^p\right]\\
&\quad+C_{p}\left(\frac{p}{p-1}\right)^p 3^{p-1}\mathbf{E}\left[\sup_{r\in[v,T]}\left|Y_r\right|^p\right]\\
&\quad +C_{p}^{*}(T-v)^{p-1}\mathbf{E}\left[\int_{v}^{T}\sup_{r\in[s,T]}|Y_r|^pds\right]+ C_{p}^{*}(T-v)^{\frac{p}{2}}\mathbf{E}\left[\left(\int_{v}^{T}|Z_s|^2ds\right)^{\frac{p}{2}}\right]\\
\end{align*}
where $C_{p}^{*}=2C_p\left(\frac{p}{p-1}\right)^p3^{p-1}5^{p-1}\lambda^p$.
Define $\delta:=T-v$, we choose an appropriate $\delta$ such that $C_{p}^{*}(T-v)^{\frac{p}{2}}<1$. Thus, for $v\in[T-\delta,T]$,
\begin{align*}
&\mathbf{E}\left[\left(\int_{v}^{T}|Z_s|^2ds\right)^{\frac{p}{2}}\right]\\
&\leq \frac{C_p\left(\frac{p}{p-1}\right)^p3^{p-1}}{1-C_{p}^{*}(T-v)^{\frac{p}{2}}}\mathbf{E}\left[\left|\xi\right|^p\right]+ \frac{C_p\left(\frac{p}{p-1}\right)^p3^{p-1}5^{p-1}}{1-C_{p}^{*}(T-v)^{\frac{p}{2}}}\mathbf{E}\left[\left(\int_{v}^{T}\left|f(s,0, \delta_0,0,\delta_0)\right|ds\right)^p\right]\\
&\quad + \frac{C_p\left(\frac{p}{p-1}\right)^p3^{p-1}}{1-C_{p}^{*}(T-v)^{\frac{p}{2}}}\mathbf{E}\left[\sup_{r\in[v,T]}\left|Y_r\right|^p\right] + \frac{C_{p}^{*}(T-v)^{p-1}}{1-C_{p}^{*}(T-v)^{\frac{p}{2}}}\mathbf{E}\left[\int_{v}^{T}\sup_{r\in[s,T]}|Y_r|^pds\right]\\   
\end{align*}
Therefore, combining with , for $\delta$ satisfies $(\frac{p}{p-1})^p2^{p-1}\frac{C_{p}^{*}(T-v)^{\frac{p}{2}}}{1-C_{p}^{*}(T-v)^{\frac{p}{2}}}<1$, we have 
\begin{align*}
\mathbf{E}\left[\sup_{t\in[v,T]}\left|Y_t\right|^p\right]\leq C_p^{'}\mathbf{E}\left[\left|\xi\right|^p+\left(\int_{v}^{T}\left|f(s,0,\delta_0,0,\delta_0)\right|ds\right)^p\right]+C_p^{'}\mathbf{E}\left[\int_{v}^{T}\sup_{r\in[s,T]}\left|Y_r\right|^pds\right]
\end{align*}
Applying Gronwall inequality, we have 
\begin{align*}
\mathbf{E}\left[\sup_{t\in[v,T]}\left|Y_t\right|^p\right]\leq C_p^{'}\mathbf{E}\left[\left|\xi\right|^p+\left(\int_{0}^{T}\left|f(s,0,\delta_0,0,\delta_0)\right|ds\right)^p\right], \ \forall v\in[T-\delta,T].
\end{align*}
Similarly, we have 
\begin{align*}
\mathbf{E}\left[\sup_{t\in[v,T]}\left|Y_t\right|^p\right]&
\leq C_p^{'}\mathbf{E}\left[\left|Y_{T-\delta}\right|^p+\left(\int_{0}^{T}\left|f(s,0,\delta_0,0,\delta_0)\right|ds\right)^p\right]\\
&\leq C_p^{'}\mathbf{E}\left[\left|\xi\right|^p+\left(\int_{0}^{T}\left|f(s,0,\delta_0,0,\delta_0)\right|ds\right)^p\right], \ \forall v\in[T-2\delta,T-\delta].\\
\end{align*}
Therefore, for $v\in[T-2\delta,T]$
\begin{align*}
\mathbf{E}\left[\sup_{t\in[v,T]}\left|Y_t\right|^p\right]
&\leq C_p^{'}\mathbf{E}\left[\left|\xi\right|^p+\left(\int_{0}^{T}\left|f(s,0,\delta_0,0,\delta_0)\right|ds\right)^p\right].\\
\end{align*}
Iterating this argument $N$ times such that $N\delta>T$, we have 
\begin{align*}
\mathbf{E}\left[\sup_{t\in[0,T]}\left|Y_t\right|^p\right]
&\leq C_p^{'}\mathbf{E}\left[\left|\xi\right|^p+\left(\int_{0}^{T}\left|f(s,0,\delta_0,0,\delta_0)\right|ds\right)^p\right].
\end{align*}
\end{proof}
Next, for $T\in(0,\varepsilon]$, we have the following theorem for small enough time horizon $[0,T]$.
\begin{theorem}\label{th3-1}
Let Assumptions $\ref{ass1ipterminalp}-\ref{assgk}$ hold. For $T\in(0,\varepsilon]$, the doubly mean reflected MFBSDE \eqref{nonlinearyz} with nonlinear resistance admits a unique solution $(Y,Z,K)\in \mathcal{S}^p \times \mathcal{H}^{p} \times BV[0,T].$   
\end{theorem}
\fi

\section*{Acknowledgement}
    This work was supported  by the National Natural Science Foundation of China (No. 12301178), the Natural Science Foundation of Shandong Province for Excellent Young Scientists Fund Program (Overseas) (No. 2023HWYQ-049),  the Natural Science Foundation of Shandong Province (No. ZR2023ZD35),  the Fundamental Research Funds for the Central Universities and  the Qilu Young Scholars Program of Shandong University. 
%\end{acknowledgement}

\iffalse
\renewcommand\thesection{Appendix A}
\section{Construction by penalization in a special case}\label{append:A}
\renewcommand\thesection{A}

\iffalse
{\color{red}$\bar{b}$ is increasing is not enough. We need to make sure that $\widetilde{b}(0)=0$ (it clearly holds), $\lim_{k\rightarrow+\infty}\widetilde{b}(k)=+\infty$ and $\widetilde{b}$ is continuous in $k$.}
\subsection{Mean field $\text{BSDEs}$}
  Define $\hat{f}(s,y,z,\mu,k):=f(s,y,z,\mu,k)+k$.  Given two processes $k^1,k^2$, it is easy to check that $(Y^1,Z^1)$ and $(Y^2,Z^2)$ are solutions to the following MFBSDE, respectively,
    \begin{equation*}
    \begin{split}
      &Y_t^1=\xi+\int_{t}^{T} \hat{f}(s,Y_s^1,\mathbf{E}[Y_s^1],Z_s^1,k_s^1)ds-\int_{t}^{T}Z^1_sdB_s,\\ 
       &Y_t^2=\xi+\int_{t}^{T} \hat{f}(s,Y_s^2,\mathbf{E}[Y_s^2],Z_s^2,k_s^2)ds-\int_{t}^{T}Z^2_sdB_s.\\   
    \end{split} 
    \end{equation*}
    Since $\frac{df}{dk}>-1$, {\color{red}$\hat{f}$} is increasing in $k$. Therefore, if $k^1_t\geq k^2_t$, then $\hat{f}(t,y,z,\mu,k^1_t)\geq \hat{f}(t,y,z,\mu,k^2_t)$. According to Theorem 3.5 in \cite{CXZ}, we have $Y^1_t\geq Y^2_t$. Similarly, if $k^1_t\leq k^2_t$, then $Y^1_t\leq Y^2_t$. Thus, for any $t\in[0,T]$, the following equation always holds true:
    \begin{equation}\label{k1k2y1y2}
        (k^1_t-k^2_t)(Y^1_t-Y^2_t)\geq 0.
    \end{equation}
     {\color{red}  By the Lipschitz assumption for $f$, it is easy to check that 
    \begin{align*}
        & \mathbf{E}\left[|Y_t^1-Y_t^2|^2\right]+\mathbf{E}\left[\int_{t}^{T}|Z_s^1-Z_s^2|^2ds\right]\\
        \leq &(2\lambda+\lambda^2)\E\left[\int_t^T|Y^1_s-Y^2_s|^2ds\right]+\mathbf{E}\left[\int_{t}^{T}|Z_s^1-Z_s^2|^2ds\right]\\
        &+2(\lambda+1)\E\left[\int_t^T|(Y^1_s-Y^2_s)(k^1_s-k^2_s)|ds\right].
    \end{align*}
 
  }
 %   \begin{equation}
  %      \begin{split}
   %        & \int_{0}^{T}(k_s^1-k_s^2)(\mathbf{E}[Y^1_s]-\mathbf{E}[Y^2_s])ds \\
    %       &=\int_{0}^{T}(k_s^{1,+}-k_s^{2,+})(\mathbf{E}[Y^1_s]-l_s+l_s-\mathbf{E}[Y^2_s])ds \\
     %      &\quad+\int_{0}^{T}(k_s^{2,-}-k_s^{1,-})(\mathbf{E}[Y^2_s]-r_s+r_s-\mathbf{E}[Y^1_s])ds \leq 0. \\
     %   \end{split}
    %\end{equation}
% Therefore, we have $(k_t^1-k_t^2)(Y^1_t-Y^2_t)=0,\ a.e.$ on $[0,T]$. Furthermore, 
Roughly speaking, the mean field of $\text{BSDE}$ is a kind of $\text{BSDE}$ whose driver depends on the law of the solution. More precisely, we consider the following type of mean field $\text{BSDE}$
\begin{align}\label{MFBSDE}
Y_t=\xi+\int_t^T f(s,\P_{(Y_s,Z_s)},Y_s,Z_s)ds-\int_t^T Z_s dB_s,
\end{align}
where $\P_{(Y_s,Z_s)}$ is the joint law of $(Y_s,Z_s)$. The driver $f:[0,T]\times\Omega\times \mathcal{P}_2(\mathbb{R}^2)\times \mathbb{R}\times\mathbb{R}\rightarrow \mathbb{R}$ is assumed to satisfy the following assumptions:

\begin{itemize}
\item[(H1)] $\{f(t,\omega,\delta_{\textbf{0}},0,0)\}_{t\in[0,T]}\in \mathcal{H}^2$, where $\delta_{\textbf{0}}$ is the Dirac measure with mass at $\textbf{0}\in \mathbb{R}^2$.
\item[(H2)] There exists a constant $L>0$, such that for all $\mu,\mu'\in\mathcal{P}_2(\mathbb{R}^2)$, $y,y',z,z'\in\mathbb{R}$,
\begin{align*}
|f(t,\mu,y,z)-f(t,\mu',y',z')|\leq L(W_2(\mu,\mu')+|y-y'|+|z-z'|).
\end{align*}
\end{itemize}

\begin{theorem}[\cite{LLZ}]\label{existence and uniqueness for MFBSDE}
Let $\xi\in L^2(\mathcal{F}_T)$. Suppose that $f$ satisfies (H1) and (H2). Then, the mean field BSDE \eqref{MFBSDE} has a unique solution $(Y,Z)\in \mathcal{S}^2\times \mathcal{H}^2$.
\end{theorem}

To investigate the comparison theorem for mean field BSDEs, the driver $f$ is assumed to be independent from the law of $Z$. That is, we consider the following type of mean field BSDE:
\begin{align}\label{MFBSDE'}
Y_t=\xi+\int_t^T f(s,\P_{Y_s},Y_s,Z_s)ds-\int_t^T Z_s dB_s.
\end{align}

\begin{theorem}[\cite{LLZ}]\label{comparison for MFBSDE}
Let $f_i=f_i(t,\omega,\mu,y,z)$, $i=1,2$ be two drivers satisfying (H1). Moreover, we suppose that one of the drivers satisfies (H2) and for any $\theta_1,\theta_2\in L^2(\mathcal{F}_T)$ and any $(t,y,z)\in[0,T]\times \mathbb{R}\times\mathbb{R}$
\begin{align*}
f_i(s,\P_{\theta_1},y,z)-f_i(s,\P_{\theta_2},y,z)\leq L(\E[|(\theta_1-\theta_2)^+|^2])^{1/2}.
\end{align*}
Let $\xi_i\in L^2(\mathcal{F}_T)$, $i=1,2$ and denote by $(Y^i,Z^i)$ the solution of the mean field BSDE \eqref{MFBSDE'} with parameters $(\xi_i,f_i)$, $i=1,2$. Then, if $\xi_1\leq \xi_2$ and $f_1(t,\mu,y,z)\leq f_2(t,\mu,y,z)$ for all $(\mu,y,z)$, we have $Y^1_t\leq Y^2_t$, $t\in[0,T]$.
%\begin{itemize}
%\item[(i)] one of the drivers satisfies (H2),
%\item[(ii)] one of the drivers satisfies
\end{theorem}
\fi


In this section, we apply a penalization method to construct the solution to the $\text{MFBSDE}$ with double mean reflections. However, due to the fact that the constraints only made for the distribution of $Y$ but not for the pointwise value of $Y$, we cannot make sure that the sequence of penalized MFBSDEs is monotone since the comparison theorem may not hold. For this reason, we only consider the following $\text{MFBSDE}$ with two linear reflections, whose parameters are denoted by $(\xi,f,l,r)$,
\begin{equation}\label{linearcase}
\begin{cases}
Y_t=\xi+\int_t^T f\left(s,Y_s,\P_{Y_s},Z_s,\P_{Z_s}\right)ds-\int_t^T Z_s dB_s+K_T-K_t, \\
l_t\leq \E\left[Y_t\right]\leq r_t, \\
K_t=K^l_t-K^r_t, \ K^l,K^r\in I[0,T]\textrm{ and } \int_0^T \left(\E[Y_t]-l_t\right)dK_t^l=\int_0^T \left(r_t-\E[Y_t]\right)dK^r_t=0.
\end{cases}
\end{equation}
Here, we assume that the obstacles $l,r$ satisfy the following assumption.

\begin{itemize}
\item[($H_{rl}$)] $r_t=\int_0^t a_s ds$, $l_t=\int_0^t b_s ds$ with $\int_0^T|a_t|^2 dt <\infty$, $\int_0^T|b_t|^2 dt<\infty$ and $l_t\leq r_t$, $t\in[0,T]$.
\end{itemize}

%\begin{definition}
%A triple of processes $(Y,Z,K)\in \mathcal{S}^2\times \mathcal{H}^2\times \mathcal{A}_{BV}$ with $K=K^l-K^r$ and $K^l,K^r\in \mathcal{A}_I$ is called the solution to the BSDE with two linear reflecting barriers with terminal value $\xi$, driver (or coefficient) $f$, barriers $l,r$ if \eqref{linearcase} is satisfied.
%\end{definition}



Now, consider the following penalized $\text{MFBSDEs}$:
\begin{equation}\begin{split}\label{panelization}
Y_t^n=&\xi+\int_t^T f(s,Y_s^n,\P_{Y_s^n},Z_s^n,\P_{Z^n_s})ds-\int_t^T Z_s^n dB_s\\
&+\int_t^T n(\E[Y_s^n]-l_s)^-ds-\int_t^T n(\E[Y_s^n]-r_s)^+ ds.
\end{split}\end{equation}
By Theorem 2.1 in \cite{LLZ}, the above equation admits a unique pair of solution $(Y^n,Z^n)\in \mathcal{S}^2\times \mathcal{H}^2$. We define $K^{n,l}_t:=\int_0^t n(\E[Y_s^n]-l_s)^- ds$, $K^{n,r}_t:=\int_0^t n(\E[Y_s^n]-r_s)^+ds$ and $K_t^n:=K^{n,l}_t-K^{n,r}_t$. We show that the solution to \eqref{linearcase} can be approximated by the solutions to \eqref{panelization}.

\begin{theorem}\label{existence and uniqueness}
Suppose that $f$ satisfies Assumption \ref{ass1} with $p=2$ and $r,l$ satisfy condition ($H_{lr}$). Given $\xi\in \mathcal{L}^2_T$ with $l_T\leq \E[\xi]\leq r_T$, the $\text{MFBSDE}$ with double mean reflections \eqref{linearcase} has a unique solution $(Y,Z,K)\in \mathcal{S}^2\times\mathcal{H}^2\times BV[0,T]$. Furthermore, $(Y,Z,K)$ is the limit of $(Y^n,Z^n,K^n)$.
\end{theorem}

Before proving Theorem \ref{existence and uniqueness}, we first present some a priori estimates similar with the classical reflected BSDEs.%, which gives an alternative way to show the uniqueness of the solution to doubly mean reflected BSDE \eqref{linearcase}. 
\begin{proposition}\label{uniquenesslinear}
Suppose that $r,l$ satisfy ($H_{rl}$) and $f^i$ satisfy Assumption \ref{ass1} with $p=2$, $i=1,2$. Given $\xi^i\in \mathcal{L}^2_T$ with $l_T\leq \E[\xi]\leq r_T$, let $(Y^i,Z^i,K^i)$ be the solution to the doubly mean reflected $\text{MFBSDE}$ with parameters $(\xi^i,f^i,l,r)$. Then, there exists a constant $C$ depending on $\lambda,T$, such that 
\begin{align*}
\E\left[\int_0^T |\hat{Y}_t|^2dt\right]+\E\left[\int_0^T |\hat{Z}_t|^2 dt\right]\leq C\E\left[|\hat{\xi}|^2+\int_0^T |\hat{f}_t|^2dt\right],
\end{align*}
where $\hat{Y}_t=Y^1_t-Y^2_t$, $\hat{Z}_t=Z^1_t-Z^2_t$, $\hat{\xi}=\xi^1-\xi^2$ and $$\hat{f}_t=f^1(t,Y^1_t,\P_{Y^1_t},Z^1_t,\P_{Z^1_t})-f^2(t,Y^1_t,\P_{Y^1_t},Z_t^1,\P_{Z_t^1}).$$
 %The solution to BSDE \eqref{linearcase} with double linear mean reflections is unique.
\end{proposition}



\begin{proof}%Suppose that $(Y^i,Z^i,K^i)$ are two solutions of \eqref{linearcase}.
 Let $\hat{K}_t=K^1_t-K^2_t$ and $\widehat{f^2_t}=f^2(t,Y^1_t,\P_{Y^1_t},Z^1_t,\P_{Z^1_t})-f^2(t,Y^2_t,\P_{Y^2_t},Z_t^2,\P_{Z^2_t})$. Applying It\^{o}'s formula to $e^{at}\hat{Y}_t^2$, where $a$ is a positive constant to be determined later, we have
\begin{equation}
\begin{split}\label{e1}
&|\hat{Y}_t|^2 e^{at}+\int_t^T ae^{as}|\hat{Y}_s|^2 ds+\int_t^T e^{as} |\hat{Z}_s|^2 ds\\
=&e^{aT}|\hat{\xi}|^2+\int_t^T 2e^{as}\hat{Y}_s \left(\hat{f}_s+\widehat{f^2_s}\right) ds+\int_t^T 2e^{as}\hat{Y}_s d\hat{K}_s-2\int_t^T e^{as}\hat{Y}_s\hat{Z}_sdB_s.
\end{split}
\end{equation}
By the assumption on $f^2$, the H\"{o}lder inequality and the property of 1-Wasserstein distance, we obtain that
\begin{equation}
\begin{split}\label{e2}
\int_t^T 2e^{as}\hat{Y}_s \widehat{f^2_s} ds\leq &\int_t^T 2\lambda e^{as}\left(|\hat{Y}_s|^2+|\hat{Y}_s\hat{Z}_s|+|\hat{Y}_s|W_1\left(\P_{Y^1_s},\P_{Y^2_s}\right)+|\hat{Y}_s|W_1\left(\P_{Z^1_s},\P_{Z^2_s}\right)\right) ds\\
\leq &\int_t^T 2\lambda e^{as}\left(|\hat{Y}_s|^2+|\hat{Y}_s\hat{Z}_s|+|\hat{Y}_s|\E[|\hat{Y}_s|]+|\hat{Y}_s|\E[|\hat{Z}_s|]\right) ds\\
\leq& \int_t^T (3\lambda+8\lambda^2)e^{as}|\hat{Y}_s|^2 ds+\int_t^T \lambda e^{as}\E[|\hat{Y}_s|^2]ds\\
&+\frac{1}{4}\int_t^T e^{as}|\hat{Z}_s|^2 ds+\frac{1}{4}\int_t^T e^{as}\E[|\hat{Z}_s|^2] ds,\\
\int_t^T 2e^{as}\hat{Y}_s \hat{f}_s ds\leq &\int_t^T e^{as}|\hat{Y}_s|^2 ds+\int_t^T e^{as}|\hat{f}_s|^2 ds.
\end{split}
\end{equation}
Set $a=\frac{3}{2}+4\lambda+8\lambda^2$. Plugging \eqref{e2} into \eqref{e1} and then taking expectations on both sides yield that 
\begin{align}\label{e3'}
\frac{1}{2}\E\left[\int_t^T e^{as}|\hat{Y}_s|^2 ds+\int_t^T e^{as} |\hat{Z}_s|^2 ds\right]\leq \E\left[e^{aT}|\hat{\xi}|^2+\int_t^T e^{as}|\hat{f}_s|^2 ds+ \int_t^T 2e^{as}\hat{Y}_s d\hat{K}_s\right].
\end{align}
It follows from the flat-off condition that 
\begin{equation}\label{e3}
\begin{split}
\E\left[\int_t^T e^{as}\hat{Y}_s d\hat{K}_s\right]&=\int_t^T e^{as}\big((\E[Y_s^1]-l_s)-(\E[Y_s^2]-l_s)\big)dK^{1,l}_s\\
&+\int_t^T e^{as}\big((r_s-\E[Y_s^1])-(r_s-\E[Y_s^2])\big)dK^{1,r}_s\\
&-\int_t^T e^{as}\big((\E[Y_s^1]-l_s)-(\E[Y_s^2]-l_s)\big)dK^{2,l}_s\\
&-\int_t^T e^{as}\big((r_s-\E[Y_s^1])-(r_s-\E[Y_s^2])\big)dK^{2,r}_s\leq 0.
\end{split}
\end{equation}
Combining Eqs. \eqref{e3'} and \eqref{e3}, we get the desired result.
\end{proof}

\begin{remark}
Note that the proof of Proposition \ref{uniquenesslinear} do not need $r,l$ are absolutely continuous. If $r,l\in C[0,T]$ satisfy $l_T\leq \E[\xi]\leq r_T$, the result of Proposition \ref{uniquenesslinear} still holds.
\end{remark}


In the following of this subsection, $C$ will always represent a positive constant  depending on $T,\lambda$, which may vary from line to line. %We first establish the estimates for $Y^n$ and $Z^n$ uniformly in $n$. 

\begin{proposition}\label{estimateYnZn}
There exists a constant $C$ independent of $n$, such that
\begin{align*}
&\sup_{t\in[0,T]}\E[\left|Y_t^n\right|^2]\leq C\left(\int_0^T a_s^2ds+\E\left[\xi^2\right]+\E\left[\int_0^T |f(s,0,\delta_0,0,\delta_0)|^2 ds\right]\right),\\
&\E\left[\int_0^T \left|Z_s^n\right|^2 ds\right]\leq C\left(\int_0^T a_s^2ds+\E\left[\xi^2\right]+\E\left[\int_0^T \left|f(s,0,\delta_0,0,\delta_0)\right|^2 ds\right]\right).
\end{align*}
\end{proposition}

\begin{proof}
Set $\widetilde{Y}^n_t=Y^n_t-r_t$. Applying It\^{o}'s formula to $e^{\beta t}|\widetilde{Y}^n_t|^2$, where $\beta$ is a positive constant to be determined later, we have
\begin{equation}\label{eq1.61}\begin{split}
&e^{\beta t}|\widetilde{Y}^n_t|^2+\int_t^T \beta e^{\beta s}|\widetilde{Y}^n_s|^2 ds+\int_t^T e^{\beta s}\left|Z_s^n\right|^2 ds\\
=&e^{\beta T}|\xi-r_T|^2+\int_t^T 2e^{\beta s}\widetilde{Y}^n_s\left(f(s,Y_s^n,\P_{Y^n_s},Z_s^n,\P_{Z^n_s})+a_s\right)ds-\int_t^T 2e^{\beta s}\widetilde{Y}_s^n Z_s^ndB_s\\
&+\int_t^T 2ne^{\beta s}\widetilde{Y}_s^n\left(\E[Y_s^n]-l_s\right)^-ds-\int_t^T 2ne^{\beta s}\widetilde{Y}_s^n(\E[\widetilde{Y}_s^n])^+ ds.
\end{split}\end{equation}
Simple calculation yields that 
\begin{equation}\label{eq1.62}\begin{split}
&2\widetilde{Y}^n_s\left(f(s,Y_s^n,\P_{Y^n_s},Z_s^n,\P_{Z^n_s})+a_s\right)\\
\leq &|f(s,r_s,\delta_{r_s},0,\delta_0)|^2+|a_s|^2+\frac{1}{4}|Z^n_s|^2+\frac{1}{4}W_1^2(\P_{Z^n_s},\delta_{0})+\lambda W_1^2(\P_{{Y}^n_s},\delta_{r_s})+(2+3\lambda+8\lambda^2)|\widetilde{Y}^n_s|^2\\
\leq &|a_s|^2+2|f(s,0,\delta_0,0,\delta_0)|^2+4\lambda^2|r_s|^2+\frac{1}{4}|Z^n_s|^2+\frac{1}{4}\E[|Z^n_s|^2]+\lambda \E[|\widetilde{Y}^n_s|^2]+(2+3\lambda+8\lambda^2)|\widetilde{Y}^n_s|^2.
\end{split}\end{equation}
Set $\beta=3+4\lambda+8\lambda^2$. Plugging Eq. \eqref{eq1.62} to Eq. \eqref{eq1.61} and taking expectations on both sides, noting that \begin{displaymath}
\E[\widetilde{Y}^n_s](\E[Y_s^n]-l_s)^-\leq 0,\ \E[\widetilde{Y}^n_s](\E[\widetilde{Y}_s^n])^+\geq 0,
\end{displaymath}
 we finally derive that
\begin{align*}
\E\left[|\widetilde{Y}_t^n|^2+\int_t^T |Z^n_s|^2 ds\right]\leq &C\E\left[|\xi-r_T|^2+\int_t^T |f(s,0,\delta_0,0,\delta_0)|^2ds+\int_t^T r_s^2ds+\int_t^T a_s^2ds\right]\\
\leq &C\E\left[|\xi|^2+\int_0^T |f(s,0,\delta_0,0,\delta_0)|^2ds+\int_0^T a_s^2ds\right].
\end{align*}
Recalling the definition of $\widetilde{Y}^n_t$, we obtain the desired result.
\end{proof}

In the following, we show that $K^{n,l}_T$ and $K^{n,r}_T$ are uniformly bounded.
\begin{proposition}\label{estimateYn-r}
There exists a constant $C$ independent of $n$, such that
\begin{align*}
n^2\int_0^T\left|(\E[Y_t^n]-r_t)^+\right|^2 dt\leq {C},\
n^2\int_0^T\left|(\E[Y_t^n]-l_t)^-\right|^2 dt\leq {C}.
\end{align*}
\end{proposition}

\begin{proof}
We only prove the first inequality since the second one can be proved in a same manner. We define $y^n_t:=\E[Y^n_t]$. Taking expectations on both sides of Eq. \eqref{panelization}, we have
\begin{align*}
y_t^n=\E[\xi]+\int_t^T \E\left[f\left(s,Y^n_s,\P_{Y^n_s},Z^n_s,\P_{Z^n_s}\right)\right]ds+\int_t^T n(y^n_s-l_s)^- ds-\int_t^T n(y^n_s-r_s)^+ds.
\end{align*} 
It is easy to check that 
\begin{align*}
d\left|(y^n_t-r_t)^+\right|^2=-2(y^n_t-r_t)^+ \left[ \E\left[f(t,Y^n_t,\P_{Y^n_t},Z^n_t,\P_{Z^n_t})\right]+a_t+n(y^n_t-l_t)^- -n(y^n_t-r_t)^+\right]dt.
\end{align*}
Noting that $y^n_T-r_T\leq 0$ and $(y^n_t-r_t)^+(y^n_t-l_t)^-\leq 0$, we have
\begin{align*}
&|(y^n_0-r_0)^+|^2+2n\int_0^T |(y^n_t-r_t)^+|^2dt\\
=&2\int_0^T (y^n_t-r_t)^+\left(\E\left[f(t,Y^n_t,\P_{Y^n_t},Z^n_t,\P_{Z^n_t})\right]+a_t\right) dt+2n\int_0^T (y^n_t-r_t)^+(y^n_t-l_t)^-dt\\
\leq &n\int_0^T  \left|(y^n_t-r_t)^+\right|^2dt+\frac{1}{n}\int_0^T \left(\E[f(t,Y^n_t,\P_{Y^n_t},Z^n_t,\P_{Z^n_t})]+a_t\right) ^2dt\\
\leq &n\int_0^T  \left|(y^n_t-r_t)^+\right|^2dt+\frac{C}{n}\int_0^T \left(\E[\left|f(t,0,\delta_0,0,\delta_0)\right|^2]+\E[\left|Y_t^n\right|^2]+\E[\left|Z_t^n\right|^2]+a^2_t\right)dt.
\end{align*}
Combining the above inequality and Proposition \ref{estimateYnZn} yields the desired result.
\end{proof}

Applying Proposition \ref{estimateYnZn}, \ref{estimateYn-r} and the B-D-G inequality, we obtain the following uniform estimate for $Y^n$.
\begin{proposition}\label{estimateKn}
There exists a constant $C$ independent of $n$, such that
\begin{align*}
\E\left[\sup_{t\in[0,T]}\left|Y_t^n\right|^2\right]\leq C.
\end{align*}
\end{proposition}

%\begin{proof}
%Proposition \ref{estimateYn-r} indicates that $|K^r_T|\leq C$ and $|K^l_T|\leq C$. Applying the B-D-G inequality implies the uniform estimate for $Y^n$.
%\end{proof}



\begin{proof}[Proof of Theorem \ref{existence and uniqueness}]
Uniqueness is a direct consequence of Theorem \ref{th-1} with $p=2$. It remains to show the existence. We prove that there exists a triple of processes $(Y,Z,K)\in\mathcal{S}^2\times \mathcal{H}^2\times BV[0,T]$ such that
\begin{align}\label{equation6.25}
\E\left[\sup_{t\in[0,T]}\left|Y_t^n-Y_t\right|^2\right]+\E\left[\int_0^T \left|Z_s^n-Z_s\right|^2 ds\right]+\sup_{t\in[0,T]}\left|K^n_t-K_t\right|\rightarrow 0, \textrm{ as } n\rightarrow \infty.
\end{align} 
Besides, $(Y,Z,K)$ is the solution to \eqref{linearcase}.

For this purpose, for any positive integer $n$,  set $\hat{X}:=X^{n+1}-X^n$ for $X=Y,Z,K$. By a similar analysis as Eqs. \eqref{e1} and \eqref{e2}, choosing $a=\frac{1}{2}+4(\lambda+2\lambda^2)$, we have
\begin{align*}
&e^{at}|\hat{Y}_t|^2+\left(\lambda+\frac{1}{2}\right)\int_t^T e^{as} |\hat{Y}_s|^2ds +\frac{3}{4}\int_t^T e^{as}|\hat{Z}_s|^2ds\\
\leq &\lambda\int_t^T e^{as} \E[|\hat{Y}_s|^2]ds +\frac{1}{4}\int_t^T e^{as}\E[|\hat{Z}_s|^2]ds+2\int_t^T e^{as}\hat{Y}_sd\hat{K}_s-2\int_t^T e^{as}\hat{Y}_s\hat{Z}_sdB_s,
\end{align*}
Taking expectations on both sides yields that 
\begin{equation}\label{e30}
\sup_{t\in[0,T]}\E[|\hat{Y}_t|^2]+\E\left[\int_0^T (|\hat{Y}_s|^2+|\hat{Z}_s|^2)ds\right]\leq C\sup_{t\in[0,T]}\E\left[\int_t^T e^{as}\hat{Y}_sd\hat{K}_s\right].
\end{equation}

It is easy to check that for any $x,y\in\mathbb{R}$, we have
\begin{align*}
-nx^2+(2n+1)xy-(n+1)y^2=-n(x-\frac{2n+1}{2n}y)^2+\frac{y^2}{4n}.
\end{align*}
Set $y^n_t=\E[Y^n_t]$, $v^n_t=(y^n_t-l_t)^-$ and $u^n_t=(y_t^n-r_t)^+$. Then, we may calculate that
\begin{align*}%\label{eq1.82}\begin{split}
&[(y^{n+1}_s-l_s)-(y^n_s-l_s)]((n+1)v_s^{n+1}-nv_s^n)\\
\leq&-n|v_s^n|^2+(2n+1)v^n_s v^{n+1}_s-(n+1)|v_s^{n+1}|^2\leq \frac{|v^{n+1}_s|^2}{4n}
\end{align*}%\end{equation}
and
\begin{align*}%\label{eq1.83}\begin{split}
&-[(y^{n+1}_s-r_s)-(y^n_s-r_s)]((n+1)u_s^{n+1}-nu_s^n)\\
\leq&-n|u_s^n|^2+(2n+1)u^n_s u^{n+1}_s-(n+1)|u_s^{n+1}|^2\leq \frac{|u^{n+1}_s|^2}{4n}.
\end{align*}%\end{equation}
Then, we obtain that
\begin{equation}\label{eq1.81}\begin{split}
\E\left[\int_t^T e^{as}\hat{Y}_sd\hat{K}_s\right]=&\int_t^T e^{as}(y^{n+1}_s-y_s^n)\bigg(\big[(n+1)v_s^{n+1}-nv_s^n\big]-\big[(n+1)u_s^{n+1}-nu_s^n\big]\bigg)ds\\
 =&\int_t^T e^{as}\left[(y^{n+1}_s-l_s)-(y^n_s-l_s)\right]((n+1)v_s^{n+1}-nv_s^n)ds\\
 &-\int_t^T e^{as}\left[(y^{n+1}_s-r_s)-(y^n_s-r_s)\right]((n+1)u_s^{n+1}-nu_s^n)ds\\
 \leq &\frac{C}{n}\int_0^T \left(|v^{n+1}_s|^2+|u^{n+1}_s|^2\right)ds\leq \frac{C}{n^3},
\end{split}\end{equation}
where we have used Proposition \ref{estimateYn-r} in the last inequality. 
%Combining Eq. \eqref{eq1.81}-\eqref{eq1.83}, by Proposition \ref{estimateYn-r}, we have
%\begin{align*}
%\E[\int_t^T e^{as}\hat{Y}_sd\hat{K}_s]\leq \frac{C}{n}\int_0^T (|v^{n+1}_s|^2+|u^{n+1}_s|^2)ds\leq \frac{C}{n^3}.
%\end{align*}
Plugging this estimate into \eqref{e30} indicates that
\begin{align}\label{eq1.84}
\sup_{t\in[0,T]}\E[|\hat{Y}_t|^2]+\E\left[\int_0^T \left(|\hat{Y}_s|^2+|\hat{Z}_s|^2\right)ds\right]\leq \frac{C}{n^3}.
\end{align}
Since $K^n$ is deterministic, it follows from \eqref{panelization} that
\begin{align*}
\hat{K}_T-\hat{K_t}=\E[\hat{Y}_t]-\E\left[\int_t^T \left(f\left(s,Y^{n+1}_s,\P_{Y^{n+1}_s},Z^{n+1}_s,\P_{Z^{n+1}_s}\right)-f\left(s,Y_s^n,\P_{Y^n_s},Z_s^n,\P_{Z^n_s}\right)\right)ds\right]
\end{align*}
and 
\begin{align}\label{eq1.84'}
\hat{Y}_t=\E_t\left[\int_t^T \left(f\left(s,Y^{n+1}_s,\P_{Y^{n+1}_s},Z^{n+1}_s,\P_{Z^{n+1}_s}\right)-f\left(s,Y_s^n,\P_{Y^n_s},Z_s^n,\P_{Z^n_s}\right)\right)ds\right]+\hat{K}_T-\hat{K_t}.
\end{align}
Using the H\"{o}lder inequality, the property of 1-Wasserstein distance and the estimate \eqref{eq1.84}, we obtain that% implies that
\begin{align*}
\sup_{t\in[0,T]}|\hat{K}_T-\hat{K_t}|^2\leq \frac{C}{n^3}.
\end{align*}
Noting the representation for $\hat{Y}$ in \eqref{eq1.84'}, by the B-D-G inequality, we have
\begin{align*}
\E\left[\sup_{t\in[0,T]}|\hat{Y}_t|^2\right]\leq \frac{C}{n^3}.
\end{align*} 
All the above analysis indicates that \eqref{equation6.25} holds.

It remains to prove the limit processes $(Y,Z,K)$ is indeed the solution to \eqref{linearcase}. The proof is similar with the proof of Theorem 5.5 in \cite{L}, so we omit it. The proof is complete.
\end{proof}

\begin{remark}
Compared with the penalty method for the reflected MFBSDE studied in \cite{CHM,DEH}, the generator $f$ in our case may depend on the distribution of both the solution $Y$ and $Z$.  Besides, we do not need any monotonicity condition for the generator $f$. %, i.e., we do not need to restrict ourselves to the case that the driver can o
%(i) Compared with the condition proposed for the reflecting boundaries in Theorem \ref{main}, we do not need to assume that $\inf_{r\in[0,T]}(r_t-l_t)>0$, that is, the upper obstacle and lower obstacle do not need to be completely separated. %\cite{FS}, we give a different approach, i.e., the penalization method which is frequently used in the classical reflected case, to the construction of solution to BSDEs with two linear mean reflecting barriers.
%\noindent (ii) The penalization method in our paper extends the one in \cite{BEH} to the following aspects. First, we could deal with the case of two reflecting barriers. Second, we do not need to assume the benchmark function $l,r$ are constant.
\end{remark}

\iffalse
{\bf Step 2.}  We show that for any $p\geq 1$,
\begin{equation}\label{uniformbounded}
\sup _{m \geq 0} \mathbf{E}\left[\exp \left\{p \gamma \sup _{t \in[t_{0}, t_{0}+h]}\left|Y_t^{(m)}\right|\right\}+\left(\int_{t_{0}}^{t_{0}+h}\left|Z_t^{(m)}\right|^2 d t\right)^p+\left|K_{t_{0}+h}^{(m)}\right|\right]<\infty .
\end{equation}
To this end, we first show that $$\sup_{m\geq 0}\mathbf{E}\left [ \exp\left \{ p\gamma\sup_{t \in[t_{0}, t_{0}+h]}\left | Y_t^{(m)} \right |  \right \}   \right ] <\infty.$$
According to  \eqref{re-2} , we have   
$$
\begin{aligned}
\left|Y^{(m)}_t\right|&=\left|y^{(m)}_t+\bar{K}^{(m)}_{t_0+h-t}-\bar{K}^{(m)}_{0}\right|\\
  &=\left|y^{(m)}_t+(\bar{K}^{(m)}_{t_0+h-t}-\bar{K}^{3}_{t_0+h-t})+(\bar{K}^{3}_{t_0+h-t}-\bar{K}^3_0)+(\bar{K}^{3}_{0}-\bar{K}^{(m)}_{0})\right|\text {, }\\
\end{aligned}
$$
where $\bar{K}^3$ is the second component of the solution to the Skorokhod problem $\mathbb{SP}_{\bar{l}^3}^{\bar{r}^3}(\bar{s}^3)$. Here $\bar{s}^3=\mathbf{E}[\eta]$ and 
\begin{align*}
&\bar{l}^3(t,x)=l^3(t_0+h-t,x)=\mathbf{E}\left[L(t_0+h-t,x)\right], \\
&\bar{r}^3(t,x)=r^3(t_0+h-t,x)=\mathbf{E}\left[R(t_0+h-t,x)\right].
\end{align*}
In view of Proposition 3.4 in $\cite{Li}$, we have
$$
\begin{aligned}
\left|\bar{K}^{(m)}_{t_0+h-t}-\bar{K}^{3}_{t_0+h-t}\right|&\leq\frac{C}{c}\left|\bar{s}^{(m)}_{t_0+h-t}-\bar{s}^{3}_{t_0+h-t}\right|+\frac{1}{c}\left(\left|\bar{l}^{(m)}(t_0+h-t)-\bar{l}^{3}(t_0+h-t)\right|\right.\\
&\left.\vee \left|\bar{r}^{(m)}(t_0+h-t)-\bar{r}^{3}(t_0+h-t)\right|\right),\\%\leq \frac{3C}{c}\sup_{t\in[0,T]}\mathbf{E}\left[\left|y_t^2\right| \right]+\frac{C}{c}\mathbf{E}[\left|\xi\right|].    
\end{aligned}
$$
Recalling \eqref{barsbarlbarr}, we have %that $\bar{s}^{(m)}_{t}=\mathbf{E}[y^{(m)}_{t_0+h-t}],\bar{s}^0=\mathbf{E}[\eta]$, we have 
$$\left|\bar{s}^{(m)}_{t_0+h-t}-\bar{s}^{3}_{t_0+h-t}\right|\leq\left|\mathbf{E}[y^{(m)}_{t}]-\mathbf{E}[\eta]\right|\leq\mathbf{E}\left[\left|y^{(m)}_{t}\right|\right]+\mathbf{E}\left[\left|\eta\right|\right]$$
and
%Recalling that 
%$$\bar{l}^0(t,x)=l^0(t_0+h-t,x)=\mathbf{E}\left[L(t_0+h-t,x)\right]$$  $$\bar{l}^{(m)}(t,x)=l^{(m)} (t_0+h-t,x)=\E\left[L(t_0+h-t,y^{(m)}_{t_0+h-t}-\E[y^{(m)}_{t_0+h-t}]+x)\right]$$ and applying Assumption $\ref{ass-1}$,
%we can obtain that 
$$
\begin{aligned}
  \left|\bar{l}^{(m)}(t_0+h-t)-\bar{l}^{3}(t_0+h-t)\right|%\\=\left|l^{(m)}(t,x)-l^{0}(t,x)\right|\\
  = \left|\mathbf{E}\left[L(t,y^{(m)}_{t}-\E[y^{(m)}_{t}]+x)\right]-\mathbf{E}\left[L(t,x)\right]\right|
 % &\leq C\mathbf{E}\left[\left|y_{t}^{(m)}-\E\left[y_{t}^{(m)}\right]\right|\right]
 \leq 2C\E\left[\left|y_{t}^{(m)}\right|\right].
\end{aligned}
$$
Similar estimate hold for $\left|\bar{r}^{(m)}(t_0+h-t)-\bar{r}^{3}(t_0+h-t)\right|$. The above analysis imply that 
\begin{equation}\label{t_0+h-t}
\left|\bar{K}^{(m)}_{t_0+h-t}-\bar{K}^{3}_{t_0+h-t}\right|\leq 3\frac{C}{c}\mathbf{E}\left[\left|y^{(m)}_{t}\right|\right]+\frac{C}{c}\mathbf{E}\left[\left|\eta\right|\right].    
\end{equation}

For the same reason, we have %{\color{red}some mistakes in the following equation}
\begin{equation}\label{0}
 \left|\bar{K}^{(m)}_{0}-\bar{K}^{3}_{0}\right|\leq 3\frac{C}{c}\mathbf{E}\left[\left|y^{(m)}_{t_0+h}\right|\right]+\frac{C}{c}\mathbf{E}\left[\left|\eta\right|\right].   
\end{equation}

Therefore,$$
\begin{aligned}
\left|Y^{(m)}_t\right|&=\left|y^{(m)}_t+\bar{K}^{(m)}_{t_0+h-t}-\bar{K}^{(m)}_{0}\right|\\
&\leq \left|y^{(m)}_t\right|+3\frac{C}{c}\mathbf{E}\left[\left|y^{(m)}_t\right|\right]+2\frac{C}{c}\mathbf{E}\left[\left|\eta\right|\right]+3\frac{C}{c}\mathbf{E}\left[\left|y^{(m)}_{t_0+h}\right|\right]+\left|\bar{K}^3_{t_0+h-t}-\bar{K}^3_0\right|,\\
\end{aligned}
$$
Consequently, we have
%and for any $t\in [t_0, t_0+h]$, 
\begin{equation}\label{supY}
    \sup_{t\in[t_0, t_0+h]}\left|Y^{(m)}_t\right|\leq\sup_{t\in[t_0, t_0+h]}\left|y^{(m)}_t\right|+6\frac{C}{c}\sup_{t\in[t_0, t_0+h]}\mathbf{E}\left[\left|y^{(m)}_t\right|\right]+B_2
\end{equation}
where $B_2:=2\frac{C}{c}\mathbf{E}\left[\left|\eta\right|\right]+2\sup_{t\in[t_0, t_0+h]}\left|\bar{K}^3_t\right|$. Similar to Eq. (24) and Eq. (25) in \cite{HMW},  we have 
$$\exp\left\{\gamma\left|y_{t}^{(m)}\right|\right\}\leq \mathbf{E}_{t}\left[\exp\left\{\gamma\left(\widetilde{\eta}+\lambda h\left(\sup_{t\in[t_0, t_0+h]}\left|Y^{(m-1)}_t\right|+\sup_{t\in[t_0, t_0+h]}\mathbf{E}\left[\left|Y^{(m-1)}_t\right|\right]\right)\right)\right\}\right]$$
and for any $m\geq 1, p\geq 1$, $t\in[t_0,t_0+h]$,
\begin{equation}\label{yineq}
\begin{split}
&\mathbf{E}\left[\exp\left\{p\gamma\sup_{t\in[t_0,t_0+h]}\left|y_{t}^{(m)}\right|\right\}\right]\\
\leq& 4\mathbf{E}\left[\exp\left\{p\gamma\left(\widetilde{\eta} +\lambda h\left(\sup_{t\in[t_0, t_0+h]}\left|Y^{(m-1)}_t\right|+\sup_{t\in[t_0, t_0+h]}\mathbf{E}\left[\left|Y^{(m-1)}_t\right|\right]\right)\right)\right\}\right] \\   
\leq& 4\mathbf{E}\left[\exp\left\{p\gamma\left(\widetilde{\eta}+\lambda h\sup_{t\in[t_0, t_0+h]}\left|Y^{(m-1)}_t\right|\right)\right\}\right]\mathbf{E}\left[\exp\left\{p\gamma\lambda h\sup_{t\in[t_0,t_0+h]}\left|Y^{(m-1)}_t\right|\right\}\right]\\
\leq& 4\mathbf{E}\left[\exp\left\{2p\gamma\widetilde{\eta}\right\}\right]^{\frac{1}{2}}\mathbf{E}\left[\exp\left\{2p \gamma \lambda h \sup_{t\in[t_0,t_0+h]}\left|Y^{(m-1)}_t\right|\right\}\right],\\
\end{split}  
\end{equation}
where $\widetilde{\eta}=|\eta |+\int_{t_0}^{t_0+h} \alpha_s ds $. Applying Lemma $\ref{le-2}$,  we obtain that
$$
\begin{aligned}
	 \mathbf{E}\left[\exp \left\{p \gamma \sup _{t\in[t_0,t_0+h]}\left|Y_t^{(m)}\right|\right\}\right] &\leq \widetilde{C}_2 \mathbf{E}\left[\exp \left\{\left(4+\frac{24C}{c}\right) p \gamma \lambda h \sup _{t\in[t_0,t_0+h]}\left|Y_t^{(m-1)}\right|\right\}\right] \\
 &\leq \widetilde{C}_2 \mathbf{E}\left[\exp \left\{ p \gamma  \sup _{t\in[t_0,t_0+h]}\left|Y_t^{(m-1)}\right|\right\}\right]^{\left(4+\frac{24C}{c}\right)\lambda h},
\end{aligned}
$$
where $$\widetilde{C}_2=4 \exp\left\{p \gamma B_2\right\} \mathbf{E}\left[\exp\left\{\left(4+\frac{24C}{c}\right)p\gamma\widetilde{\eta}\right\}\right]^{\frac{1}{2}}.$$
We set $\rho=\frac{1}{1-\left(4+24\frac{C}{c}\right)\beta h}$.  Then iterating the above procedure $m$ times yields that
$$\mathbf{E}\left[\exp \left\{p \gamma \sup _{t\in[t_0,t_0+h]}\left|Y_t^{(m)}\right|\right\}\right]\leq
\widetilde{C}_2^{\rho} \mathbf{E}\left[\exp\left\{24\frac{C}{c}p\gamma\left(\left|\eta\right|+\int_{t_0}^{t_0+h}\alpha_s ds\right)\right\}\right]^{\frac{\rho}{2}}<\infty.$$
It then follows that $$\sup_{m\geq 0}\mathbf{E}\left [ \exp\left \{ p\gamma\sup_{t\in[t_0,t_0+h]}\left | Y_t^{(m)} \right |  \right \}   \right ] <\infty.$$

Second, according to  inequality $(\ref{yineq})$, we have
$$\sup _{m \geq 0} \mathbf{E}\left[\exp \left\{p \gamma \sup _{t\in[t_0,t_0+h]}\left|y_t^{(m)}\right|\right\}\right]<\infty, \quad \forall p \geq 1.$$
Recalling that \eqref{t_0+h-t} and \eqref{0},  we have %{\color{red}it is better to recall where we can get the following estimate}
$$
\begin{aligned}
\sup _{m \geq 0} \left|K_{t_0+h}^{(m)}\right| 
&=\sup _{m \geq 0} \left|\bar{K}_{h}^{(m)}-\bar{K}_{0}^{(m)}\right|\\
&=\sup _{m \geq 0} \left|\bar{K}_{h}^{(m)}-\bar{K}_{h}^{3}+\bar{K}_{h}^{3}-\bar{K}_{0}^{3}+\bar{K}_{0}^{3}-\bar{K}_{0}^{(m)}\right|\\
&\leq 6\frac{C}{c}\sup_{m \geq 0}\mathbf{E}\left[\sup _{t\in[t_0,t_0+h]}\left|y_t^{(m)}\right|\right]+2\frac{C}{c}\mathbf{E}\left[\left|\eta\right|\right]+2\sup_{t\in[0,T]}\left|\bar{K}^3_t\right|<\infty \\
\end{aligned}
$$
Finally, since $Z^{(m)}=z^{(m)}$ (see Eq. \eqref{Zmzm}), applying  Corollary 4 in $\cite{BH08}$ to the quadratic $\text{BSDE}$ \eqref{ytm}, we obtain 
$$
\sup _{m \geq 0} \mathbf{E}\left[\left(\int_{t_0}^{t_0+h}\left|Z_t^{(m)}\right|^2 d t\right)^p\right]<\infty, \forall p \geq 1.
$$
\fi

\renewcommand\thesection{Appendix B}
\section{Supplementary proofs}\label{append:B}
\renewcommand\thesection{B}

\begin{proof}[Proof of Eq. \eqref{uniformbounded}]
We first show that $$\sup_{m\geq 0}\mathbf{E}\left [ \exp\left \{ p\gamma\sup_{t \in[t_{0}, t_{0}+h]}\left | Y_t^{(m)} \right |  \right \}   \right ] <\infty.$$
According to  \eqref{re-2} , we have   
$$
\begin{aligned}
\sup_{t\in[t_0,t_0+h]}\left|Y^{(m)}_t\right|&=\sup_{t\in[t_0,t_0+h]}\left|y^{(m)}_t+\bar{K}^{(m)}_{t_0+h-t}\right|\\
  &=\sup_{t\in[t_0,t_0+h]}\left|y^{(m)}_t+(\bar{K}^{(m)}_{t_0+h-t}-\bar{K}^{3}_{t_0+h-t})+\bar{K}^{3}_{t_0+h-t}\right|\\
  &\leq\sup_{t\in[t_0,t_0+h]}\left|y^{(m)}_t\right|+\sup_{t\in[t_0,t_0+h]}\left|\bar{K}^{(m)}_{t_0+h-t}-\bar{K}^{3}_{t_0+h-t}\right|+\sup_{t\in[t_0,t_0+h]}\left|\bar{K}^{3}_{t_0+h-t}\right|\text {,}\\
\end{aligned}
$$
where $\bar{K}^3$ is the second component of the solution to the Skorokhod problem $\mathbb{SP}_{\bar{l}^3}^{\bar{r}^3}(\bar{s}^3)$. Here $\bar{s}^3=\mathbf{E}[\eta]$ and 
\begin{align*}
&\bar{l}^3(t,x)=l^3(t_0+h-t,x)=\mathbf{E}\left[L(t_0+h-t,x)\right], \\
&\bar{r}^3(t,x)=r^3(t_0+h-t,x)=\mathbf{E}\left[R(t_0+h-t,x)\right].
\end{align*}

In view of Proposition $\ref{continuity}$ , we have
$$
\begin{aligned}
\sup_{t\in[t_0,t_0+h]}\left|\bar{K}^{(m)}_{t_0+h-t}-\bar{K}^{3}_{t_0+h-t}\right|&=\sup_{t\in[0,h]}\left|\bar{K}^{(m)}_{t}-\bar{K}^{3}_{t}\right|\\
&\leq\frac{C}{c}\sup_{t\in[0,h]}\left|\bar{s}^{(m)}_{t}-\bar{s}^{3}_{t}\right|+\frac{1}{c}\left(\bar{L}^{3}_h\vee\bar{R}^{3}_h\right)\text{,}\\     
\end{aligned}
$$
where $\bar{L}^3_h=\sup_{(t,x)\in[0,h]\times\mathbb{R}}\left|\bar{l}^{(m)}(t,x)-\bar{l}^{3}(t,x)\right|$ and
$\bar{R}^3_h=\sup_{(t,x)\in[0,h]\times\mathbb{R}}\left|\bar{r}^{(m)}(t,x)-\bar{r}^{3}(t,x)\right|$.\\

Recalling \eqref{barsbarlbarr}, we have %that $\bar{s}^{(m)}_{t}=\mathbf{E}[y^{(m)}_{t_0+h-t}],\bar{s}^0=\mathbf{E}[\eta]$, we have 
$$\sup_{t\in [0,h]}\left|\bar{s}^{(m)}_{t}-\bar{s}^{3}_{t}\right|\leq\sup_{t\in [0,h]}\left|\mathbf{E}[y^{(m)}_{t_0+h-t}]-\mathbf{E}[\eta]\right|\leq\sup_{t\in [t_0,t_0+h]}\mathbf{E}[|y^{(m)}_{t}|]+\mathbf{E}\left[\left|\eta\right|\right]$$
and
%Recalling that 
%$$\bar{l}^0(t,x)=l^0(t_0+h-t,x)=\mathbf{E}\left[L(t_0+h-t,x)\right]$$  $$\bar{l}^{(m)}(t,x)=l^{(m)} (t_0+h-t,x)=\E\left[L(t_0+h-t,y^{(m)}_{t_0+h-t}-\E[y^{(m)}_{t_0+h-t}]+x)\right]$$ and applying Assumption $\ref{ass-1}$,
%we can obtain that 
$$
\begin{aligned}
\bar{L}_h^3&=\sup_{(t,x)\in[0,h]\times\mathbb{R}}\left|\bar{l}^{(m)}(t,x)-\bar{l}^{3}(t,x)\right|\\
  &=\sup_{(t,x)\in[0,h]\times\mathbb{R}} \left|\mathbf{E}\left[L(t_0+h-t,y^{(m)}_{t_0+h-t}-\E[y^{(m)}_{t_0+h-t}]+x)\right]-\mathbf{E}\left[L(t_0+h-t,x)\right]\right|\\
 &\leq 2C\sup_{t\in [0,h]}\E[|y_{t_0+h-t}^{(m)}|]=2C\sup_{t\in [t_0,t_0+h]}\mathbf{E}[|y^{(m)}_{t}|].\\
\end{aligned}
$$
Similar estimate holds for $\bar{R}^{3}_h$. The above analysis imply that 
\begin{equation}\label{t_0+h-t}
\sup_{t\in[t_0,t_0+h]}\left|\bar{K}^{(m)}_{t_0+h-t}-\bar{K}^{3}_{t_0+h-t}\right|\leq 3\frac{C}{c}\sup_{t\in [t_0,t_0+h]}\mathbf{E}[|y^{(m)}_{t}|]+\frac{C}{c}\mathbf{E}\left[\left|\eta\right|\right].    
\end{equation}

Consequently, we have
\begin{equation}\label{supY}
    \sup_{t\in[t_0, t_0+h]}|Y^{(m)}_t|\leq\sup_{t\in[t_0, t_0+h]}|y^{(m)}_t|+3\frac{C}{c}\sup_{t\in[t_0, t_0+h]}\mathbf{E}[|y^{(m)}_t|]+B_2
\end{equation}
where $B_2:=\frac{C}{c}\mathbf{E}\left[\left|\eta\right|\right]+\sup_{t\in[0, T]}\left|\bar{K}^3_t\right|$. Similar to Eq. (24) and Eq. (25) in \cite{HMW},  we have 
$$\exp\left\{\gamma|y_{t}^{(m)}|\right\}\leq \mathbf{E}_{t}\left[\exp\left\{\gamma\left(\widetilde{\eta}+\lambda h\left(\sup_{t\in[t_0, t_0+h]}\left|Y^{(m-1)}_t\right|+\sup_{t\in[t_0, t_0+h]}\mathbf{E}\left[\left|Y^{(m-1)}_t\right|\right]\right)\right)\right\}\right]$$
and for any $m\geq 1, p > 1$, $t\in[t_0,t_0+h]$,
\begin{equation}\label{yineq}
\begin{split}
&\mathbf{E}\left[\exp\left\{p\gamma\sup_{t\in[t_0,t_0+h]}|y_{t}^{(m)}|\right\}\right]\\
\leq& 4\mathbf{E}\left[\exp\left\{p\gamma\left(\widetilde{\eta} +\lambda h\left(\sup_{t\in[t_0, t_0+h]}\left|Y^{(m-1)}_t\right|+\sup_{t\in[t_0, t_0+h]}\mathbf{E}\left[\left|Y^{(m-1)}_t\right|\right]\right)\right)\right\}\right] \\   
\leq& 4\mathbf{E}\left[\exp\left\{p\gamma\left(\widetilde{\eta}+\lambda h\sup_{t\in[t_0, t_0+h]}\left|Y^{(m-1)}_t\right|\right)\right\}\right]\mathbf{E}\left[\exp\left\{p\gamma\lambda h\sup_{t\in[t_0,t_0+h]}\left|Y^{(m-1)}_t\right|\right\}\right]\\
\leq& 4\mathbf{E}\left[\exp\left\{2p\gamma\widetilde{\eta}\right\}\right]^{\frac{1}{2}}\mathbf{E}\left[\exp\left\{2p \gamma \lambda h \sup_{t\in[t_0,t_0+h]}\left|Y^{(m-1)}_t\right|\right\}\right],\\
\end{split}  
\end{equation}
where $\widetilde{\eta}=|\eta |+\int_{t_0}^{t_0+h} \alpha_s ds $. Applying Lemma $\ref{le-2}$,  we obtain that
$$
\begin{aligned}
	 \mathbf{E}\left[\exp \left\{p \gamma \sup _{t\in[t_0,t_0+h]}\left|Y_t^{(m)}\right|\right\}\right] &\leq \widetilde{C}_2 \mathbf{E}\left[\exp \left\{\left(4+\frac{24C}{c}\right) p \gamma \lambda h \sup _{t\in[t_0,t_0+h]}\left|Y_t^{(m-1)}\right|\right\}\right] \\
 &\leq \widetilde{C}_2 \mathbf{E}\left[\exp \left\{ p \gamma  \sup _{t\in[t_0,t_0+h]}\left|Y_t^{(m-1)}\right|\right\}\right]^{\left(4+\frac{24C}{c}\right)\lambda h},
\end{aligned}
$$
where $$\widetilde{C}_2=4 \exp\left\{p \gamma B_2\right\} \mathbf{E}\left[\exp\left\{\left(4+\frac{24C}{c}\right)p\gamma\widetilde{\eta}\right\}\right]^{\frac{1}{2}}.$$
We set $\rho=\frac{1}{1-\left(4+\frac{24C}{c}\right)\beta h}$.  Then iterating the above procedure $m$ times yields that
$$\mathbf{E}\left[\exp \left\{p \gamma \sup _{t\in[t_0,t_0+h]}\left|Y_t^{(m)}\right|\right\}\right]\leq
\widetilde{C}_2^{\rho} \mathbf{E}\left[\exp\left\{24\frac{C}{c}p\gamma\left(\left|\eta\right|+\int_{t_0}^{t_0+h}\alpha_s ds\right)\right\}\right]^{\frac{\rho}{2}}<\infty.$$
It then follows that $$\sup_{m\geq 0}\mathbf{E}\left [ \exp\left \{ p\gamma\sup_{t\in[t_0,t_0+h]}\left | Y_t^{(m)} \right |  \right \}   \right ] <\infty.$$

Second, according to  inequality $(\ref{yineq})$, we have
$$\sup _{m \geq 0} \mathbf{E}\left[\exp \left\{p \gamma \sup _{t\in[t_0,t_0+h]}|y_t^{(m)}|\right\}\right]<\infty, \quad \forall p > 1.$$
Similar to \eqref{t_0+h-t} ,  we have %{\color{red}it is better to recall where we can get the following estimate}
$$
\begin{aligned}
\sup _{m \geq 0} \left|K_{t_0+h}^{(m)}\right| 
&=\sup _{m \geq 0} \left|\bar{K}_{h}^{(m)}-\bar{K}_{0}^{(m)}\right|\\
&=\sup _{m \geq 0} \left|\bar{K}_{h}^{(m)}-\bar{K}_{h}^{3}+\bar{K}_{h}^{3}\right|\\
&\leq 3\frac{C}{c}\sup_{m \geq 0}\mathbf{E}[\sup _{t\in[t_0,t_0+h]}|y_t^{(m)}|]+\frac{C}{c}\mathbf{E}\left[\left|\eta\right|\right]+\sup_{t\in[0,T]}\left|\bar{K}^3_t\right|<\infty \\
\end{aligned}
$$
Finally, since $Z^{(m)}=z^{(m)}$ (see Eq. \eqref{Zmzm}), applying  Corollary 4 in $\cite{BH08}$ to the quadratic $\text{BSDE}$ \eqref{ytm}, we obtain 
$$
\sup _{m \geq 0} \mathbf{E}\left[\left(\int_{t_0}^{t_0+h}\left|Z_t^{(m)}\right|^2 d t\right)^p\right]<\infty, \forall p > 1.
$$  
The proof is complete.
\end{proof}

%\renewcommand\thesection{Appendix C}
%\section{The proof of step 3 in existence of Theorem 4.4 }\label{append:C}
%\renewcommand\thesection{C}
\begin{proof}[Proof of Eq. \eqref{pip}]
 Actually, the proof is similar to the one for Eq. \eqref{bddexpdetltay}. In the following, we mainly highlight the differences. We first show that for any $m,q \geq 1$ , $p>1$ and $\theta\in\left(0,1\right)$
$$
\begin{aligned}
\mathbf{E} \left[\exp\left\{p\gamma\sup_{s\in[t_0,t_0+h]}\delta_{\theta}\bar{y}_s^{(m,q)}\right\}\right]&\leq4\mathbf{E}\left[\exp\left\{16p\gamma\zeta^{(m, q)} \right\}\right]^{\frac{1}{2}}\mathbf{E}\left[\exp\left\{16p\gamma\left(|\eta|+\chi^{(m, q)}\right)\right\}\right]^{\frac{1}{4}}\\
&\quad \times \mathbf{E}\left[\exp\left\{16p\gamma\lambda h\sup _{s \in[t_0, t_0+h]}\left|\delta_\theta Y_s^{(m-1, q)}\right|\right\}\right],   
\end{aligned}
$$
where the pair of processes $\left(\delta_\theta y^{(m, q)}, \delta_\theta z^{(m, q)}\right)$ are defined similarly as Eq. ($\ref{thetaY}$) and
$$
\begin{aligned}
	\zeta^{(m, q)} & =|\eta|+2\lambda h  \sup_m \mathbf{E}\left[\sup _{s \in[t_0, t_0+h]}\left|Y_s^{(m)}\right|\right]+\int_t^{t_0+h} \alpha_s d s+\lambda h\left(\sup _{s \in[t_0, t_0+h]}\left|Y_s^{(m-1)}\right|\right.\\
  &\left.\qquad+\sup _{s \in[t_0, t_0+h]}\left|Y_s^{(m+q-1)}\right|\right), \\
	\chi^{(m, q)} & =4 \lambda h\sup_m \mathbf{E}\left[\sup _{s \in[t_0, t_0+h]}\left|Y_s^{(m)}\right|\right] +\int_t^{t_0+h} \alpha_s d s+2 \lambda h\left(\sup _{s \in[t_0, t_0+h]}\left|Y_s^{(m-1)}\right|\right.\\
 &\left.\qquad+\sup _{s \in[t_0, t_0+h]}\left|Y_s^{(m+q-1)}\right|\right).
 \end{aligned}
 $$

 Recalling \eqref{ytm}, it is easy to check that $\left(\delta_\theta y^{(m, q)}, \delta_\theta z^{(m, q)}\right)$
satisfies the following $\text{BSDE}$
$$
\delta_\theta y_t^{(m, q)}=-\eta+\int_t^{t_0+h}\left(\delta_\theta f^{(m, q)}\left(s, \delta_\theta z_s^{(m, q)}\right)+\delta_\theta f_0^{(m, q)}(s)\right) d s-\int_t^{t_0+h} \delta_\theta z_s^{(m, q)} d B_s,
$$
where
$$
\begin{aligned} 
 \delta_\theta f_0^{(m, q)}(t)&= \frac{1}{1-\theta}\left(f\left(t, Y_t^{(m+q-1)}, \mathbf{P}_{Y_t^{(m+q-1)}}, z_t^{(m)}\right)-f\left(t, Y_t^{(m-1)}, \mathbf{P}_{Y_t^{(m-1)},}, z_t^{(m)}\right)\right), \\
\delta_\theta f^{(m, q)}(t, z)&=  \frac{1}{1-\theta}\left(\theta f\left(t, Y_t^{(m+q-1)}, \mathbf{P}_{Y_t^{(m+q-1)}}, z_t^{(m+q)}\right)\right. \\
& \left.\qquad \qquad \qquad \qquad \qquad-f\left(t, Y_t^{(m+q-1)}, \mathbf{P}_{Y_t^{(m+q-1)}},-(1-\theta) z+\theta z_t^{(m+q)}\right)\right) .\\
\end{aligned}
$$   
Similar to the proof of Lemma 9 in \cite{HMW}, we  obtain that for all $m,q \geq 1$ and $p>1$ 
$$
\begin{aligned}
& \mathbf{E} \left[\exp\left\{p\gamma\sup_{s\in[t_0,t_0+h]}\delta_{\theta}\bar{y}_s^{(m,q)}\right\}\right] \\ 
\leq& 4\mathbf{E}\left[\exp\left\{16p\gamma\zeta^{(m, q)} \right\}\right]^{\frac{1}{2}}\mathbf{E}\left[\exp\left\{8p\gamma\left(|\eta|+\chi^{(m, q)}+\lambda h\sup _{s \in[t_0, t_0+h]} \delta_\theta Y_s^{(m-1, q)}\right)\right\}\right]^{\frac{1}{2}}\\
& \times \mathbf{E}\left[\exp\left\{8p\gamma\lambda h\sup _{s \in[t_0, t_0+h]}\delta_\theta Y_s^{(m-1, q)}\right\}\right]^{\frac{1}{2}} \\
\leq &4\mathbf{E}\left[\exp\left\{16p\gamma\zeta^{(m, q)} \right\}\right]^{\frac{1}{2}}\mathbf{E}\left[\exp\left\{16p\gamma\left(|\eta|+\chi^{(m, q)}\right)\right\}\right]^{\frac{1}{4}}\mathbf{E}\left[\exp\left\{16p\gamma\lambda h\sup _{s \in[t_0, t_0+h]}\delta_\theta Y_s^{(m-1, q)}\right\}\right]^{\frac{1}{4}}\\
& \times \mathbf{E}\left[\exp\left\{16p\gamma\lambda h\sup _{s \in[t_0, t_0+h]}\delta_\theta Y_s^{(m-1, q)}\right\}\right]^{\frac{1}{2}} \\
\leq&4\mathbf{E}\left[\exp\left\{16p\gamma\zeta^{(m, q)} \right\}\right]^{\frac{1}{2}}\mathbf{E}\left[\exp\left\{16p\gamma\left(|\eta|+\chi^{(m, q)}\right)\right\}\right]^{\frac{1}{4}}\mathbf{E}\left[\exp\left\{16p\gamma\lambda h\sup _{s \in[t_0, t_0+h]}\delta_\theta Y_s^{(m-1, q)}\right\}\right].
\end{aligned}
$$

Secondly, through some elementary calculations,  we obtain that $$\delta_{\theta}\bar{Y}_{s}^{(m,q)}\leq \delta_{\theta}\bar{y}_{s}^{(m,q)}+6\frac{C}{c}\sup_{s\in[t_0,t_0+h]}\mathbf{E}[\delta_{\theta}\bar{y}_{s}^{(m,q)}]+2B_3, $$
where $B_3:=2\sup_{s\in[t_0,t_0+h]}\left|\bar{K}^3_s\right|+\left(\frac{C}{c}+1\right)\mathbf{E}\left[\left|\eta\right|\right]+12\frac{C}{c}\sup_{m}\mathbf{E}[\sup_{s\in[t_0,t_0+h]}|y_s^{(m)}|]<\infty$. Applying Lemma $\ref{le-2}$, we have 
$$
\begin{aligned}
 \mathbf{E}\left[\exp\left\{p\gamma\sup_{s\in[t_0,t_0+h]}\delta_{\theta}\bar{Y}_{s}^{(m,q)}\right\}\right]&\leq\widetilde{C}_3 \mathbf{E}\left[\exp\left\{2(1+6\frac{C}{c})16p\gamma\lambda h\sup_{s\in[t_0,t_0+h]}\delta_\theta \bar{Y}_s^{(m-1, q)}\right\}\right]\\
 & \leq\widetilde{C}_3\mathbf{E}\left[\exp\left\{p\gamma\sup_{s\in[t_0,t_0+h]}\delta_\theta \bar{Y}_s^{(m-1, q)}\right\}\right]^{(32+192\frac{C}{c})\lambda h},
\end{aligned}
$$
where $$\widetilde{C}_3=4\mathbf{E}\left[\exp\left\{16p\gamma\zeta^{(m, q)} \right\}\right]^{\frac{1}{2}} \exp\left\{2p\gamma B_3\right\}\mathbf{E}\left[16p\gamma(16+96\frac{C}{c})\left(\left|\eta\right|+\chi^{(m,q)}\right)\right]^{\frac{1}{4}}.$$ 
Set $\widetilde{\rho}=\frac{1}{1-(32+192\frac{C}{c})\lambda h}$. Iterating the above procedure $m$ times yields that
$$
\begin{aligned}
 &\mathbf{E}\left[\exp\left\{p\gamma\sup_{s\in[t_0,t_0+h]}\delta_{\theta}\bar{Y}_{s}^{(m,q)}\right\}\right]
 \leq\widetilde{C}_3^{\widetilde{\rho}} \mathbf{E}\left[\exp\left\{p\gamma\sup_{s\in[t_0,t_0+h]}\delta_\theta \bar{Y}_s^{(1, q)}\right\}\right]^{(32\lambda h+192\frac{C}{c}\lambda h)^{m-1}}.
\end{aligned}
$$
By the uniform estimate in  \eqref{uniformbounded}, we know that for any $\theta\in(0,1)$
$$\lim_{m\to\infty}\sup_{q\geq 1}\mathbf{E}\left[\exp\left\{p\gamma\sup_{s\in[t_0,t_0+h]}\delta_\theta \bar{Y}_s^{(1, q)}\right\}\right]^{(32\lambda h+192\frac{C}{c}\lambda h)^{m-1}}=1.$$
Recalling Eq. \eqref{uniformbounded} again, we obtain that 
$$
\begin{aligned}
\Pi(p)=\sup_{\theta\in(0,1)}\lim_{m\to\infty}\sup_{q\geq 1}\mathbf{E}\left[\exp\left\{p\gamma\sup_{s\in[t_0,t_0+h]}\delta_{\theta}\bar{Y}_{s}^{(m,q)}\right\}\right]\leq \sup_{m,q\geq 1} \widetilde{C}_3^{\widetilde{\rho}} < \infty.
\end{aligned}$$ 
\end{proof}

\fi






\begin{thebibliography}{99} 

\bibitem{BEK} Bank, P. and El Karoui, N. (2004) A stochastic representation theorem with applications to optimization and obstacle problems. The Annals of Probability, 32(1B), 1030-1067.


%\bibitem{BDHPS}  Briand, P., Delyon, B., Hu, Y., Pardoux, E. and Stoica, L. (2003). $L^p$ solutions of backward stochastic differential equations. Stochastic Processes and their Applications, 108(1), 109-129.

\bibitem{BDL} Buckdahn, R., Djehiche, B., Li, J. and Peng, S. (2009). Mean-field backward stochastic differential equations: a limit approach. The Annals of Probability, 37, 1524–1565.

\bibitem{BE} Briand, P., Elie, R. (2013). A simple constructive approach to quadratic BSDEs with or without delay. Stochastic processes and their applications, 123(8), 2921-2939.

\bibitem{BEH}  Briand, P., Elie, R. and Hu, Y. (2018). $\text{BSDEs}$ with mean reflection. The Annals of Applied Probability, 28(1), 482-510.


\bibitem{BH}  Briand, P. and Hibon, H. (2021). Particle systems for mean reflected \text{BSDEs}. Stochastic Processes and their Applications, 131, 253-275.

%\bibitem{BH08} Briand, P. and Hu, Y. (2008). Quadratic \text{BSDEs} with convex generators and unbounded terminal conditions. Probability Theory and Related Fields, 141: 543-567.

\bibitem{BLP} Buckdahn, R., Li, J. and Peng, S. (2009). Mean-field backward stochastic differential equations and related partial differential equations. Stochastic processes and their Applications, 119(10), 3133-3154.

%\bibitem{BY} Bayraktar, E.,Yao, S. (2015). Doubly reflected BSDEs with integrable parameters and related Dynkin games. Stochastic Processes and their Applications, 125(12), 4489-4542.

%\bibitem{BKR} Burdzy, K., Kang, W. and Ramanan, K. (2009) The Skorokhod problem in a time-dependent interval. Stochastic Processes and their Applications, 119: 428-452.

%\bibitem{CCX} Chen, Yajie, Chuanzhi Xing, and Xiao Zhang. "L p solution of general mean-field BSDEs with continuous coefficients." Acta mathematica scientia 40.4 (2020): 1116-1140.

\bibitem{CHM} Chen, Y., Hamadène, S. and  Mu, T. (2023). Mean-field doubly reflected backward stochastic differential equations. Numerical Algebra, Control and Optimization, 13(3-4), 431-460

%\bibitem{CXZ} Chen, Y., Xing, C. and Zhang, X. (2020). $L^p$ solution of general mean-field BSDEs with continuous coefficients. Acta mathematica scientia, 40(4), 1116-1140.

%\bibitem{CM} Crépey, S. and Matoussi, A. (2008). Reflected and doubly reflected \text{BSDEs} with jumps: a priori estimates and comparison. Ann. Appl. Probab. 18(5): 2041-2069

%\bibitem{CK}  Cvitani\'{c}, J. and Karatzas, I. (1996). Backward stochastic differential equations with reflectionand Dynkin games. Ann. Probab., 24(4): 2024-2056.

%\bibitem{DI} Delong, Ł., Imkeller, P. (2010). Backward stochastic differential equations with time delayed generators—results and counterexamples.

\bibitem{DD} Djehiche, B., Dumitrescu, R. (2025). Zero-sum mean-field Dynkin games: characterization and convergence. Mathematics of Operations Research, available online.

\bibitem{DDZ} Djehiche, B., Dumitrescu, R. and Zeng, J. (2025). A propagation of chaos result for weakly interacting nonlinear Snell envelopes. Stochastic Processes and their Applications, 188, 104669.

\bibitem{DEH} Djehiche, B., Elie, R.and Hamadène, S. (2023). Mean-field reflected backward stochastic differential equations. The Annals of Applied Probability, 33(4), 2493-2518.

%\bibitem{DRZ} Dos Reis, G., Réveillac, A., Zhang, J. (2011). FBSDEs with time delayed generators: Lp-solutions, differentiability, representation formulas and path regularity. Stochastic processes and their applications, 121(9), 2114-2150.
%\bibitem{DQS} Dumitrescu, R., Quenez, M.C. and Sulem, A. (2016) Generalized Dynkin games and doubly reflected BSDEs with jumps. Electronic Journal of Probability, 21: 1-32.

\bibitem{EKPPQ} El Karoui, N., Kapoudjian, C., Pardoux, E., Peng, S. and Quenez, M.C. (1997). Reflected solutions
of backward SDE’s, and related obstacle problems for PDE’s. The Annals of Probability, 25(2), 702-737.

\bibitem{EPQ} El Karoui, N., Peng, S. and Quenez, M. C. (1997). Backward stochastic differential equations in finance. Mathematical finance, 7(1), 1-71.

%\bibitem{EPQ1} El Karoui, N., Pardoux, E. and Quenez, M.C. (1997) {Reflected backward SDE's and  American options}. Numerical Methods in Finance (Cambridge Univ. Press), 1997: 215-231.

%\bibitem{EPQ2} El Karoui, N., Peng, S. and Quenez M.C. (2001) A dynamic maximum principle for the optimization of recursive utilities under constraints.  Ann. Appl. Probab., 664-693.

%\bibitem{EPQ} El Karoui N, Pardoux E, Quenez M.C. Backward stochastic differential equations in finance. Math. Finance, 1997, 7: 1-71

%\bibitem{FS21} Falkowski, A. and Slomi\'{n}ski, L. (2021) Mean reflected stochastic differential equations with  two constraints. Stochastic Processes and their Applications, 141: 172-196.

\bibitem{FS} Falkowski, A. and Słomiński, L. (2022). Backward stochastic differential equations with mean reflection and two constraints. Bulletin des Sciences Mathématiques, 176, 103117.

%\bibitem{GIOOQ} Grigorova, M., Imkeller, P., Offen, E., Ouknine, Y.and Quenez, M. C. (2017). Reflected \text{BSDEs} when the obstacle is not right-continuous and optimal stopping. Ann. Appl. Probab. 27(5): 3153-3188

%\bibitem{GIOQ} Grigorova, M., Imkeller, P., Ouknine, Y. and Quenez, M.C. (2018) Doubly reflected BSDEs and $\mathcal{E}^f$-Dynkin games: beyond the right-continuous case. Electronic Journal of Probability, 23: 1-38.

%\bibitem{GLX} Gu, Z., Lin, Y. and Xu, K. (2024). Mean reflected \text{BSDE} driven by a marked point process and application in insurance risk management. ESAIM: Control, Optimisation and Calculus of Variations, 30, 51.

%\bibitem{HL} Hamadène, S. and Lepeltier, J. P. (2000). Reflected \text{BSDEs} and mixed game problem. Stochastic processes and their applications, 85(2), 177-188.

%\bibitem{HLM} Hamadene, S.,  Lepeltier, J.-P. and Matoussi, A. (1997) Double barrier backward SDEs with continuous coefficient. in: El Karoui, L. Mazliak (Eds.), in : Pitman Res. Notes  Math. Ser., 364: 161-171.

\bibitem{HHLLW} Hibon, H., Hu, Y., Lin, Y., Luo, P.and Wang, F. (2017). Quadratic \text{BSDEs} with mean reflection. Mathematical Control and Related Fields, 8, 721-738.

\bibitem{HHTW} Hao, T., Hu, Y., Tang, S., Wen, J. (2025). Mean-field backward stochastic differential equations and nonlocal PDEs with quadratic growth. The Annals of Applied Probability, 35(3), 2128-2174.

\bibitem{HMW'} Hu, Y., Moreau, R. and Wang, F. (2022). Quadratic mean-field reflected \text{BSDEs}. Probability, Uncertainty and Quantitative Risk, 7(3): 169–194.

\bibitem{HMW} Hu, Y., Moreau, R.and Wang F. (2024). General mean reflected backward stochastic differential equations. Journal of Theoretical Probability, 37(1): 877-904.

%\bibitem{K1} Klimsiak, T. (2012). Reflected \text{BSDEs} with monotone generator. Electronic Journal of Probability, 17, 1-25.

%\bibitem{K2} Klimsiak, T. (2013). \text{BSDEs} with monotone generator and two irregular reflecting barriers. Bulletin des Sciences Mathématiques, 137(3), 268-321.

%\bibitem{KLQT} Kobylański, M., Lepeltier, J. P., Quenez, M. C. and Torres, S. (2002). Reflected \text{BSDE} with superlinear quadratic coefficient. Probability and Mathematical Statistics, 22.

\bibitem{Li} Li, H. (2023).  The Skorokhod problem with two nonlinear constraints. Probability and Mathematical Statistics, 43(2): 207-239.
	
\bibitem{L} Li, H. (2024). Backward stochastic differential equations with double mean reflections. Stochastic Processes and their Applications, 173, 104371.

\bibitem{LS} Li, H., Shi, J. (2025). Mean Field Backward Stochastic Differential Equations with Double Mean Reflections. arXiv preprint arXiv:2501.10939.

\bibitem{Li'} Li, J. (2014). Reflected mean-field backward stochastic differential equations. Approximation and associated nonlinear \text{PDEs}. Journal of Mathematical Analysis and Applications, 413(1), 47-68.

%\bibitem{LLZ} Li, J., Liang, H. and Zhang, X. (2018). General mean-field \text{BSDEs} with continuous coefficients. Journal of Mathematical Analysis and Applications, 466(1), 264-280.

 %\bibitem{LW} Liu, G. and Wang, F. (2019). \text{BSDEs} with mean reflection driven by G-Brownian motion. Journal of Mathematical Analysis and Applications, 470(1), 599-618.

 \bibitem{LX} Li, J. and Xing, C. (2022). General mean-field BDSDEs with continuous coefficients. Journal of Mathematical Analysis and Applications, 506(2), 125699.

\bibitem{LW} Liu, G. and Wang, F. (2019). BSDEs with mean reflection driven by G-Brownian motion. Journal of Mathematical Analysis and Applications, 470(1), 599-618.
 
 \bibitem{LP}Luo, P. (2024). Mean-field backward stochastic differential equations with mean reflection and nonlinear resistance. Stochastics, 96(8), 2037-2061.

 \bibitem{MW} Ma, J. and Wang, Y. (2009) On variant reflected backward SDEs, with applications. J. Appl. Math. Stoch. Anal., Article ID 854768. 

\bibitem{NQW} Niu, Y., Qu, B. and Wang, F. (2025). Lp-solutions of multi-dimensional BSDEs with mean reflection. Stochastic Processes and their Applications, 187, 104663.

\bibitem{PP} Pardoux, E. and Peng, S. (1990). Adapted solution of a backward stochastic differential equation. Systems \& control letters, 14(1), 55-61.

%\bibitem{LM} Lepeltier, J-P., and Mingyu Xu. "Penalization method for reflected backward stochastic differential equations with one rcll barrier." Statistics  probability letters 75.1 (2005): 58-66.

%\bibitem{MPZ} Matoussi A, Possamai D, Zhou C. Second order reflected backward stochastic differential equations[J]. 2013.

%\bibitem{P} Peng S. Backward stochastic differential equations and applications to optimal control[J]. Applied Mathematics and Optimization, 1993, 27(2): 125-144.

%\bibitem{MM} Mwigilwa W F, Mhlanga F J. BSDE with Jumps When Mean Reflection Is Nonlinear[J]. International Journal of Mathematics and Mathematical Sciences, 2024, 2024(1): 9963889.

%\bibitem{PP} Pardoux E, Peng S. Adapted solution of a backward stochastic differential equation[J]. Systems control letters, 1990, 14(1): 55-61.

%\bibitem{PP91} Pardoux E, Peng S. Backward stochastic differential equations and quasilinear parabolic partial differential equations[C]//Stochastic Partial Differential Equations and Their Applications: Proceedings of IFIP WG 7/1 International Conference University of North Carolina at Charlotte, NC June 6–8, 1991. Berlin, Heidelberg: Springer Berlin Heidelberg, 2005: 200-217.

%\bibitem{PX} Peng, S. and Xu, M. (2005). The smallest $ g $-supermartingale and reflected \text{BSDE} with single and double ${L}^{2} $ obstacles. In Annales de l'IHP Probabilités et statistiques (Vol. 41, No. 3, pp. 605-630).

\bibitem{QF} Qu, B. and Wang, F. (2023). Multi-dimensional \text{BSDEs} with mean reflection. Electronic Journal of Probability, 28, 1-26.

\bibitem{QX} Qian, Z. and Xu, M. (2018). Reflected backward stochastic differential equations with resistance. The Annals of Applied Probability, 28(2), 888-911.

%\bibitem{XZ} Xing, Z. (2022). Mean-field backward stochastic differential equations with reflection and related nonlocal PDEs in a convex domain. Systems \& Control Letters, 159, 105081.

%\bibitem{ZJ} Zhang, J. (2017). Backward stochastic differential equations. Springer New York.

\end{thebibliography}
\end{document}
